\chapter{Cardiology}

\medterm{ACE Inhibitor}-- A drug that inhibits angiotensin-converting enzyme, reducing formation of angiotensin II and degradation of bradykinin. ACE inhibitors reduce blood pressure and afterload and improve outcomes in selected patients with heart failure, chronic kidney disease, diabetes, or post-myocardial infarction; adverse effects include cough, hyperkalemia, hypotension, and angioedema.

\medterm{Acute Coronary Syndrome}-- An operational umbrella term encompassing a clinical continuum of acute myocardial ischemic states precipitated by the sudden, critical reduction of coronary arterial blood flow, almost universally initiated by the rupture or erosion of an unstable, thin-cap fibroatheroma (TCFA) followed by platelet activation and intravascular thrombus formation. ACS is triaged into three distinct clinical and diagnostic entities based on 12-lead ECG findings and serial cardiac biomarker assays: (1) ST-Elevation Myocardial Infarction (STEMI): characterized by complete, persistent, transmural thrombotic occlusion of an epicardial coronary artery, producing persistent ST-segment elevations (or new LBBB/true posterior MI criteria) and elevated cardiac troponins, requiring emergent reperfusion therapy within strictly mandated time windows (primary PCI with door-to-balloon time $<90$ minutes, or fibrinolysis within 30 minutes if PCI is unavailable within 120 minutes); (2) Non-ST-Elevation Myocardial Infarction (NSTEMI): characterized by subtotal or intermittent occlusive thrombosis causing subendocardial ischemic necrosis, manifesting with ischemic symptoms, ST-segment depressions, T-wave inversions, or transient ST elevations, accompanied by a dynamic rise and/or fall of cardiac troponin (cTnI/cTnT) above the 99th percentile upper reference limit; and (3) Unstable Angina (UA): ischemic chest pain occurring at rest, of new onset, or crescendo in nature, without objective biochemical evidence of myocardial necrosis (negative troponins). Initial medical management across all ACS includes chewable aspirin (162--325 mg), a potent $P2Y_{12}$ receptor inhibitor (ticagrelor, prasugrel, or clopidogrel), anticoagulation (unfractionated heparin, enoxaparin, or bivalirudin), high-intensity statin therapy, anti-ischemic agents (nitrates, beta-blockers), and risk-stratified invasive coronary angiography.

\medterm{Acute Limb Ischemia}-- Sudden reduction in arterial blood flow to a limb that threatens tissue viability, usually caused by embolism, thrombosis, or graft occlusion. The six classic findings are pain, pallor, pulselessness, paresthesia, paralysis, and poikilothermia; urgent vascular assessment and revascularization are required.

\medterm{Acute mesenteric ischemia}-- A life-threatening vascular emergency characterized by the sudden, critical reduction of blood flow through the mesenteric arterial circulation (most commonly the superior mesenteric artery [SMA]), leading to intestinal cellular hypoperfusion, mucosal sloughing, transmural bowel necrosis, peritonitis, sepsis, and death. Etiologies are categorized into: (1) Superior Mesenteric Artery Embolism (~50\% of cases): cardiac emboli (atrial fibrillation, post-MI mural thrombus, endocarditis) lodge at anatomical narrowings, characteristically 3 to 8 cm distal to the SMA origin beyond the middle colic artery, sparing the proximal jejunum; (2) SMA Thrombosis (~25\%): acute superimposed thrombosis on preexisting severe atherosclerotic plaque at the SMA ostium; (3) Non-Occlusive Mesenteric Ischemia (NOMI, ~20\%): severe splanchnic vasoconstriction triggered by low-flow states (cardiogenic shock, sepsis, high-dose vasopressor infusions); and (4) Mesenteric Venous Thrombosis (MVT, ~5\%): hypercoagulable states. The clinical hallmark is severe, agonizing periumbilical abdominal pain out of proportion to physical examination findings (the abdomen remains soft and non-tender initially). Delayed diagnosis results in transmural infarction, marked by peritoneal signs (guarding, rigidity), bloody diarrhea, and profound lactic acidosis. Computed tomography angiography (CTA) of the abdomen and pelvis is the definitive diagnostic modality of choice. Immediate management requires systemic heparinization, broad-spectrum antibiotics, hemodynamic resuscitation, and emergent revascularization (endovascular thrombectomy/thrombolysis or surgical embolectomy/bypass) alongside resection of non-viable bowel with planned second-look laparotomy.

\medterm{Afterload}-- The resistance against which the ventricles must eject blood during systole. Left ventricular afterload is primarily determined by systemic vascular resistance (SVR); right ventricular afterload is determined by pulmonary vascular resistance (PVR). Afterload is estimated by blood pressure or calculated as wall stress. Increased afterload (as in hypertension or aortic stenosis) reduces stroke volume and increases myocardial oxygen demand.

\medterm{Alcoholic Cardiomyopathy}-- Dilated cardiomyopathy caused by chronic excessive alcohol consumption, one of the most common reversible causes of dilated cardiomyopathy. Direct toxic effect of ethanol and its metabolite acetaldehyde on cardiomyocytes. Presentation is identical to other dilated cardiomyopathies. Abstinence is the cornerstone of treatment and can lead to significant recovery of LV function. Other treatments follow standard heart failure guidelines.

\medterm{Amiodarone}-- A highly lipophilic, iodine-rich benzofuran derivative that serves as a broad-spectrum antiarrhythmic agent. Formally classified as a Vaughan Williams Class III agent due to its primary action of blocking rapid delayed rectifier potassium currents ($I_{Kr}$), which prolongs the myocardial action potential duration and effective refractory period across atrial, ventricular, and accessory pathway tissues. However, amiodarone uniquely exhibits electrophysiological properties across all four Vaughan Williams classes: Class I (inactivated sodium channel blockade with rapid kinetics), Class II (non-competitive $\alpha$- and $\beta$-adrenergic receptor antagonism), and Class IV (L-type calcium channel blockade, slowing SA and AV nodal conduction). Indicated for the acute termination and suppression of life-threatening ventricular arrhythmias (VT, VF) and the rhythm or rate control of supraventricular arrhythmias (atrial fibrillation, flutter), especially in patients with structural heart disease or reduced ejection fraction where other antiarrhythmics are proarrhythmic or contraindicated. Because each 200-mg tablet contains ~75 mg of iodine and the drug has a massive volume of distribution with a prolonged elimination half-life (40--60 days), chronic administration requires vigilant multi-organ toxicity monitoring: pulmonary fibrosis (most fatal, presenting with cough and dyspnea with reduced DLCO), thyroid dysfunction (amiodarone-induced hypothyroidism via Wolff-Chaikoff effect, or type 1/2 thyrotoxicosis via Jod-Basedow effect or destructive thyroiditis), hepatotoxicity (elevated transaminases), corneal microdeposits (vortex keratopathy, nearly universal but benign), optic neuropathy, peripheral sensory-motor neuropathy, blue-gray slate skin discoloration, and photosensitivity. It is a potent inhibitor of CYP3A4, CYP2C9, and P-glycoprotein, doubling serum levels of digoxin and warfarin (requiring a 50\% dose reduction).

\medterm{Angina Pectoris}-- Chest discomfort caused by myocardial ischemia without infarction. Stable angina occurs predictably with exertion and resolves with rest. Unstable angina occurs at rest or with minimal exertion and is an acute coronary syndrome. Prinzmetal (variant) angina occurs at rest due to coronary vasospasm. Diagnosis involves ECG (may show ST depression during episodes) and stress testing. Treatment includes nitrates, beta-blockers, and revascularization when indicated.

\medterm{Anticoagulant Monitoring}-- Unfractionated heparin is monitored with aPTT (target 1.5-2.5 x control). Low molecular weight heparins generally do not require monitoring. Warfarin is monitored with INR (target 2-3 for most indications, 2.5-3.5 for mechanical mitral valve). DOACs (rivaroxaban, apixaban, dabigatran) have predictable pharmacokinetics and do not require routine monitoring, though anti-Xa levels can be measured for rivaroxaban/apixaban and dilute thrombin time for dabigatran.

\medterm{Aortic Aneurysm}-- Permanent dilation of the aorta exceeding 50\% of normal diameter. Abdominal aortic aneurysm (AAA) is more common than thoracic aortic aneurysm and is usually caused by atherosclerosis. AAA >5.5 cm or rapid expansion (>0.5 cm/year) warrants surgical repair (open or endovascular). Thoracic aortic aneurysms may be caused by Marfan syndrome, bicuspid aortic valve, or hypertension. Screening ultrasound is recommended for men aged 65-75 who have ever smoked.

\medterm{Aortic Coarctation}-- Congenital narrowing of the aorta, usually near the insertion of the ductus arteriosus. It causes higher blood pressure and stronger pulses in the arms than in the legs, with a radio-femoral delay; severe neonatal disease may present with shock when the ductus closes. Diagnosis uses echocardiography or cross-sectional vascular imaging, and treatment may involve surgery or catheter intervention.

\medterm{Aortic Dissection}-- A tear in the intima of the aorta allowing blood to dissect through the media, creating a false lumen. It is classified as Stanford Type A (ascending aorta involvement, requiring emergency surgery) or Type B (descending aorta only, managed medically unless complicated). Risk factors include hypertension, Marfan syndrome, Ehlers-Danlos syndrome, bicuspid aortic valve, and aortic instrumentation. Symptoms include sudden, severe, tearing chest pain radiating to the back.

\medterm{Aortic embolism}-- The acute mechanical occlusion of a distal systemic arterial lumen by an embolus originating within the aorta, arising from disrupted mural thrombus, ulcerated atherosclerotic plaque debris, or catastrophic aortic thromboembolism. A classic and dramatic clinical manifestation is the Saddle Embolus, where a large, organized thromboembolus lodges directly astride the aortic bifurcation (carina) into the common iliac arteries. The source is most commonly cardiac (atrial fibrillation, left ventricular mural thrombus following anterior MI, dilated cardiomyopathy, or left atrial myxoma) or directly aortic (severe shaggy or mobile thoracic/abdominal aortic atheroma, or aneurysmal mural thrombus). Complete occlusion of the terminal aorta produces bilateral acute lower-extremity ischemia, presenting with the classic '6 P's': Pain, Pallor, Pulselessness (absent femoral, popliteal, and pedal pulses bilaterally), Paresthesias, Poikilothermia (cold extremities), and Paralysis. Severe ischemia can cause paraplegia mimicking acute spinal cord infarction, alongside sensory loss in a stocking distribution. Diagnostic confirmation is established by emergent computed tomography angiography (CTA) of the abdominal aorta and lower extremities. Aortic saddle embolism represents an extreme, time-critical surgical emergency: immediate intravenous unfractionated heparin must be administered, followed by emergent bilateral femoral arterial cutdowns and retrograde transfemoral balloon catheter embolectomy (Fogarty catheter thrombectomy) or open surgical intervention to restore lower limb perfusion within 4 to 6 hours and prevent irreversible neuromuscular necrosis, myoglobinuria, and fatal reperfusion metabolic syndrome.

\medterm{Aortic Regurgitation}-- Incompetence of the aortic valve causing abnormal diastolic backflow of blood from the aorta into the left ventricle. Etiologies include aortic root dilation (Marfan syndrome, bicuspid aortic valve, thoracic aortic aneurysm, chronic hypertension), degenerative calcific valve disease, infective endocarditis, and acute aortic dissection (type A). Acute AR causes sudden pulmonary edema and cardiogenic shock due to an unprepared, non-compliant left ventricle. Chronic AR induces chronic volume and pressure overload, producing massive eccentric left ventricular hypertrophy and chamber dilation ('cor bovinum'). Auscultation reveals a high-pitched, blowing early diastolic decrescendo murmur heard best along the left sternal border at end-expiration with the patient leaning forward. In severe chronic AR, an Austin Flint murmur (a low-pitched, mid-to-late diastolic rumbling apical murmur caused by the regurgitant jet impinging on the anterior mitral valve leaflet) may be present. Marked widening of pulse pressure generates classic peripheral eponymic physical signs: Corrigan's pulse ('water-hammer' bounding carotid pulse with rapid collapse), Duroziez's sign (systolic and diastolic murmurs over the femoral artery with stethoscope compression), Quincke's sign (capillary pulsations in nail beds), and de Musset's sign (rhythmic systolic head bobbing). Echocardiography establishes severity (regurgitant volume $\ge 60$ mL/beat, effective regurgitant orifice area $\ge 0.30$ cm$^2$). Surgical aortic valve replacement is indicated for symptomatic severe AR or asymptomatic patients with LVEF $\le 50\%$ or LV end-systolic diameter (LVESD) $>50$ mm.

\medterm{Aortic Stenosis}-- Narrowing of the aortic valve opening, obstructing left ventricular outflow during systole. Etiologies comprise calcific degeneration of a normal tricuspid valve (typically presenting in patients aged 70--80 years), calcification of a congenital bicuspid aortic valve (presenting earlier in life, aged 40--60 years), and post-inflammatory rheumatic heart disease. Chronic outflow obstruction increases systolic afterload, driving compensatory concentric left ventricular hypertrophy according to Laplace's Law to normalize wall stress, which progressively impairs diastolic relaxation and increases myocardial oxygen demand. The cardinal symptoms emerge with exertion: dyspnea (heart failure, 50\% 2-year mortality if untreated), angina (50\% 5-year mortality), and syncope (50\% 3-year mortality; SAD triad). Physical examination reveals a harsh, late-peaking crescendo-decrescendo systolic ejection murmur heard best at the right second intercostal space radiating to the carotids, a diminished and delayed carotid pulse (pulsus parvus et tardus), and a soft or absent aortic component of the second heart sound ($A_2$). Severity is graded echocardiographically: mild (valve area $>1.5$ cm$^2$, mean gradient $<20$ mmHg), moderate (area 1.0--1.5 cm$^2$, mean gradient 20--40 mmHg), and severe (aortic valve area $<1.0$ cm$^2$, mean pressure gradient $\ge 40$ mmHg, and peak jet velocity $\ge 4.0$ m/s). Definitive treatment for severe symptomatic AS is valve replacement via Transcatheter Aortic Valve Replacement (TAVR) or Surgical Aortic Valve Replacement (SAVR).

\medterm{ARB-Neprilysin Inhibitor (ARNI)}-- Sacubitril/valsartan combines an ARB with a neprilysin inhibitor, increasing natriuretic peptide levels while blocking the renin-angiotensin system. The PARADIGM-HF trial demonstrated superior reduction in cardiovascular death and heart failure hospitalization compared to enalapril. ARNI is now a cornerstone of HFrEF therapy, recommended as replacement for ACE inhibitors in patients who remain symptomatic on optimal medical therapy.

\medterm{Arrhythmia}-- Any abnormality of cardiac rate, rhythm, or site of impulse origin, arising through enhanced automaticity, reentry, or triggered activity. Classified by rate as tachyarrhythmia or bradyarrhythmia and by site as supraventricular or ventricular. Diagnosis rests on the 12-lead ECG and ambulatory monitoring when intermittent. Management follows ACLS algorithms and ranges from drugs and ablation to cardioversion, defibrillation, and pacing.

\medterm{Arrhythmogenic right ventricular cardiomyopathy}-- A genetic cardiomyopathy characterized by fibrofatty replacement of the right ventricular myocardium, predisposing to ventricular arrhythmias and sudden cardiac death, particularly during exercise. ECG shows T-wave inversions in V1-V3 and epsilon waves; diagnosis uses echocardiography, cardiac MRI, and genetic testing.

\medterm{Arterial insufficiency}-- A clinical and hemodynamic state characterized by inadequate arterial blood delivery to tissues or organ systems, failing to meet baseline metabolic and oxygen demands. In the peripheral extremities, it is the defining feature of Peripheral Artery Disease (PAD), predominantly caused by systemic atherosclerosis of the aorta, iliac, femoropopliteal, and infrapopliteal tibial arteries. Progressive luminal narrowing causes exertional muscle ischemia, presenting classically as Intermittent Claudication (reproducible, cramp-like muscle pain in the calf, thigh, or buttock brought on by walking and relieved within 10 minutes of rest). As stenosis progresses, resting blood flow becomes deficient, advancing to Chronic Limb-Threatening Ischemia (CLTI / critical limb ischemia), characterized by ischemic rest pain (burning pain in the distal foot and toes, worsened by leg elevation in bed and relieved by hanging the leg over the bedside), non-healing arterial ischemic ulcers (punched-out, painful ulcers with pale necrotic bases on the toes, heels, and bony prominences), and ischemic gangrene. Physical signs include diminished or absent distal pulses, dependent rubor with pallor on elevation (Buerger sign), atrophic hairless shiny skin, dystrophic toenails, cool extremities, and delayed capillary refill ($>3$ seconds). Initial non-invasive screening utilizes the Ankle-Brachial Index (ABI): normal 1.0--1.3; mild-to-moderate PAD 0.41--0.90; severe ischemia $\le 0.40$; and non-compressible calcified vessels $>1.40$. Management relies on smoking cessation, antiplatelet therapy (aspirin or clopidogrel), high-intensity statins, structured exercise therapy, cilostazol, and revascularization (endovascular balloon angioplasty/stenting or surgical bypass grafting).

\medterm{Arteriosclerosis}-- A broad, generic clinicopathological term denoting the thickening, hardening, loss of elasticity, and fibrosis of arterial walls ('hardening of the arteries'). Arteriosclerosis encompasses three distinct pathological and anatomical entities: (1) Atherosclerosis: the most common and clinically dominant form, characterized by lipid-rich fibrofatty atheromatous plaques within the tunica intima of large elastic arteries (aorta, carotid) and medium-sized muscular arteries (coronary, popliteal, cerebral), driving ischemic heart disease, stroke, and gangrene; (2) Arteriolosclerosis: pathological thickening and luminal narrowing of small arteries and arterioles ($<100$ $\mu$m diameter), presenting either as hyaline arteriolosclerosis (homozygous plasma protein leakage and smooth muscle matrix deposition driven by chronic, moderate systemic hypertension and diabetes mellitus) or hyperplastic arteriolosclerosis (concentric, laminated, 'onion-skin' smooth muscle and basement membrane proliferation with necrotizing arteriolitis driven by severe, malignant hypertension); and (3) Mönckeberg Medial Calcific Sclerosis: dystrophic calcification localized entirely within the internal elastic lamina and tunica media of medium-sized muscular arteries (predominantly in elderly individuals, patients with diabetes mellitus, and end-stage renal disease); because it spares the tunica intima and does not occlude the vessel lumen, it does not directly cause luminal narrowing or tissue ischemia, but causes stiff, pipe-like non-compressible arteries resulting in an artifactually elevated Ankle-Brachial Index ($\text{ABI} > 1.40$).

\medterm{Aspirin}-- Acetylsalicylic acid, a foundational antiplatelet agent and nonsteroidal anti-inflammatory drug (NSAID) that exerts an irreversible inhibitory effect on platelet function. Aspirin covalently acetylates a specific serine residue (Ser529) within the active-site channel of cyclooxygenase-1 (COX-1), permanently abolishing its catalytic activity. Because platelets are anucleate fragments that lack the transcriptional machinery to synthesize new enzymes, COX-1 inhibition persists for the entire biological lifespan of the platelet (7 to 10 days). Blockade of COX-1 halts the conversion of arachidonic acid to prostaglandin $H_2$, eliminating the synthesis of thromboxane $A_2$ ($TXA_2$), a potent platelet-activating agonist and vasoconstrictor. In acute coronary syndromes (STEMI, NSTEMI, unstable angina), immediate chewable non-enteric aspirin (162--325 mg) achieves rapid platelet inhibition within 20 minutes, reducing 30-day vascular mortality by ~23\%. For secondary cardiovascular prevention (following MI, ischemic stroke, or coronary stenting), low-dose daily aspirin (75--100 mg) is maintained indefinitely, balancing antithrombotic efficacy with minimized gastric toxicity. Major adverse effects include gastrointestinal ulceration and bleeding (due to gastric mucosal $PGE_2$ suppression), aspirin-exacerbated respiratory disease (AERD / Samter triad: asthma, nasal polyps, and bronchospasm driven by leukotriene shunting), and Reye syndrome (fulminant hepatic failure and encephalopathy in children treated with aspirin during viral illnesses like influenza or varicella).

\medterm{Asystole}-- Complete cessation of cardiac electrical activity, presenting as a flat line on the ECG with no QRS complexes, P waves, or other organized electrical activity. Referred to as "cardiac standstill." Causes: end-stage of any terminal rhythm, profound hypoxia, severe electrolyte disturbance (hyperkalemia), cardiac tamponade, massive pulmonary embolism, drug overdose, or terminal cardiac disease. Confirmed in two leads with correct lead connections. Management: high-quality CPR, IV epinephrine (1 mg every 3-5 minutes), and search for reversible causes (H's and T's); unsynchronized defibrillation is not indicated. Asystole has very poor prognosis; pulseless electrical activity or ventricular fibrillation is excluded before calling it.

\medterm{Atherosclerosis}-- A chronic, progressive, systemic inflammatory fibroproliferative disease of the elastic and large-to-medium muscular arteries (aorta, coronary, carotid, cerebral, and peripheral arteries) characterized by the focal accumulation of lipids, inflammatory leukocytes, vascular smooth muscle cells, and extracellular matrix components within the arterial intima, forming fibrofatty atherosclerotic plaques. Pathogenesis follows the 'response-to-injury' model: (1) Endothelial Dysfunction: Chronic hemodynamic shear stress, cigarette smoke toxins, hyperglycemia, and elevated circulating ApoB-containing lipoproteins (LDL, VLDL) injure the vascular endothelium, increasing permeability and inducing adhesion molecules (VCAM-1, ICAM-1, E-selectin); (2) Fatty Streak Formation: Subendothelial LDL particles become trapped in the extracellular matrix and undergo enzymatic and oxidative modification (oxLDL). Circulating monocytes adhere, extravasate into the intima, and differentiate into macrophages that express scavenger receptors (SR-A, CD36), engulfing oxLDL unchecked to become lipid-laden foam cells; (3) Plaque Progression: T-lymphocytes (Th1) and foam cells secrete pro-inflammatory cytokines (IFN-$\gamma$, TNF-$\alpha$, IL-1), triggering vascular smooth muscle cell (VSMC) migration from the media into the intima. VSMCs proliferate and secrete extracellular collagen and proteoglycans, creating a protective fibrous cap over a necrotic lipid core formed by apoptotic foam cells; and (4) Plaque Complication: Activated macrophages secrete matrix metalloproteinases (MMPs) that degrade collagen in the fibrous cap. A thin-cap fibroatheroma (TCFA, cap $<65 \ \mu\text{m}$) becomes vulnerable to rupture, erosion, or intraplaque hemorrhage, exposing thrombogenic subendothelial core components (tissue factor) to flowing blood, triggering acute platelet adhesion, thrombus propagation, and acute arterial occlusion (MI, ischemic stroke).

\medterm{Atrial Fibrillation}-- The most common sustained cardiac arrhythmia, characterized by rapid, irregular atrial activity (350-600 bpm) with irregular ventricular response. ECG shows absent P waves, irregularly irregular RR intervals, and narrow QRS complexes (unless aberrant conduction). AFib increases stroke risk 5-fold due to atrial stasis and thrombus formation, particularly in the left atrial appendage. Treatment targets rate control, rhythm control, and anticoagulation.

\medterm{Atrial Fibrillation Ablation}-- Catheter-based procedure to eliminate atrial fibrillation by creating pulmonary vein isolation (PVI), electrically disconnecting the pulmonary veins from the left atrium. Indicated for symptomatic AFib refractory to or intolerant of antiarrhythmic drugs. Success rates: 70-80\% for paroxysmal AFib, 50-60\% for persistent AFib. Complications include pulmonary vein stenosis, pericardial tamponade, atrio-esophageal fistula (rare but life-threatening), and stroke. Anticoagulation is continued post-procedure.

\medterm{Atrial Flutter}-- A rapid, regular atrial rhythm (typically 250-350 bpm) caused by a macro-reentrant circuit, most commonly in the right atrium. ECG shows sawtooth flutter waves best seen in leads II, III, and aVF, with regular or irregular ventricular response depending on AV conduction ratio (commonly 2:1 or 4:1). Atrial flutter carries similar stroke risk to AFib and requires anticoagulation. Treatment includes cardioversion and catheter ablation.

\medterm{Atrial Septal Defect}-- Congenital opening in the interatrial secundum allowing left-to-right shunting of blood. Small defects may close spontaneously. Larger defects cause right heart volume overload, right atrial and ventricular dilation, pulmonary hypertension, and eventually right heart failure. Paradoxical embolism may occur with right-to-left shunting. Systolic crescendo-decrescendo murmur at left upper sternal border with fixed split S2. Echocardiography confirms diagnosis and assesses shunt magnitude. Transcatheter closure is indicated for moderate-to-large defects.

\medterm{Atrial Tachycardia}-- Supraventricular tachycardia arising from an atrial focus outside the sinus node, often focal with rates 100-250 beats per minute with discrete P waves of different morphology from sinus P waves. Causes: triggered activity (digoxin toxicity with automatic atrial tachycardia), focal atrial tachycardia (structural disease, post-ablation reentry), multifocal atrial tachycardia (MAT) seen in severe COPD; distinct lone P-wave morphologies (at least three). Treatment: address underlying cause; beta-blockers, calcium channel blockers (rate control), class IC or III antiarrhythmics, and catheter ablation. Carotid sinus massage may not terminate automatic atrial tachycardia.

\medterm{Atrioventricular Node}-- Located in the posterior-inferior region of the interatrial septum near the coronary sinus. It receives impulses from the SA node via internodal pathways and delays them by approximately 120-200 milliseconds before transmitting them to the ventricles. This delay allows atrial contraction to complete before ventricular contraction. The AV node is the backup pacemaker with an intrinsic rate of 40-60 beats per minute.

\medterm{Austin Flint Murmur}-- A mid-diastolic, low-pitched rumbling murmur heard at the apex, caused by the regurgitant jet from severe aortic regurgitation striking the anterior mitral leaflet. This functional murmur can mimic mitral stenosis but is distinguished by the absence of an opening snap, presence of wide pulse pressure, and early diastolic decrescendo murmur of aortic regurgitation. It resolves after aortic valve replacement.

\medterm{Balloon Aortic Valvuloplasty}-- Catheter-based procedure using balloon inflation to dilate a stenotic aortic valve. Provides temporary hemodynamic improvement in severe AS but restenosis is common within 6-12 months. Used as a bridge to TAVR or surgical AVR in hemodynamically unstable patients, or as palliation in patients not eligible for definitive valve replacement. Complications include aortic regurgitation, stroke, and vascular access complications.

\medterm{Basilar artery aneurysm}-- A intracranial saccular (berry) or fusiform dilation of the basilar artery, representing the most common aneurysm of the posterior cerebral circulation (accounting for approximately 5--8\% of all intracranial aneurysms), typically arising at the basilar bifurcation ('basilar tip' or basilar apex, where it divides into the posterior cerebral arteries) or at the junction with the superior cerebellar or anterior inferior cerebellar arteries. Aneurysms at the basilar apex are exposed to intense, continuous hemodynamic impingement from converging vertebral arterial flows. Unruptured aneurysms may be asymptomatic or produce compressive cranial neuropathies and brainstem signs: compression of the third cranial nerve (oculomotor nerve palsy: ipsilateral ptosis, pupillary dilation, and 'down-and-out' eye deviation) or brainstem ischemia from parent vessel or perforator branch thrombosis. Rupture of a basilar artery aneurysm precipitates devastating aneurysmal Subarachnoid Hemorrhage (aSAH), presenting with an explosive 'thunderclap' headache ('worst headache of life'), vomiting, meningismus, coma, and high immediate mortality ($>50\%$). Posterior circulation aneurysms carry a significantly higher intrinsic rupture risk per millimeter of diameter than anterior circulation aneurysms. Diagnosis is established by non-contrast head CT followed by CT angiography (CTA) or digital subtraction angiography (DSA). Treatment is predominantly endovascular (microcoil embolization, stent-assisted coiling, or flow-diverting stents) due to the hazardous, complex surgical exposure required for open microsurgical clipping around the ventral brainstem.

\medterm{Beck Triad}-- A classic constellation of three physical examination findings described in 1935 by the American cardiovascular surgeon Claude Beck that indicates acute Cardiac Tamponade: (1) Systemic Arterial Hypotension with narrow pulse pressure (resulting from severe reduction in left ventricular stroke volume and cardiac output); (2) Elevated Jugular Venous Pressure with prominent jugular venous distension (caused by impaired venous return into the compressed right atrium, exhibiting a prominent $x$ descent and absent or blunted $y$ descent); and (3) Muffled, Distant Heart Sounds on cardiac auscultation (caused by the acoustic dampening effect of the insulated layer of fluid surrounding the heart within the pericardial sac). While Beck's triad is pathognomonic when fully present, all three components are detected concurrently in only approximately 30--40\% of confirmed tamponade cases. Its absence does not exclude the diagnosis, especially in subacute or low-pressure tamponade; emergency bedside echocardiography (POCUS) is the definitive modality of choice, followed by immediate needle pericardiocentesis.

\medterm{Beta-Adrenergic Blocker}-- A class of cardiovascular therapeutics that competitively antagonize endogenous catecholamines (epinephrine and norepinephrine) at $\beta$-adrenergic receptor subtypes ($\beta_1$, $\beta_2$, and $\beta_3$). Class II antiarrhythmics in the Vaughan Williams system. Cardioselective agents (metoprolol tartrate/succinate, bisoprolol, atenolol, nebivolol) selectively block $\beta_1$-receptors in the myocardium and juxtaglomerular apparatus, decreasing intracellular cAMP and PKA phosphorylation of L-type calcium channels. This produces negative chronotropy (decreased heart rate at the SA node), negative inotropy (reduced contractility), negative dromotropy (delayed AV nodal conduction, prolonging PR interval), and reduced renin secretion (lowering systemic blood pressure). Non-selective agents (propranolol, nadolol, timolol) block both $\beta_1$ and $\beta_2$ receptors, while third-generation agents possess vasodilatory properties via concurrent $\alpha_1$-blockade (carvedilol, labetalol) or endothelial nitric oxide release (nebivolol). In heart failure with reduced ejection fraction (HFrEF), three specific agents (carvedilol, metoprolol succinate, and bisoprolol) constitute a mandatory pillar of guideline-directed medical therapy (GDMT), blunting chronic sympathetic neurohormonal cardiotoxicity, improving ejection fraction, and reducing mortality by 30--35\%. Also indicated for post-myocardial infarction secondary prevention, rate control in atrial fibrillation, chronic stable angina, and aortic dissection (reducing ventricular ejection force $dP/dt$). Contraindications include cardiogenic shock, decompensated acute heart failure, second- or third-degree AV block without a pacemaker, severe symptomatic bradycardia, and severe active reactive airway disease (bronchospasm).

\medterm{Bicuspid Aortic Valve}-- The most common congenital cardiac anomaly, present in 1-2\% of the population. The aortic valve has two leaflets instead of three, predisposing to aortic stenosis (usually by age 50-60), aortic regurgitation, and infective endocarditis. Associated with aortopathy (aortic root dilation and dissection risk) and coarctation of the aorta. Regular echocardiographic surveillance is recommended. Aortic valve replacement and/or aortic root replacement may be required.

\medterm{Bowel ischemia}-- A spectrum of intestinal vascular disorders occurring when mesenteric arterial or venous blood flow is insufficient to satisfy the metabolic requirements of the small or large intestine, resulting in cellular hypoxia, mucosal sloughing, and potentially transmural necrosis. Divided anatomically and clinically into: (1) Mesenteric Ischemia (predominantly small bowel), which may be acute (embolic or thrombotic occlusion of the superior mesenteric artery, or non-occlusive mesenteric ischemia [NOMI], presenting with severe periumbilical pain out of proportion to physical exam) or chronic ('intestinal angina', marked by severe dull postprandial epigastric pain, fear of eating, and marked weight loss); and (2) Ischemic Colitis (colon ischemia, the most common form, ~75\% of all intestinal ischemia), typically occurring in elderly patients with systemic atherosclerosis following transient systemic hypoperfusion, hypotension, or heart failure. Colonic ischemia selectively involves vascular watershed regions with limited collateral blood supply: Griffith's point at the splenic flexure (junction between SMA and IMA territories) and Sudeck's point at the rectosigmoid junction. Ischemic colitis presents acutely with crampy left lower quadrant abdominal pain followed within 24 hours by mild, hematochezia (bloody diarrhea). Abdominal CT shows segmental colonic wall thickening and 'thumbprinting' (submucosal edema and hemorrhage); colonoscopy reveals pale, edematous mucosa with petechiae and ulceration. Most cases of ischemic colitis are transient and resolve with conservative medical therapy (IV fluids, bowel rest, antibiotics), but transmural infarction with peritonitis demands emergent surgical resection.

\medterm{Bradycardia}-- A heart rate <60 beats per minute. May be physiologic (athletes) or pathologic due to sinus node dysfunction, AV block, or medication effects. Symptomatic bradycardia causes fatigue, dizziness, syncope, or heart failure. Evaluation includes ECG, Holter monitoring, and assessment of reversible causes. Treatment: pacing (temporary or permanent) for symptomatic or high-grade AV block.

\medterm{Brugada Syndrome}-- An inherited, autosomal dominant cardiac channelopathy characterized by a high predisposition to polymorphic ventricular tachycardia (PVT), ventricular fibrillation (VF), and sudden cardiac death (SCD) in the complete absence of gross structural heart disease, classically presenting in young adult males (male-to-female ratio 8:1) of Southeast Asian ancestry (where it is historically termed Lai Tai or sudden unexpected nocturnal death syndrome, SUNDS). In approximately 20--25\% of cases, it is caused by loss-of-function mutations in the SCN5A gene on chromosome 3, which encodes the $\alpha$-subunit of the cardiac voltage-gated sodium channel ($Na_v1.5$). The reduction in inward sodium current ($I_{Na}$) creates a marked transmural voltage gradient across the right ventricular outflow tract (RVOT) epicardium during phase 1 and early phase 2 of the action potential (where outward transient potassium current $I_{to}$ is densest), predisposing to phase 2 reentry and VF. The diagnostic hallmark is the Type 1 Brugada ECG pattern: coved-type ST-segment elevation $\ge 2$ mm ($0.2$ mV) followed by a negative T-wave in one or more right precordial leads ($V_1, V_2$), often accompanied by an incomplete right bundle branch block (RBBB) appearance. The ECG abnormality is dynamic and unmasked or exacerbated by fever (febrile illness), high vagal tone (sleep, postprandial), and Class IC sodium channel blockers (ajmaline, flecainide, procainamide, used for diagnostic provocation challenge). Patients present with nocturnal agonal respirations, syncope, or cardiac arrest. The only definitively proven therapy to prevent sudden cardiac death in high-risk patients (those with prior arrest, documented VT, or syncope) is an Implantable Cardioverter-Defibrillator (ICD); quinidine (which blocks $I_{to}$) and RVOT epicardial catheter ablation are used for recurrent electrical storms.

\medterm{Bundle Branch Block}-- Delay or block in the electrical conduction through the bundle branches. Left Bundle Branch Block (LBBB): QRS >120 ms, broad notched R waves in lateral leads (I, V5-V6), deep S waves in V1-V2. Right Bundle Branch Block (RBBB): QRS >120 ms, rSR' pattern in V1-V2, wide S waves in I, V5-V6. RBBB can occur in healthy individuals or with RV strain. LBBB often indicates structural heart disease, cardiomyopathy, or anterior MI. Management: evaluate underlying cause; consider cardiac resynchronization therapy for LBBB with heart failure.

\medterm{Bundle of His}-- A narrow band of specialized cardiac tissue that conducts electrical impulses from the AV node through the fibrous cardiac skeleton to the interventricular septum. It divides into the right and left bundle branches. The bundle of His is the only electrical connection between atria and ventricles through the fibrous skeleton. Damage here causes complete (third-degree) heart block.

\medterm{Calcium Channel Blockers}-- Agents that block L-type calcium channels, reducing calcium influx into cardiac and smooth muscle cells. Dihydropyridines (amlodipine, nifedipine) cause vasodilation and are used for hypertension and angina. Non-dihydropyridines (verapamil, diltiazem) reduce heart rate and contractility, used for rate control in atrial fibrillation and angina. Side effects include peripheral edema, constipation (verapamil), and bradycardia. Non-dihydropyridines are contraindicated in HFrEF.

\medterm{Caput Medusae}-- Dilated periumbilical veins radiating outward from the umbilicus, seen in portal hypertension with recanalization of the umbilical vein. Not strictly a cardiac finding but relevant in right heart failure causing hepatic congestion and portal hypertension. It must be distinguished from caval collateral circulation. The sign is named after the Gorgon Medusa of Greek mythology.

\medterm{Cardiac allograft}-- A donor heart harvested from a brain-dead organ donor, transplanted into a recipient suffering from end-stage, refractory terminal heart failure (such as stage D HFrEF refractory to maximal guideline-directed medical therapy, biventricular failure, or intractable life-threatening ventricular arrhythmias) or complex congenital heart disease. Orthotopic heart transplantation (OHT) is the surgical standard, wherein the diseased native heart is excised and the donor organ is implanted via bicaval and left atrial anastomoses, connecting the donor aorta and pulmonary artery. Because the donor heart is completely surgically denervated (lacking sympathetic and parasympathetic autonomic inputs), it beats at an intrinsically elevated resting heart rate (~90--110 bpm) and does not respond to atropine or carotid sinus massage; heart rate acceleration during exercise depends on circulating adrenal catecholamines rather than neural reflexes. Lifelong multi-drug immunosuppression is mandatory to prevent allograft rejection, typically comprising a calcineurin inhibitor (tacrolimus or cyclosporine), an antimetabolite (mycophenolate mofetil), and corticosteroids. Surveillance for allograft rejection relies on serial transvenous Endomyocardial Biopsies (EMB) of the right ventricular septum. Long-term post-transplant survival is limited by Cardiac Allograft Vasculopathy (CAV, an accelerated, concentric, diffuse obliterative intimal hyperplasia affecting the entire coronary arterial tree, driven by immune-mediated endothelial injury; because the heart is denervated, CAV-induced ischemia is typically silent without angina), malignancy (post-transplant lymphoproliferative disorder [PTLD] driven by EBV), chronic kidney disease (calcineurin inhibitor nephrotoxicity), and opportunistic infections.

\medterm{Cardiac Amyloidosis}-- An infiltrative cardiomyopathy caused by the extracellular deposition of insoluble, misfolded amyloid protein fibrils within the myocardial interstitium, producing progressive ventricular wall stiffening, restrictive ventricular filling physiology, and advanced heart failure. The two predominant forms affecting the heart are: (1) Transthyretin Cardiac Amyloidosis (ATTR), caused by the misfolding of hepatic transthyretin (a tetrameric carrier protein for thyroxine and retinol-binding protein), presenting either as wild-type ATTR (ATTRwt, historically 'senile amyloidosis', predominantly in elderly men $>70$ years) or hereditary / variant ATTR (ATTRv, autosomal dominant mutations in the TTR gene, such as Val122Ile in African Americans); and (2) Light-Chain Amyloidosis (AL), a plasma cell dyscrasia producing monoclonal immunoglobulin light chains with rapid, fulminant cardiac involvement. Pathognomonic diagnostic features include a discordance between marked left ventricular wall thickening ('pseudo-hypertrophy' on echocardiography) and low QRS voltages on 12-lead ECG. Cardiac MRI demonstrates diffuse, circumferential subendocardial or transmural Late Gadolinium Enhancement (LGE) with abnormal myocardial T1 mapping and expanded extracellular volume (ECV). Non-invasive diagnosis of ATTR is definitively established by bone scintigraphy ($^{99\text{m}}\text{Tc}$-PYP, DPD, or HMDP showing grade 2 or 3 cardiac uptake) in conjunction with negative serum/urine immunofixation and normal serum free light chain ratios to definitively exclude AL amyloidosis. Endomyocardial biopsy with Congo red staining displaying apple-green birefringence under polarized light remains the gold standard. Disease staging relies on the Mayo and NAC staging systems using NT-proBNP and troponin levels. Targeted therapeutics for ATTR include the oral TTR tetramer stabilizer tafamidis and TTR gene-silencing therapies (patisiran, vutrisiran, inotersen); AL amyloidosis requires clone-directed chemotherapy (daratumumab, bortezomib, dexamethasone).

\medterm{Cardiac Arrest}-- The sudden, abrupt cessation of effective mechanical cardiac pumping activity, resulting in the instantaneous loss of systemic arterial perfusion, immediate cerebral hypoperfusion, loss of consciousness, and cessation of normal breathing (apnea or agonal gasps), which is fatal within minutes without immediate resuscitation. Electrophysiologically divided into shockable rhythms (Ventricular Fibrillation [VF] and Pulseless Ventricular Tachycardia [pVT], driven by acute coronary ischemia and responsive to immediate electrical defibrillation) and non-shockable rhythms (Asystole and Pulseless Electrical Activity [PEA], driven by reversible systemic causes summarized by the 'Hs and Ts': Hypovolemia, Hypoxia, Hydrogen ion acidosis, Hypo/hyperkalemia, Hypothermia, Tension pneumothorax, Tamponade, Toxins, and Thrombosis). The Chain of Survival emphasizes immediate bystander recognition, rapid EMS activation, high-quality CPR (compression rate 100--120/min, depth 5--6 cm, full chest recoil), early defibrillation with automated external defibrillators (AEDs), and advanced post-resuscitation care. Globally, Out-of-Hospital Cardiac Arrest (OHCA) carries a survival-to-hospital-discharge rate of approximately 8--12\%, heavily dependent on witnessed arrest and immediate bystander CPR. Standardized reporting follows the Utstein guidelines. Post-resuscitation care focuses on early coronary angiography, hemodynamic stabilization, and targeted temperature management (32--36$^\circ$C) to mitigate ischemic-reperfusion brain injury, followed by multimodal neuroprognostication at 72 hours.

\medterm{Cardiac Axis}-- The mean net vector and spatial direction of electrical depolarization through the ventricular myocardium in the frontal plane during the cardiac cycle, calculated from the hexaxial reference system of the six limb leads (I, II, III, aVR, aVL, aVF) on a standard 12-lead ECG. The QRS complex represents the summation of all ventricular depolarization vectors, with the thick left ventricle dominating the net electrical signal. Normal cardiac axis lies between $-30^\circ$ and $+90^\circ$ (leads I and II are both predominantly positive). Deviations are classified as: (1) Left Axis Deviation (LAD, between $-30^\circ$ and $-90^\circ$): lead I is positive, while leads II and aVF are negative; etiologies include left anterior fascicular block (LAFB), left bundle branch block (LBBB), left ventricular hypertrophy (LVH), inferior wall MI, and mechanical elevation of the diaphragm (pregnancy, ascites, obesity); (2) Right Axis Deviation (RAD, between $+90^\circ$ and $+180^\circ$): lead I is negative, while lead aVF is positive; normal variant in tall, thin individuals and neonates, but pathological in right ventricular hypertrophy (RVH), acute pulmonary embolism (cor pulmonale), chronic COPD, left posterior fascicular block (LPFB), and lateral wall MI; and (3) Extreme Axis Deviation ('Northwest' or indeterminate axis, between $-90^\circ$ and $\pm 180^\circ$): both leads I and aVF are predominantly negative; highly suggestive of ventricular tachycardia (VT, strongly favoring VT over SVT with aberrancy), apical ventricular pacing, or severe hyperkalemia.

\medterm{Cardiac Catheterization}-- Invasive procedure involving insertion of catheters into the heart via peripheral arteries or veins. Diagnostic catheterization measures intracardiac pressures, saturations, and performs coronary angiography. Indications include evaluation of CAD, valvular disease, cardiomyopathy, and congenital heart disease. Therapeutic catheterization includes PCI, valvuloplasty, and device closure. Complications include bleeding, vessel injury, arrhythmias, and contrast-induced nephropathy.

\medterm{Cardiac Conduction System}-- The specialized tissue responsible for generating and conducting electrical impulses through the heart. Components include the sinoatrial (SA) node, internodal pathways, atrioventricular (AV) node, bundle of His, right and left bundle branches, and Purkinje fibers. The system ensures coordinated contraction: atria contract first, followed by ventricles, with appropriate delay at the AV node.

\medterm{Cardiac CT Angiography}-- Noninvasive imaging using CT to visualize coronary arteries and cardiac structures. It has high negative predictive value for excluding coronary artery disease, making it useful for low-to-intermediate risk patients with chest pain. Calcium scoring quantifies coronary artery calcification as a marker of atherosclerosis. Limitations include heart rate control requirements, radiation exposure, and inability to assess hemodynamic significance of stenosis. It also evaluates aortic and congenital heart disease.

\medterm{Cardiac Cycle}-- The sequence of events occurring during one complete heartbeat, consisting of systole (contraction) and diastole (relaxation). During diastole, the ventricles fill with blood; during systole, they contract and eject blood. The cardiac cycle includes isovolumetric relaxation, ventricular filling, isovolumetric contraction, and ventricular ejection phases. The cycle takes approximately 0.8 seconds at a heart rate of 75 bpm.

\medterm{Cardiac Index}-- Cardiac output normalized to body surface area (BSA), calculated as CO / BSA. Normal cardiac index is 2.5-4.2 L/min/m2. It provides a more accurate assessment of cardiac function than cardiac output alone, accounting for body size. A cardiac index below 2.2 L/min/m2 indicates cardiogenic shock. Cardiac index is used to guide therapy in critically ill patients.

\medterm{Cardiac Metastases}-- More common than primary cardiac tumors. Most frequent sources: lung cancer, breast cancer, melanoma, lymphoma, and renal cell carcinoma. Most commonly involve the pericardium (effusion, tamponade), but can affect myocardium and endocardium. Presentation depends on location: pericardial effusion, heart failure, arrhythmias, or embolic events. Diagnosis by echocardiography, cardiac MRI, or CT. Treatment focuses on systemic cancer therapy; cardiac surgery rarely indicated.

\medterm{Cardiac MRI}-- Cardiovascular magnetic resonance (CMR) imaging, the acknowledged non-invasive reference gold standard for the comprehensive anatomical, functional, and tissue-characterization assessment of the heart and thoracic vessels, operating without exposure to ionizing radiation. CMR utilizes powerful magnetic fields (1.5 or 3.0 Tesla) and radiofrequency pulses to generate cine balanced steady-state free precession (bSSFP) images, providing unmatched spatial resolution and reproducibility for quantifying left and right ventricular volumes, wall mass, stroke volume, and ejection fraction (LVEF and RVEF), completely independent of geometric assumptions or poor acoustic windows that limit echocardiography. Its most revolutionary capability is non-invasive myocardial tissue characterization: (1) Late Gadolinium Enhancement (LGE): intravenous gadolinium-based contrast accumulates in expanded extracellular spaces (fibrosis, necrosis, scar); the pattern of LGE distinguishes ischemic scar (subendocardial or transmural enhancement conforming to a coronary vascular territory) from non-ischemic cardiomyopathies (mid-wall enhancement in dilated cardiomyopathy, patchy mid-wall/epicardial enhancement in hypertrophic cardiomyopathy, patchy subepicardial LGE in acute myocarditis, and diffuse subendocardial LGE in cardiac amyloidosis); (2) Parametric T1 and T2 Mapping: quantifies diffuse interstitial myocardial fibrosis, extracellular volume fraction (ECV, massively elevated in amyloidosis), and myocardial edema/inflammation (elevated T2 in acute myocarditis or acute rejection); and (3) Phase-Contrast Velocity Encoding: quantifies transvalvular regurgitant fractions and shunt ratios ($Q_p/Q_s$).

\medterm{Cardiac Output Measurement}-- Methods include: thermodilution via pulmonary artery catheter (bolus of cold saline and temperature change measurement), Fick principle (O2 consumption / arteriovenous O2 difference), echocardiographic Doppler (flow velocity time integral x valve area x heart rate), and non-invasive cardiac output monitors (bioimpedance, bioreactance). CO measurement guides management of heart failure, shock, and perioperative hemodynamic optimization.

\medterm{Cardiac PET Imaging}-- Positron emission tomography using radiotracers (Rb-82, N-13 ammonia, O-15 water) for quantitative myocardial perfusion assessment. Advantages over SPECT: superior resolution, shorter scan time, attenuation correction, and absolute myocardial blood flow quantification (mL/g/min). FDG PET identifies hibernating myocardium (viable but ischemic tissue that may benefit from revascularization). Used for viability assessment, sarcoidosis evaluation, and cardiac infection imaging.

\medterm{Cardiac Radiation Injury}-- Late complication of mediastinal radiation therapy (Hodgkin lymphoma, breast cancer, lung cancer). Manifestations include pericardial disease (constrictive pericarditis, effusion), coronary artery disease (accelerated atherosclerosis), valvular disease (fibrosis), myocardial fibrosis, and conduction abnormalities. Risk increases with higher radiation doses and larger treatment fields. Long-term surveillance with echocardiography is recommended for survivors.

\medterm{Cardiac Rehabilitation}-- A medically supervised, comprehensive, multidisciplinary secondary prevention program designed to optimize cardiovascular health, improve functional exercise capacity, and reduce recurrent morbidity and mortality in patients with established cardiovascular disease. Core components encompass: (1) Structured, individualized exercise training (combining aerobic conditioning and resistance training with continuous or telemetry ECG monitoring); (2) Aggressive cardiovascular risk factor modification (lipid management, blood pressure control, smoking cessation, and diabetes optimization); (3) Nutritional counseling and dietary guidance; (4) Medication adherence education; and (5) Psychosocial assessment and support for depression, anxiety, and stress. Cardiac rehab is a guideline-mandated Class I recommendation following acute myocardial infarction, percutaneous coronary intervention (PCI), coronary artery bypass graft (CABG) surgery, heart valve repair or replacement, stable chronic heart failure (HFrEF), heart transplantation, and peripheral artery disease. Meta-analyses demonstrate that participation reduces all-cause mortality by approximately 13--15\% and cardiovascular mortality by 26--30\%, with marked reductions in unplanned hospital readmissions. Despite proven clinical benefits, global and national participation rates remain sub-optimal (20--30\%), driving recent implementation of virtual, home-based, and mobile health cardiac rehabilitation platforms.

\medterm{Cardiac Resynchronization Therapy}-- Biventricular pacing that synchronizes contraction of the left and right ventricles in patients with heart failure, reduced EF, and prolonged QRS duration (especially LBBB). CRT improves symptoms, exercise capacity, and quality of life, and reduces heart failure hospitalizations and mortality. Response may be incomplete in 30\% of patients. Selection criteria include EF 35\% or less, NYHA II-IV despite optimal medical therapy, and QRS 150 ms or greater.

\medterm{Cardiac Sarcoidosis}-- Infiltrative disease characterized by non-caseating granulomas in the myocardium, affecting 25-30\% of patients with systemic sarcoidosis. Presentation includes conduction abnormalities (AV block), ventricular arrhythmias, and heart failure. Diagnosis requires cardiac MRI showing late gadolinium enhancement (typically basal septal and epicardial), PET showing active inflammation, or endomyocardial biopsy. Treatment: corticosteroids for active inflammation, ICD for arrhythmia prevention.

\medterm{Cardiac Skeleton}-- A dense connective tissue framework composed of four fibrous rings that surround the heart valves and provide structural support. The cardiac skeleton includes the annuli fibrosi (valve rings), the membranous septum, and the left and right fibrous trigones. It serves as electrical insulation between atria and ventricles, forcing impulses through the AV node, and provides attachment points for valve leaflets and cardiac muscle.

\medterm{Cardiac Stress Testing}-- Assessment of cardiac function under physiological or pharmacological stress to detect ischemia. Exercise stress testing involves treadmill or bicycle exercise with continuous ECG monitoring. Pharmacological stress (dobutamine, adenosine, regadenoson) is used when patients cannot exercise. Stress testing combined with imaging (nuclear perfusion, echocardiography, cardiac MRI) improves sensitivity and specificity. Indications include evaluation of chest pain, preoperative risk assessment, and prognosis.

\medterm{Cardiac Tamponade Physiology}-- A life-threatening cardiovascular emergency characterized by acute or subacute compression of all four cardiac chambers caused by the abnormal accumulation of fluid, pus, blood, or gas within the non-compliant, fibroelastic pericardial space under elevated intrapericardial pressure. The severity of hemodynamic compromise depends not merely on absolute effusion volume, but on the rate of accumulation: a rapid accumulation of as little as 150--200 mL (e.g., myocardial free-wall rupture, aortic dissection type A, cardiac trauma) causes catastrophic tamponade, whereas slow, chronic accumulation (malignancy, uremia, tuberculosis) allows the pericardium to stretch, accumulating $>1.5--2.0$ L before pressures rise precipitously. As intrapericardial pressure rises to equal and exceed normal right-sided diastolic filling pressures, diastolic ventricular filling is progressively impaired. The heart becomes volume-constrained, leading to chamber collapse and equalisation of end-diastolic pressures across all four chambers (RAP $\approx$ RVEDP $\approx$ PCWP $\approx$ LVEDP). Stroke volume and cardiac output plummet. The clinical hallmark is Beck's Triad: (1) Systemic arterial hypotension; (2) Jugular venous distension (elevated JVP with prominent $x$ descent and absent $y$ descent due to arrested early diastolic filling); and (3) Muffled, distant heart sounds. Pathophysiologically, enhanced ventricular interdependence produces Pulsus Paradoxus (an inspiratory drop in systolic blood pressure $>10$ mmHg: during inspiration, increased venous return expands the right ventricle within the rigid pericardium, bowing the interventricular septum leftward and compressing the LV, acutely reducing LV stroke volume). Diagnostic ECG shows sinus tachycardia, low QRS voltage, and Electrical Alternans (beat-to-beat variation in QRS amplitude caused by the heart swinging pendulum-like within the fluid). Echocardiography confirms early diastolic right ventricular free-wall collapse (specific), late systolic right atrial collapse (sensitive), and plethoric IVC with no inspiratory collapse. Definitive emergency treatment is immediate percutaneous pericardiocentesis or surgical creation of a pericardial window.

\medterm{Cardiac Tumors}-- Neoplasms arising within or involving the heart, categorized into primary cardiac tumors (originating within cardiac tissues, rare, incidence $<0.1\%$) and secondary metastatic tumors (20 to 40 times more common, metastasizing from lung carcinoma, breast carcinoma, malignant melanoma, lymphoma, or renal cell carcinoma). Approximately 75--80\% of primary cardiac tumors are histologically benign: (1) Atrial Myxoma: the most common primary tumor in adults (~50\%), typically solitary, pedunculated, gelatinous masses arising from the fossa ovalis in the left atrium (75\%), presenting with intracardiac mechanical obstruction ('tumor plop' murmur mimicking mitral stenosis), systemic embolization (stroke), and constitutional symptoms (fever, weight loss, elevated ESR from IL-6 secretion); (2) Rhabdomyoma: the most common pediatric cardiac tumor, strongly associated with tuberous sclerosis (TSC1/TSC2 mutations), typically regressing spontaneously; (3) Fibroma: intramural ventricular masses associated with Gorlin syndrome and ventricular arrhythmias; and (4) Papillary fibroelastoma: small, sea-anemone-like benign growths on valvular endocardium carrying high stroke risk. Malignant primary tumors account for 20--25\%, predominantly aggressive sarcomas (Angiosarcoma, most common in the right atrium with early metastases and hemopericardium, and undifferentiated pleomorphic sarcoma). Diagnostic imaging comprises echocardiography (TTE/TEE) and cardiac MRI for tissue characterization (T1/T2 signal and late gadolinium enhancement). Benign symptomatic tumors are treated with urgent surgical resection; malignant tumors have a dismal prognosis requiring multimodal surgical and oncological care.

\medterm{Cardiogenic Shock}-- A critical, life-threatening clinical and hemodynamic state characterized by severe end-organ tissue hypoperfusion resulting primarily from cardiac pump failure, despite adequate or elevated intravascular volume (ventricular preload). Diagnostic hemodynamic criteria established by the SHOCK trial strictly mandate: (1) Persistent systemic arterial hypotension (systolic $\text{BP} < 90$ mmHg for $>30$ minutes or requirement for vasopressors to maintain $\text{MAP} \ge 65$ mmHg); (2) Severely depressed Cardiac Index ($\text{CI} < 2.2$ L/min/m$^2$, or $<1.8$ L/min/m$^2$ without pharmacological/mechanical support); and (3) Elevated ventricular filling pressures (Pulmonary Capillary Wedge Pressure $\text{PCWP} > 15$ mmHg and/or $\text{CVP} > 10$ mmHg). Acute Myocardial Infarction with extensive left ventricular necrosis (loss of $>40\%$ of functioning LV myocardium) accounts for $>80\%$ of cases; other etiologies include acute mechanical complications of MI (ventricular septal rupture, papillary muscle rupture with flail mitral leaflet, free wall rupture), fulminant acute myocarditis, decompensated end-stage cardiomyopathy, severe acute aortic or mitral regurgitation, and sustained ventricular tachyarrhythmias. Clinical features reflect systemic hypoperfusion (cool, clammy extremities, oliguria $<0.5$ mL/kg/h, altered mental status, and profound lactic acidosis) accompanied by pulmonary congestion (bilateral rales, acute pulmonary edema). Clinical severity is stratified using the SCAI staging classification (Stages A through E: At risk, Beginning, Classic, Deteriorating, and Extremis). Emergency management requires immediate invasive hemodynamic monitoring (arterial line and pulmonary artery catheter), inotropic support (dobutamine or milrinone to augment contractility), first-line vasopressors (norepinephrine to restore coronary perfusion pressure), and temporary Mechanical Circulatory Support (MCS), including intra-aortic balloon pump (IABP), percutaneous microaxial LV pumps (Impella), and venoarterial extracorporeal membrane oxygenation (VA-ECMO). In ischemic cardiogenic shock, emergency revascularization via primary PCI or emergency CABG is the only definitive intervention proven to reduce mortality (SHOCK trial).

\medterm{Cardiology}-- The branch of medicine devoted to the study, diagnosis, and treatment of disorders of the heart and circulatory system. Cardiologists manage conditions including coronary artery disease, heart failure, valvular heart disease, arrhythmias, cardiomyopathy, pericardial disease, congenital heart disease, and hypertension. Diagnostic tools: electrocardiography, echocardiography, cardiac catheterization, stress testing, cardiac MRI, CT angiography, and nuclear cardiac imaging. Interventional cardiology performs percutaneous coronary intervention and structural heart procedures (TAVR, MitraClip). Electrophysiology manages rhythm disorders with ablation and device therapy (pacemakers, ICDs). Preventive cardiology focuses on risk factor modification for cardiovascular disease.

\medterm{Cardiomegaly}-- Abnormal enlargement of the heart, identified clinically, on chest radiography, or with cardiac imaging. It may reflect chamber dilatation, myocardial hypertrophy, pericardial effusion, or a combination of these; causes include hypertension, valvular disease, cardiomyopathy, congenital heart disease, and high-output states.

\medterm{Cardiomyopathy Classification}-- The MOGE(S) classification system provides phenotypic and genotypic description: M (morpho-functional phenotype), O (organ/system involvement), G (genetic inheritance), E (etiology), S (functional status/NYHA class). This replaces older descriptive classifications and provides more precise characterization for prognosis and family counseling. Each cardiomyopathy can be classified as primary (genetic or acquired) or secondary (due to systemic disease).

\medterm{Cardio-Oncology}-- The subspecialty focusing on cardiovascular care of cancer patients, addressing cancer therapy-related cardiac dysfunction, prevention, and management. Key concerns include anthracycline cardiotoxicity, trastuzumab-related cardiomyopathy, immune checkpoint inhibitor myocarditis, and radiation-induced heart disease. Monitoring includes baseline and serial echocardiography during and after cancer treatment. Early detection and intervention can prevent irreversible cardiac damage and allow continuation of cancer therapy.

\medterm{Cardiopulmonary Exercise Testing}-- Comprehensive assessment of cardiovascular and pulmonary function during graded exercise on a treadmill or cycle ergometer. Measures: VO2 max, anaerobic threshold, ventilatory efficiency (VE/VCO2 slope), and hemodynamic responses. Used to evaluate dyspnea on exertion, determine heart failure prognosis (VO2 max <14 mL/kg/min indicates need for transplant evaluation), assess exercise capacity for disability evaluation, and guide cardiac rehabilitation programs.

\medterm{Cardiopulmonary Resuscitation}-- A sequential, evidence-based series of emergency interventions (Basic Life Support [BLS] and Advanced Cardiovascular Life Support [ACLS]) designed to maintain vital organ perfusion and oxygen delivery during cardiac arrest until the restoration of spontaneous circulation (ROSC). High-quality chest compressions form the bedrock of CPR: compressions must be delivered at a rate of 100 to 120 per minute, at a depth of 5 to 6 cm (2 to 2.4 inches) in adults, allowing complete chest recoil between compressions while minimizing interruptions ($>60\%$ chest compression fraction). Compressions create an artificial cardiac output (~25--30\% of normal) primarily by increasing intrathoracic pressure ('thoracic pump' mechanism) and direct cardiac compression ('cardiac pump' mechanism), generating a coronary perfusion pressure (CPP = aortic diastolic pressure minus right atrial pressure) critical for myocardial survival ($\text{CPP} \ge 15$ mmHg is required for ROSC). In un-intubated adults, compressions are paired with ventilations at a 30:2 ratio; once an advanced airway is placed, continuous compressions are delivered with asynchronous positive-pressure ventilations (1 breath every 6 seconds, 10 breaths/min, avoiding hyperventilation which increases intrathoracic pressure and reduces venous return). Resuscitation follows rhythm-specific ACLS algorithms: shockable rhythms (Ventricular Fibrillation [VF] and Pulseless Ventricular Tachycardia [pVT]) require immediate unsynchronized defibrillation, CPR, and IV epinephrine (1 mg every 3--5 min) and amiodarone (300 mg then 150 mg); non-shockable rhythms (Asystole and Pulseless Electrical Activity [PEA]) require high-quality CPR, early epinephrine, and aggressive identification and correction of reversible causes (the 5 H's: Hypovolemia, Hypoxia, Hydrogen ion/acidosis, Hypo/hyperkalemia, Hypothermia; and 5 T's: Tension pneumothorax, Tamponade, Toxins, Thrombosis pulmonary, Thrombosis coronary).

\medterm{Cardiorenal Syndrome}-- Dysfunction of the heart and kidneys where acute or chronic failure of one organ contributes to dysfunction of the other. Type 1: acute cardiac dysfunction causing acute kidney injury (e.g., cardiogenic shock). Type 2: chronic cardiac dysfunction causing progressive CKD. Type 3: acute kidney injury causing cardiac dysfunction. Type 4: chronic kidney disease causing cardiac dysfunction. Type 5: systemic conditions affecting both organs. Management is challenging and requires coordinated care.

\medterm{Cardiotoxic Drug Overdose}-- Acute poisonings and intentional or accidental toxic ingestions involving pharmaceuticals that directly perturb cardiac electrophysiology, myocardial contractility, or vascular tone, representing a critical subset of medical emergencies in intensive care and cardiovascular medicine. Major drug classes and targeted antidotal management include: (1) Beta-Adrenergic Blockers: toxicity manifests with profound bradycardia, hypotension, heart block, and cardiogenic shock; management requires intravenous atropine, high-dose intravenous glucagon (which activates adenylate cyclase directly via G-protein coupling bypassing $\beta$-receptors), High-Dose Insulin Euglycemia Therapy (HIET, providing potent inotropic support via improved myocyte glucose uptake), and intravenous lipid emulsion (ILE) for lipophilic agents; (2) Calcium Channel Blockers (verapamil, diltiazem): causes refractory vasodilation and bradyarrhythmias; treated with intravenous calcium chloride/gluconate, HIET, vasopressors (norepinephrine, epinephrine), and ILE; (3) Digoxin: cardiac glycoside toxicity causes hyperkalemia and polymorphic arrhythmias (bidirectional ventricular tachycardia, PVCs, junctional rhythms); definitive antidote is Digoxin-Specific Antibody Fragments (DigiFab); (4) Tricyclic Antidepressants (TCAs, amitriptyline): causes voltage-gated sodium channel blockade producing characteristic wide QRS complexes ($\ge 100$ ms) and prominent terminal $R$ wave in lead $aVR$, predisposing to ventricular arrhythmias and seizure; immediate treatment is hypertonic intravenous Sodium Bicarbonate (1--2 mEq/kg IV push) to alkalinize plasma (target pH 7.50--7.55) and flood channels with sodium ions; and (5) Potassium channel blockers (QT-prolonging drugs) inducing Torsades de Pointes, treated with intravenous magnesium sulfate.

\medterm{Cardiotoxicity from Chemotherapy}-- Chemotherapy agents that damage the myocardium, most commonly anthracyclines (doxorubicin, daunorubicin). Cardiotoxicity presents as dilated cardiomyopathy with reduced ejection fraction. Risk is dose-dependent (doxorubicin cumulative dose >550 mg/m2). Other cardiotoxic agents include trastuzumab (reversible cardiomyopathy), cyclophosphamide, and 5-fluorouracil. Monitoring with serial echocardiography during treatment is essential. Treatment includes ACE inhibitors and beta-blockers.

\medterm{Cardiovascular Disease}-- A broad umbrella designation encompassing all pathological disorders affecting the heart and systemic vascular circulation, representing the single leading cause of global morbidity and mortality worldwide (responsible for approximately 18--19 million deaths annually, accounting for $>32\%$ of all global deaths according to World Health Organization estimates). Major clinical classifications include: (1) Coronary Artery Disease (CAD / ischemic heart disease), comprising stable angina, unstable angina, and acute myocardial infarction; (2) Cerebrovascular Disease, encompassing ischemic stroke, transient ischemic attack (TIA), and intracranial hemorrhage; (3) Peripheral Artery Disease (PAD), presenting with claudication and critical limb-threatening ischemia; (4) Heart Failure (HFrEF and HFpEF); (5) Valvular Heart Disease (rheumatic valve disease, calcific aortic stenosis, mitral valve prolapse); (6) Cardiomyopathies (dilated, hypertrophic, restrictive, arrhythmogenic); and (7) Congenital Heart Anomalies. Foundational modifiable cardiovascular risk factors include systemic arterial hypertension (the leading attributable risk factor for premature mortality), atherogenic dyslipidemia (elevated LDL-C, apolipoprotein B, and lipoprotein(a)), cigarette smoking, type 2 diabetes mellitus, metabolic syndrome, central obesity, sedentary lifestyle, and chronic kidney disease. Primary and secondary prevention hinges on guideline-directed lifestyle modifications alongside targeted pharmacotherapies (high-intensity statins, PCSK9 inhibitors, antihypertensives, SGLT2 inhibitors, GLP-1 receptor agonists, and antiplatelet agents).

\medterm{Cardiovascular Health}-- Cardio is shorthand for cardiovascular health or exercise that elevates heart rate and respiratory effort to strengthen the heart, lungs, and circulatory system. Aerobic activities such as walking, running, cycling, and swimming improve cardiac output, lower blood pressure, and enhance metabolic efficiency. Regular cardiovascular exercise reduces the risk of heart disease, stroke, and diabetes while improving overall longevity.

\medterm{Cardiovascular System}-- The organ system comprising the heart and blood vessels (arteries, arterioles, capillaries, venules, veins) responsible for circulating blood throughout the body. Functions: delivery of oxygen and nutrients to tissues, removal of carbon dioxide and metabolic wastes, transport of hormones and immune cells, thermoregulation, and maintenance of fluid and electrolyte homeostasis. The heart is a four-chambered muscular pump (two atria, two ventricles) that generates the pressure gradient for blood flow. The vascular system forms a closed circuit: systemic circulation (oxygenated blood from left heart to body, deoxygenated blood returns to right heart) and pulmonary circulation (deoxygenated blood from right heart to lungs, oxygenated blood returns to left heart).

\medterm{Carotid Bruit}-- An audible systolic (sometimes continuous) murmur over the carotid artery due to turbulent flow, most commonly from atherosclerotic carotid stenosis. A carotid bruit raises concern for significant stenosis (>=50 percent), which confers increased stroke risk. However, sensitivity and specificity are limited; the finding prompts confirmation with carotid duplex ultrasound. Management depends on stenosis severity and symptoms: aggressive risk factor modification, antiplatelets, and carotid endarterectomy or stenting for symptomatic high-grade stenosis.

\medterm{Central Venous Pressure}-- Pressure in the superior vena cava or right atrium, measured via central venous catheter. Normal: 2-8 mmHg. Elevated CVP indicates right heart failure, fluid overload, cardiac tamponade, or constrictive pericarditis. Low CVP indicates hypovolemia. CVP monitoring guides fluid resuscitation and diuretic therapy in heart failure management. Less invasive alternatives include bedside ultrasound assessment of IVC diameter and collapsibility.

\medterm{Chagas Disease}-- A parasitic infection caused by Trypanosoma cruzi, transmitted by reduviid bugs (kissing bugs), endemic in Latin America. Acute phase: fever, swelling at bite site (chagoma), Romana sign (periorbital edema). Chronic phase (years later): Chagas cardiomyopathy (most common and serious), megaesophagus, and megacolon. Cardiomyopathy: progressive heart failure, apical aneurysm, thrombus formation, heart block, and sudden cardiac death. Diagnosis: serology (ELISA, IFA), PCR in acute phase. Treatment: benznidazole or nifurtimox for acute and early chronic phase. Late chronic cardiomyopathy: symptomatic heart failure management, anticoagulation, and pacemaker/ICD as indicated.

\medterm{Chest Pain}-- A cardinal symptom with cardiac and non-cardiac causes requiring urgent distinction. Cardiac: acute coronary syndrome (angina, MI), pericarditis, aortic dissection, myocarditis. Pulmonary: pulmonary embolism, pneumothorax, pneumonia, pleurisy. GI: GERD, esophageal spasm, peptic ulcer. Musculoskeletal and psychogenic causes also common. Evaluation: history (character, radiation, aggravating/relieving factors, risk factors), ECG, troponin, chest X-ray, and risk stratification; rule-out algorithms (HEART score, serial troponins) guide disposition.

\medterm{Cholesterol and Lipids}-- Fatty substances in the blood, including LDL cholesterol, HDL cholesterol, and triglycerides, that play critical roles in cellular function and energy metabolism. Abnormal lipid levels, particularly high LDL and triglycerides with low HDL, promote atherosclerosis and increase the risk of heart attack and stroke. Management includes dietary modification, exercise, and statins or other lipid-lowering medications.

\medterm{Cholesterol crystal embolization}-- A systemic, multi-organ atheroembolic syndrome (also termed atheroembolism or blue toe syndrome) caused by the dislodgement and showering of microscopic cholesterol crystal fragments and atheromatous debris from eroded, ulcerated atherosclerotic plaques in the aorta or major arteries into the downstream microcirculation (arterioles 100 to 200 $\mu$m in diameter). Most commonly precipitated iatrogenically following invasive vascular procedures (cardiac catheterization, aortic endovascular aneurysm repair, percutaneous coronary intervention) or vascular surgery, and less frequently following the initiation of anticoagulation (which dissolves protective fibrin mesh over ulcerated plaques) or thrombolytic therapy; occasionally occurs spontaneously. Embolized biconvex cholesterol crystals physically lodge in distal arterioles, inciting an intense foreign-body reaction characterized by endothelial proliferation, perivascular infiltration of eosinophils and macrophages, and progressive intraluminal obliterative fibrosis. The clinical hallmark is a triad of livedo reticularis (a purplish, reticulated, net-like cutaneous pattern on the legs and trunk), 'blue toe syndrome' (painful, purpuric, cyanotic digital ischemia with preserved peripheral pulses), and subacute progressive renal failure (atheroembolic renal disease, characteristically worsening weeks after catheterization). Other manifestations include intestinal ischemia (abdominal pain, gastrointestinal bleeding), retinal Hollenhorst plaques (refractile golden cholesterol crystals at arteriolar bifurcations on fundoscopy), cerebral infarction, and pancreatitis. Laboratory clues include peripheral blood eosinophilia ($>80\%$), hypocomplementemia (low C3/C4), elevated ESR, and urine eosinophils (Hansel stain). Definitive diagnosis is established by skin, renal, or muscle biopsy demonstrating pathognomonic 'ghost cells' or biconvex needle-shaped clefts where cholesterol crystals dissolved during histology preparation. Treatment is supportive; anticoagulation is strictly contraindicated as it promotes further atheroembolism.

\medterm{Chronic Venous Insufficiency}-- Impaired venous return from the lower extremities due to venous valve incompetence or obstruction. Leads to venous hypertension, edema, skin changes (lipodermatosclerosis, hyperpigmentation), and venous stasis ulcers (medial malleolus). Risk factors: varicose veins, DVT, obesity, prolonged standing. Diagnosis: clinical, venous duplex ultrasound. Treatment: compression therapy (graded stockings), leg elevation, exercise, and wound care. Severe cases: endovenous ablation, venous stenting, or surgical reconstruction.

\medterm{Circumflex Artery}-- A branch of the left coronary artery that travels in the left atrioventricular groove. It supplies the lateral wall of the left ventricle, the left atrium, and in left-dominant systems, the inferior wall and posteromedial papillary muscle. Circumflex artery occlusion causes lateral wall myocardial infarction, which may be missed on standard ECG leads.

\medterm{Class IA Antiarrhythmics}-- Sodium channel blockers that slow phase 0 depolarization and prolong the action potential duration and refractory period. Examples include quinidine, procainamide, and disopyramide. Used for atrial and ventricular arrhythmias. Quinidine can cause QT prolongation and torsades de pointes. Procainamide may cause lupus-like syndrome with long-term use. These drugs are generally reserved for cases not controlled by first-line agents due to proarrhythmic potential.

\medterm{Class IB Antiarrhythmics}-- Sodium channel blockers that shorten the action potential duration and are preferentially taken up by ischemic tissue. Examples include lidocaine and mexiletine. Lidocaine is used for ventricular tachycardia and ventricular fibrillation, particularly in the setting of acute myocardial infarction. It has minimal effect on normal myocardium. Side effects include CNS toxicity (seizures, confusion) and cardiac depression.

\medterm{Class IC Antiarrhythmics}-- Sodium channel blockers that significantly slow phase 0 depolarization without affecting action potential duration. Examples include flecainide and propafenone. Used for supraventricular arrhythmias in patients without structural heart disease. The CAST trial showed increased mortality in post-MI patients, so they are contraindicated in that setting. Flecainide is effective for atrial fibrillation and atrial flutter in structurally normal hearts.

\medterm{Class III Antiarrhythmics}-- Potassium channel blockers that prolong the action potential duration and refractory period. Examples include amiodarone, sotalol, dofetilide, and ibutilide. Amiodarone is the most commonly used but has significant toxicity. Sotalol has additional beta-blocking properties and can cause torsades de pointes. Dofetilide requires inpatient initiation with QT monitoring. These drugs are used for atrial fibrillation and ventricular arrhythmias.

\medterm{Coarctation of the Aorta}-- A congenital narrowing of the aorta, typically near the ductus arteriosus. Presents with upper extremity hypertension, diminished femoral pulses, and a blood pressure gradient between arms and legs. Associated with bicuspid aortic valve and other congenital heart defects. Repair is via surgical resection or balloon angioplasty with stent placement.

\medterm{Cocaine Cardiotoxicity}-- Cocaine blocks norepinephrine and dopamine reuptake, causing sympathomimetic effects: tachycardia, hypertension, and coronary vasospasm. Acute myocardial infarction can occur even with normal coronary arteries due to vasospasm, accelerated atherosclerosis, and prothrombotic effects. Treatment: benzodiazepines for sympatholysis. Beta-blockers are contraindicated (unopposed alpha stimulation worsens vasospasm). Calcium channel blockers and nitroglycerin for vasospasm.

\medterm{Cogan’s syndrome}-- A rare, chronic systemic autoimmune chronic inflammatory disorder characterized primarily by the combination of ocular inflammatory disease and inner ear vestibulo-auditory dysfunction, frequently classified as a systemic variable-vessel vasculitis. Classically presents in young adults with bilateral interstitial keratitis (causing eye pain, photophobia, lacrimation, and conjunctival injection) closely associated with sudden-onset Meniere-like vestibulo-auditory symptoms: severe rotational vertigo, ataxia, nausea, vomiting, tinnitus, and rapidly progressive, profound bilateral sensorineural hearing loss that can culminate in permanent deafness within weeks to months. Crucially, approximately 10--15\% of patients develop severe cardiovascular manifestations driven by large-vessel systemic vasculitis: chronic aortitis affecting the ascending thoracic aorta and aortic root, causing severe, progressive aortic regurgitation (requiring urgent aortic valve replacement or root repair), thoracic and abdominal aortic aneurysms, coronary ostial stenosis (provoking myocardial ischemia or infarction in young patients), and systemic vasculitis mimicking Takayasu arteritis. Pathogenesis involves autoimmune responses directed against inner ear and corneal autoantigens (such as SSA/Ro, cochlin, and heat shock protein 70). Early aggressive immunosuppression with high-dose systemic corticosteroids and second-line biologic agents (such as TNF inhibitors [infliximab] or tocilizumab) is essential to prevent permanent profound deafness and catastrophic cardiovascular complications.

\medterm{Collateral Circulation}-- Alternative vascular pathways that develop when coronary arteries become chronically narrowed. Collaterals connect different coronary artery territories and can partially compensate for stenosis. They develop gradually in response to ischemia and are more prominent in patients with chronic coronary artery disease. Well-developed collaterals can reduce infarct size during acute occlusion.

\medterm{Congestive Heart Failure}-- A clinical syndrome resulting from structural or functional cardiac abnormalities that impair the ability of the ventricles to fill with or eject blood, characterized by elevated intracardiac pressures and inadequate systemic perfusion at rest or during stress. Historically termed congestive heart failure (CHF) due to prominent symptoms of venous fluid congestion; modern cardiology guidelines designate this entity simply as Heart Failure (HF), categorized by left ventricular ejection fraction: Heart Failure with Reduced Ejection Fraction (HFrEF, $\text{LVEF} \le 40\%$, systolic failure); Heart Failure with Mildly Reduced Ejection Fraction (HFmrEF, LVEF 41--49\%); and Heart Failure with Preserved Ejection Fraction (HFpEF, $\text{LVEF} \ge 50\%$, diastolic failure). Under the updated 2022 AHA/ACC/HFSA guidelines, heart failure is staged progressively: Stage A (At risk for HF, e.g., hypertension, diabetes, obesity); Stage B (Pre-HF, structural disease or elevated biomarkers without symptoms); Stage C (Symptomatic HF, past or current); and Stage D (Advanced HF, refractory symptoms requiring specialized interventions such as inotropes, LVAD, or cardiac transplantation). Functional severity is graded using the New York Heart Association (NYHA) functional classification (Classes I--IV). Core pathophysiological drivers include maladaptive neurohormonal activation (renin-angiotensin-aldosterone system [RAAS], sympathetic nervous system, and endothelin), driving adverse ventricular remodeling and myocyte apoptosis. Cardinal manifestations encompass exertional dyspnea, orthopnea, paroxysmal nocturnal dyspnea, fatigue, elevated jugular venous pressure (JVP), hepatojugular reflux, pulmonary rales, and dependent peripheral edema. Guideline-directed medical therapy (GDMT) for HFrEF mandates four foundational drug classes: (1) SGLT2 inhibitors; (2) Beta-blockers (carvedilol, metoprolol succinate, bisoprolol); (3) ARNI (sacubitril/valsartan) or ACEi/ARB; and (4) Mineralocorticoid receptor antagonists (spironolactone/eplerenone), alongside loop diuretics (furosemide) for decongestion.

\medterm{Constrictive Pericarditis}-- A severe, debilitating chronic pericardial disease characterized by fibrous thickening, loss of elasticity, contracture, and dense calcification of the fibroserous pericardium, transforming it into a rigid, non-compliant casing that encases the heart and restricts diastolic ventricular filling. Etiologies include idiopathic or post-viral pericarditis, prior mediastinal radiation therapy (e.g., for Hodgkin lymphoma or breast cancer), prior cardiac surgery ('post-pericardiotomy syndrome'), tuberculosis (the dominant cause in developing countries), connective tissue diseases, and purulent pericarditis. In early diastole, ventricular filling proceeds rapidly and unimpeded; however, as the ventricular volume abruptly reaches the mechanical limit imposed by the rigid pericardial casing in mid-to-late diastole, filling is abruptly halted. This produces pathognomonic hemodynamic features: (1) Invasive ventricular pressure tracings display the classic 'dip-and-plateau' or 'square root' sign (a prominent early diastolic dip followed by an elevated, flat plateau); (2) High interdependence between the ventricles within a fixed total volume, producing discordance on respiratory Doppler echocardiography ($>25\%$ respiratory variation in mitral inflow $E$-velocity); (3) Elevation and equalisation of end-diastolic pressures across all four chambers; and (4) Annulus paradoxus on tissue Doppler (medial mitral annular $e'$ velocity is higher than lateral $e'$, reversing normal physiology). Physical examination reveals marked systemic venous congestion (jugular venous distension, massive ascites, hepatomegaly, and peripheral edema), a prominent, sharp early diastolic sound termed a Pericardial Knock (heard 60--120 ms after S2, corresponding to the abrupt arrest of ventricular filling), and Kussmaul's Sign (a paradoxical inspiratory rise or failure to fall in jugular venous pressure, caused by the inability of the encased right heart to accommodate inspiratory venous return). Unlike cardiac tamponade, pulsus paradoxus is uncommon. Definitive curative treatment is surgical Complete Pericardiectomy ('pericardial stripping').

\medterm{Continuous Murmurs}-- Murmurs that persist throughout systole and diastole without interruption, produced by blood flowing continuously through an abnormal connection. The most common cause is patent ductus arteriosus (PDA), where blood flows from the aorta to pulmonary artery throughout the cardiac cycle. Other causes include arteriovenous fistulas and rupture of sinus of Valsalva. The murmur is characteristically machine-like in quality.

\medterm{Contrast Echocardiography}-- Use of microbubble contrast agents to enhance endocardial border definition and assess myocardial perfusion. Indicated for suboptimal non-contrast echo images and assessment of myocardial viability. Microbubbles remain entirely within the intravascular space, making them ideal for perfusion assessment. First-pass perfusion imaging detects areas of reduced myocardial blood flow. Contraindicated with intracardiac shunts (risk of paradoxical embolism).

\medterm{Cor Triatriatum}-- A rare congenital heart defect where the left atrium is divided into two chambers by a fibromuscular membrane, separating the pulmonary venous chamber (posterior) from the true left atrium (anterior) containing the left atrial appendage and mitral valve. The membrane obstructs pulmonary venous return, mimicking mitral stenosis. Presents in infancy with dyspnea, failure to thrive, recurrent respiratory infections, and pulmonary edema (if severe obstruction). Milder cases may present in adulthood with exertional dyspnea and pulmonary hypertension. Diagnosis: echocardiography shows the characteristic membrane in the left atrium with color Doppler demonstrating flow obstruction. Treatment: surgical excision of the membrane. Prognosis is excellent after repair.

\medterm{Coronary Angiography}-- The invasive imaging of coronary arteries via contrast injection and X-ray fluoroscopy, performed by catheterizing the femoral, radial, or brachial artery and advancing a catheter to the coronary ostia. Gold standard for assessing coronary artery disease, stenosis severity, and anatomy before revascularization (PCI or CABG). May include ventriculography and pressure measurements (FFR). Complications: contrast nephropathy, bleeding, dissection, arrhythmia, MI, stroke, and access-site injury. Invasive; CTA is a noninvasive alternative.

\medterm{Coronary Artery Aneurysm}-- A localized dilatation of a coronary artery exceeding the normal vessel diameter by approximately 1.5 times. It may be caused by atherosclerosis, Kawasaki disease, connective-tissue disorders, infection, trauma, or an iatrogenic procedure, and can lead to thrombosis, distal embolization, myocardial ischemia, or rupture.

\medterm{Coronary Artery Bypass Grafting (CABG)}-- A surgical revascularization procedure for severe coronary artery disease (CAD). Grafts, most commonly the left internal mammary artery (LIMA) and saphenous vein grafts (SVG), are used to bypass stenotic coronary arterial segments. Indicated for multi-vessel CAD, left main coronary artery stenosis, or severe CAD in diabetic patients. The procedure is typically performed under cardiopulmonary bypass, though off-pump CABG is an alternative. CABG improves survival, relieves angina, and reduces the risk of myocardial infarction in selected patient populations.

\medterm{Coronary Artery Disease}-- Narrowing or occlusion of coronary arteries, usually due to atherosclerosis, reducing blood flow to the myocardium. Risk factors include hypertension, diabetes, hyperlipidemia, smoking, family history, and obesity. CAD can present as stable angina, acute coronary syndromes (unstable angina, NSTEMI, STEMI), or sudden cardiac death. Diagnosis includes ECG, stress testing, and coronary angiography. Treatment involves lifestyle modification, medications, and revascularization.

\medterm{Coronary Artery Spasm}-- Transient intense constriction of an epicardial coronary artery that reduces myocardial blood flow. It may cause rest angina and transient ST-segment elevation and can be triggered by smoking, stimulants, or endothelial dysfunction; calcium-channel blockers and nitrates are commonly used for prevention and relief.

\medterm{Coronary Calcium Score}-- A CT-based measurement of coronary artery calcification, quantified as the Agatston score. A score of 0 indicates very low risk of cardiovascular events; >400 indicates high burden and significant atherosclerosis. It is used for cardiovascular risk stratification in asymptomatic individuals and may guide statin therapy decisions. It does not visualize coronary stenosis directly but reflects total atherosclerotic burden.

\medterm{Coronary Circulation}-- The vascular system supplying blood to the myocardium. The left coronary artery (LCA) divides into the left anterior descending (LAD) and circumflex (LCx) arteries. The right coronary artery (RCA) supplies the right ventricle, SA node (in 60\% of people), and inferior left ventricular wall. Coronary arteries branch into arterioles, capillaries, and drain into coronary veins and the coronary sinus.

\medterm{Coronary Dominance}-- Determined by which artery gives rise to the posterior descending artery (PDA). Right coronary dominance (85\% of people): RCA gives rise to PDA. Left coronary dominance (8\%): LCx gives rise to PDA. Codominance (7\%): both RCA and LCx contribute. Dominance affects which territories are supplied by each artery and influences the pattern of myocardial infarction.

\medterm{Coronary Heart Disease}-- Also known as coronary artery disease (CAD) or ischemic heart disease (IHD), the leading cause of morbidity and mortality worldwide, characterized by the pathological narrowing or obstruction of the coronary arteries supplying the myocardium, predominantly driven by atherosclerosis. The pathophysiological hallmark is the chronic, progressive accumulation of fibrofatty atherosclerotic plaques within the intima of medium-to-large epicardial coronary arteries. Clinical manifestations encompass chronic coronary syndromes (stable exertional angina pectoris, microvascular angina, ischemic cardiomyopathy, and asymptomatic silent ischemia) and acute coronary syndromes (unstable angina, NSTEMI, STEMI, and sudden cardiac death). Major modifiable cardiovascular risk factors include cigarette smoking, systemic arterial hypertension, dyslipidemia (elevated LDL-C, low HDL-C, elevated triglycerides), diabetes mellitus, visceral obesity, physical inactivity, and atherogenic diet; non-modifiable risk factors include advancing age, male sex (or post-menopausal status in women), and a premature family history of CAD. Diagnostic modalities include non-invasive functional testing (exercise treadmill ECG, stress echocardiography, myocardial perfusion PET/SPECT, stress cardiac MRI), anatomical imaging (coronary CT angiography [CCTA] with calcium scoring), and invasive coronary catheterization with fractional flow reserve (FFR). Management relies on intensive risk-factor modification, guideline-directed medical therapy (statins, antiplatelet agents, beta-blockers, ACE inhibitors/ARBs, SGLT2 inhibitors), and myocardial revascularization via percutaneous coronary intervention (PCI) or coronary artery bypass grafting (CABG).

\medterm{Coronary Microvascular Dysfunction}-- Impairment of the small coronary arteries and arterioles causing impaired coronary flow reserve despite normal epicardial coronary arteries. Common in women, patients with diabetes, and hypertension. Causes chest pain (angina equivalent) and may contribute to heart failure with preserved ejection fraction. Diagnosis: coronary flow reserve measurement during invasive angiography or by PET/MRI. Treatment: statins, ACE inhibitors, ranolazine, and risk factor modification.

\medterm{Corticosteroids}-- Synthetic glucocorticoids with potent anti-inflammatory and immunosuppressive effects. Examples include prednisone, prednisolone, methylprednisolone, dexamethasone, and hydrocortisone. They suppress inflammatory gene expression and inhibit immune cell function. Used in autoimmune diseases, organ transplant rejection, severe allergies, and adrenal insufficiency. Side effects include osteoporosis, hyperglycemia, weight gain, avascular necrosis, cataracts, adrenal suppression, and increased infection risk.

\medterm{CYP450 System}-- A superfamily of heme-containing enzymes primarily in the liver that metabolize drugs and xenobiotics. Major isoforms include CYP3A4 (metabolizes ~50\% of drugs), CYP2D6, CYP2C9, CYP2C19, CYP1A2, and CYP2E1. CYP inducers (rifampin, carbamazepine, phenytoin) decrease drug levels; CYP inhibitors (ketoconazole, erythromycin, grapefruit juice) increase drug levels. Genetic polymorphisms in CYP enzymes affect drug metabolism, forming the basis of pharmacogenomics.

\medterm{Deep Vein Thrombosis}-- The formation of an occlusive or non-occlusive blood clot (thrombus) within the lumen of the deep venous system, occurring predominantly in the lower extremities (classified as proximal DVT: popliteal, femoral, or iliac veins; or distal / calf DVT: anterior tibial, posterior tibial, and peroneal veins) and less frequently in the upper extremities (subclavian, axillary veins, associated with central venous catheters). Pathogenesis is governed by Virchow's Triad: (1) Endothelial Injury (catheter trauma, surgery, fractures); (2) Venous Stasis (prolonged bed rest, long-distance air travel, paralysis, heart failure); and (3) Hypercoagulability (inherited thrombophilias such as Factor V Leiden, prothrombin G20210A mutation, antithrombin/protein C/S deficiencies; or acquired states such as active malignancy, pregnancy, oral contraceptives/HRT, and antiphospholipid syndrome). Clinically, patients present with unilateral calf or thigh swelling, localized erythema, warmth, pain, and a tender, palpable venous cord. Diagnostic evaluation begins with pretest probability assessment using the Wells Score for DVT: patients with low probability undergo high-sensitivity D-dimer testing (a normal level $<500$ ng/mL reliably excludes DVT with $>98\%$ negative predictive value); patients with high pretest probability or positive D-dimer undergo Compression Duplex Ultrasonography, with non-compressibility of the venous lumen representing the gold-standard diagnostic criteria. Untreated proximal DVT carries an approximate 50\% risk of embolization causing life-threatening Pulmonary Embolism (PE), and long-term risk of Post-Thrombotic Syndrome (PTS / chronic venous hypertension with chronic edema, stasis dermatitis, and intractable ulceration). First-line treatment consists of immediate therapeutic anticoagulation for at least 3 months, with Direct Oral Anticoagulants (DOACs: apixaban, rivaroxaban, dabigatran, edoxaban) preferred over warfarin, and low-molecular-weight heparin (LMWH) reserved for pregnancy.

\medterm{Defibrillation}-- The delivery of a therapeutic, unsynchronized high-energy electrical shock across the myocardium to terminate chaotic, life-threatening ventricular tachyarrhythmias, specifically Ventricular Fibrillation (VF) and Pulseless Ventricular Tachycardia (pVT). Defibrillation works by passing an electrical current through the thorax to simultaneously depolarize a critical mass of myocardium (approximately 75--90\% of cardiac myocytes), rendering them refractory. This extinguishes all competing, chaotic reentrant wavelets and heterotopic electrical foci, permitting the heart's natural intrinsic pacemaker (the sinoatrial node) to resume organized sinus rhythm and coordinated mechanical contraction. Defibrillation is strictly unsynchronized, delivering energy immediately upon button depression regardless of ECG timing (distinguishing it from synchronized electrical cardioversion). Modern biphasic defibrillators deliver current in two phases (reversing polarity halfway through the waveform), achieving superior termination rates at lower energy levels (typically 120--200 Joules biphasic vs. 360 Joules for obsolete monophasic devices) with reduced post-shock myocardial stunning. Defibrillation success declines by 7 to 10\% for every minute of delay without CPR; immediate deployment of automated external defibrillators (AEDs) by bystanders is the single most powerful predictor of survival in out-of-hospital cardiac arrest.

\medterm{Dextrocardia}-- A congenital cardiac positional anomaly where the heart is located in the right hemithorax with the apex pointing to the right. Most cases are due to situs inversus totalis (complete mirror-image reversal of all thoracic and abdominal organs, affecting 1 in 10,000 births), which carries a normal life expectancy if cardiac anatomy is otherwise normal. Dextrocardia with situs solitus (heart on right but abdominal organs normally positioned) is strongly associated with congenital heart disease, particularly corrected transposition of the great arteries.

\medterm{Diastole}-- The phase of the cardiac cycle when the heart muscle relaxes, allowing the chambers to fill with blood. Ventricular diastole begins with isovolumetric relaxation (all valves closed, pressure falls), followed by rapid ventricular filling, slow filling (diastasis), and atrial contraction (atrial kick, contributing ~20\% of ventricular filling in normal hearts). During diastole, coronary blood flow to the left ventricular myocardium occurs (systolic compression of coronary arteries prevents flow during contraction). Diastolic dysfunction (impaired relaxation or reduced compliance) is a hallmark of heart failure with preserved ejection fraction (HFpEF) and restrictive cardiomyopathy.

\medterm{Diastolic Dysfunction}-- Impaired relaxation, compliance, or filling of the left or right ventricle during diastole. It can raise filling pressures and cause exertional dyspnoea, pulmonary congestion, or heart failure with preserved ejection fraction; echocardiography assesses relaxation, chamber structure, and estimated filling pressure.

\medterm{Diastolic Murmurs}-- Murmurs occurring between S2 and the next S1. Early diastolic decrescendo murmurs indicate semilunar valve regurgitation (aortic regurgitation, pulmonary regurgitation). Mid-diastolic low-pitched rumbles indicate AV valve stenosis (mitral stenosis, tricuspid stenosis). Late diastolic murmurs may indicate mitral stenosis with sinus rhythm. Diastolic murmurs are almost always pathological and require echocardiographic evaluation.

\medterm{Digital ischemia}-- A clinical condition characterized by compromised arterial perfusion to the fingers or toes, occurring when arterial inflow fails to meet metabolic demands, presenting with digital pain, coldness, pallor, cyanosis, and, in severe or prolonged cases, ischemic digital ulceration and dry gangrene. Etiologies are categorized into: (1) Vasospastic Disorders: severe primary Raynaud phenomenon or secondary Raynaud syndrome (associated with systemic sclerosis / scleroderma, mixed connective tissue disease, SLE, where structural microvascular intimal fibroelastosis and platelet activation compound vasospasm); (2) Thromboembolic Occlusion: microemboli from cardiac sources, aortic atheroma, or upper extremity arterial aneurysms (subclavian aneurysm in thoracic outlet syndrome), and cholesterol crystal embolization; (3) Inflammatory and Obliterative Vasculopathies: Thromboangiitis Obliterans (Buerger disease, an inflammatory, non-atherosclerotic segmental vaso-occlusive disease strongly linked to tobacco use in young smokers, causing digital ischemia and gangrene); and (4) Iatrogenic or Traumatic: radial artery catheterization thrombosis, hypothenar hammer syndrome (repetitive blunt trauma to the palmar ulnar artery against the hook of the hamate, forming an aneurysm or thrombosis), and frostbite. Physical evaluation requires assessment of peripheral pulses and the Allen Test prior to radial artery cannulation. Non-invasive diagnosis includes digital photoplethysmography (PPG), Doppler ultrasound, and digital subtraction arteriography. Management includes cold avoidance, smoking cessation, oral calcium channel blockers (amlodipine, nifedipine), phosphodiesterase-5 inhibitors (sildenafil), topical or intravenous prostaglandins (iloprost), digital sympathectomy, and catheter-directed interventions.

\medterm{Digitalis toxicity}-- A potentially life-threatening adverse effect of cardiac glycoside therapy (digoxin, digitoxin) characterized by a narrow therapeutic index (0.5-2.0 ng/mL for digoxin). Toxicity occurs through Na+/K+-ATPase inhibition, increasing intracellular sodium, which slows calcium efflux via the Na+/Ca2+ exchanger, causing increased intracellular calcium and enhanced cardiac contractility, while also increasing vagal tone and slowing AV conduction. Risk factors include advanced age, renal impairment (digoxin is renally cleared),.

\medterm{Dilated Cardiomyopathy}-- A primary myocardial disorder characterized by left ventricular (or biventricular) dilation and systolic contractile dysfunction (left ventricular ejection fraction $\text{LVEF} < 40\%$) in the complete absence of abnormal loading conditions (such as severe systemic hypertension or significant valvular stenosis) or severe coronary artery disease sufficient to cause global systolic impairment. DCM is a leading indication for cardiac transplantation. Etiologies are categorized into genetic / familial (~30--50\% of cases, most frequently autosomal dominant truncating variants in the Titin gene [TTN], Lamin A/C [LMNA, associated with high risk of sudden death and AV block], and sarcomeric/desmosomal genes) and non-genetic causes (viral myocarditis [coxsackievirus, parvovirus B19, adenovirus], chronic alcohol toxicity [alcoholic cardiomyopathy], anthracycline chemotherapy [doxorubicin], peripartum cardiomyopathy, autoimmune disease, and thyroid dysfunction). The pathological hallmark is marked ventricular chamber enlargement with eccentric hypertrophy and patchy interstitial fibrosis. Patients present with progressive congestive heart failure (exertional dyspnea, orthopnea, peripheral edema), functional mitral/tricuspid regurgitation due to annular dilation, ventricular arrhythmias, and mural thrombi. Prognosis and survival have dramatically improved with Guideline-Directed Medical Therapy (GDMT 'four pillars': ARNI/ACEi, beta-blockers, MRAs, SGLT2 inhibitors). Primary prevention Implantable Cardioverter-Defibrillators (ICDs) are indicated for persistent $\text{LVEF} \le 35\%$ despite $\ge 3$ months of optimal medical therapy, and Cardiac Resynchronization Therapy (CRT) for patients with LBBB and QRS duration $\ge 150$ ms.

\medterm{Direct Oral Anticoagulants}-- DOACs directly inhibit specific coagulation factors. Direct thrombin inhibitors: dabigatran. Factor Xa inhibitors: rivaroxaban, apixaban, edoxaban, betrixaban. They have predictable pharmacokinetics, fixed dosing, and routine monitoring is not required. They have fewer drug interactions than warfarin. Dabigatran and rivaroxaban are reversed by specific agents (idarucizumab, andexanet alfa). DOACs are preferred over warfarin for non-valvular atrial fibrillation and VTE.

\medterm{Dobutamine Stress Echocardiography}-- Pharmacological stress test using dobutamine (beta-1 agonist) to increase heart rate, contractility, and myocardial oxygen demand. New or worsening wall motion abnormalities indicate ischemia. Preferred when patients cannot exercise. Beta-blockers should be held before testing. Sensitivity 80-85\% and specificity 80-88\% for detecting significant coronary artery disease. Additional atropine may be administered if target heart rate not achieved.

\medterm{Double Outlet Right Ventricle}-- A congenital heart defect where more than 50\% of both great arteries arise from the right ventricle, with a VSD as the only outlet from the left ventricle. Classified by the position of the VSD relative to the great arteries: subaortic VSD (TOF-type physiology), subpulmonic VSD (Taussig-Bing malformation, TGA-type physiology), doubly committed VSD, and non-committed VSD (remote from both arteries). Presents with variable physiology: cyanosis or heart failure depending on the degree of pulmonary blood flow. The Taussig-Bing malformation presents with cyanosis and early heart failure. Diagnosis: echocardiography showing both great arteries arising from the RV. Treatment: surgical repair tailored to VSD position — arterial switch operation for Taussig-Bing, intracardiac tunnel repair for subaortic VSD.

\medterm{Dressler Syndrome}-- A delayed pericarditis occurring 2-10 weeks after myocardial infarction or cardiac surgery. It is an autoimmune-mediated inflammation caused by antibodies against myocardial antigens released during tissue injury. Symptoms include fever, pleuritic chest pain, and pericardial friction rub. ECG may show diffuse ST elevation. Treatment includes NSAIDs, colchicine, and avoidance of anticoagulants if possible. It is a type of post-cardiac injury syndrome.

\medterm{Drug Bioavailability}-- The fraction of an administered drug that reaches systemic circulation unchanged. Oral bioavailability is affected by first-pass metabolism, absorption, and drug formulation. IV administration has 100\% bioavailability. Bioavailability influences drug dosing and route selection. Drugs with low oral bioavailability (e.g., morphine ~25\%) require higher oral doses compared to IV doses. Food can increase or decrease bioavailability depending on the drug.

\medterm{Drug Clearance}-- The volume of plasma completely cleared of drug per unit time. Hepatic clearance depends on hepatic blood flow and enzyme activity; renal clearance depends on glomerular filtration, tubular secretion, and reabsorption. Total clearance determines maintenance dose. Clearance decreases in hepatic or renal impairment, requiring dose adjustment. Drug clearance is independent of concentration for first-order kinetics but constant for zero-order (saturable) kinetics.

\medterm{Drug Hypersensitivity Reactions}-- Immune-mediated adverse drug reactions classified by Gell and Coombs. Type I (immediate, IgE-mediated: anaphylaxis, urticaria); Type II (cytotoxic, IgG/IgM-mediated: hemolytic anemia, thrombocytopenia); Type III (immune complex-mediated: serum sickness, vasculitis); Type IV (delayed, T-cell mediated: contact dermatitis, Stevens-Johnson syndrome). Management includes drug discontinuation, supportive care, and in severe cases, immunosuppression. Cross-reactivity between related drugs must be considered.

\medterm{Dual Antiplatelet Therapy}-- The combination of aspirin and a P2Y12 inhibitor (clopidogrel, prasugrel, or ticagrelor) for prevention of thrombotic events after acute coronary syndrome and percutaneous coronary intervention. Duration varies: 1 month after bare-metal stent, 6-12 months after drug-eluting stent, and 12 months after ACS. Prasugrel is more potent but increases bleeding risk; ticagrelor is reversible but requires twice-daily dosing. DAPT score guides duration of therapy.

\medterm{Ductus Arteriosus}-- A fetal blood vessel connecting the pulmonary artery to the descending aorta, allowing blood to bypass the non-functioning fetal lungs. In utero, the ductus arteriosus is kept patent by prostaglandin E2 and low oxygen tension. After birth, the ductus normally closes within 24--48 hours (functional closure) and becomes the ligamentum arteriosum by 2--3 weeks (anatomic closure). Failure of closure results in patent ductus arteriosus (PDA), a left-to-right shunt causing continuous machinery murmur, bounding pulses, and volume overload of the left heart. Prostaglandin E1 (alprostadil) is used to maintain ductal patency in ductal-dependent congenital heart disease (e.g., pulmonary atresia, coarctation of the aorta).

\medterm{Dyslipidemia}-- An abnormal concentration or composition of circulating lipoproteins, including elevated low-density lipoprotein cholesterol, elevated triglycerides, reduced high-density lipoprotein cholesterol, or mixed abnormalities. It contributes to atherosclerotic cardiovascular disease and may be primary, inherited, or secondary to another disease or exposure.

\medterm{Dyspnea on Exertion}-- Shortness of breath precipitated by physical activity, a cardinal symptom of cardiac and pulmonary disease. Cardiac causes: heart failure, coronary artery disease (angina equivalent), valvular disease; pulmonary: COPD, asthma, ILD, anemia, deconditioning. Quantified by NYHA functional class (heart failure) and mMRC scale (dyspnea). Evaluation: history (onset, associated symptoms), physical exam, ECG, chest X-ray, echo, spirometry, and exercise testing. The symptom reflects the demand/supply mismatch of oxygen delivery.

\medterm{Ebstein Anomaly}-- A rare congenital heart defect accounting for less than 1 percent of all congenital heart disease, characterized by apical displacement of the septal and posterior leaflets of the tricuspid valve into the right ventricle, resulting in atrialization of the proximal right ventricle and a small functional right ventricular chamber. The anterior leaflet is usually large and sail-like. The tricuspid valve is often regurgitant, causing right atrial enlargement and right-to-left shunting across an associated atrial septal defect or patent foramen ovale.

\medterm{Echocardiogram}-- The primary, non-invasive ultrasonic imaging modality in clinical cardiology, providing real-time anatomical, functional, and hemodynamic evaluation of the cardiac chambers, valves, great vessels, and pericardium. Standard clinical modalities include: (1) Transthoracic Echocardiography (TTE): performed via the anterior thoracic cage across four standard acoustic windows (parasternal, apical, subcostal, and suprasternal); (2) Transesophageal Echocardiography (TEE): utilizes an endoscopic ultrasonic transducer positioned within the esophagus and stomach directly behind the left atrium, avoiding ribs and lungs to provide exquisite, high-resolution visualization of the mitral valve, aortic root, left atrial appendage (excluding thrombus prior to cardioversion or ablation), and small infective endocarditis vegetations; (3) M-Mode and 2D/3D Imaging: measures chamber dimensions, myocardial wall thickness (diagnosing concentric vs. eccentric LVH), and computes left ventricular ejection fraction via the biplane Simpson method of discs; (4) Spectral and Color Flow Doppler: utilizes the Doppler shift principle to assess direction and velocity of blood flow, enabling the non-invasive calculation of pressure gradients across stenotic valves via the simplified Bernoulli equation ($\Delta P = 4v^2$), quantification of valve areas via the continuity equation, and assessment of valvular regurgitation; and (5) Tissue Doppler and Strain Imaging: evaluates myocardial diastolic relaxation velocities ($e'$) and global longitudinal strain (GLS, sensitive for early subclinical systolic dysfunction in amyloidosis and cardiotoxicity).

\medterm{Echocardiography}-- Ultrasound imaging of the heart, the primary modality for assessing cardiac structure and function. Transthoracic echocardiography (TTE) is the most common approach; transesophageal echocardiography (TEE) provides better views of posterior structures (left atrium, mitral valve, aorta). Doppler echocardiography assesses blood flow velocities, valve function, and pressure gradients. Echocardiography is used to evaluate ejection fraction, valve disease, pericardial effusion, and structural abnormalities.

\medterm{Ectopic Heartbeat}-- A premature, extra cardiac depolarization originating from an abnormal, irritable pacemaker focus situated outside the dominant sinoatrial (SA) node within the cardiac electrical conduction system or myocardium, perceived by patients as a 'skipped beat', flutter, or post-pause 'thump' in the chest. Ectopic beats arise electrophysiologically through enhanced abnormal automaticity, triggered activity (early or delayed afterdepolarizations), or micro-reentry. Classified according to anatomical site of origin: (1) Premature Atrial Contractions (PACs): arise from an ectopic atrial focus, producing an early, morphologically abnormal P-wave ($P'$) that typically conducts through the AV node with a normal, narrow QRS complex, followed by an incomplete compensatory pause as the sinus node is reset; (2) Premature Junctional Contractions (PJCs): arise from the AV junction, producing a narrow QRS without preceding P-waves (or with retrograde inverted P-waves); and (3) Premature Ventricular Contractions (PVCs): arise from ventricular myocardium below the bifurcation of the bundle of His, producing an early, wide, bizarre QRS complex ($\ge 120$ ms) with discordant ST-T segment changes, followed by a full compensatory pause (the sinus cycle is undisturbed). Ectopic beats are ubiquitous and frequently benign in healthy individuals, exacerbated by caffeine, alcohol, nicotine, sympathomimetic decongestants, emotional stress, lack of sleep, or electrolyte imbalances (hypokalemia, hypomagnesemia). However, in the presence of structural heart disease (prior myocardial infarction, cardiomyopathy, channelopathies), frequent or complex PVCs (couplets, triplets, R-on-T phenomenon) signify heightened electrical instability and predispose to sustained ventricular tachycardia or ventricular fibrillation; high-burden PVCs ($>10--15\%$ of total daily beats) can directly cause a reversible PVC-induced cardiomyopathy.

\medterm{Edema (historically Dropsy)}-- An historic term for edema, referring to the abnormal accumulation of serous fluid in the tissues or body cavities. Types: peripheral edema (dependent: feet, ankles, legs, sacrum), pulmonary edema (fluid in lung interstitium and alveoli--causes include left heart failure, ARDS, high altitude), ascites (fluid in peritoneal cavity--cirrhosis, heart failure, malignancy), anasarca (generalized severe edema). Pathophysiology: increased hydrostatic pressure (heart failure, venous obstruction), decreased plasma oncotic pressure (nephrotic syndrome, cirrhosis, malnutrition), increased capillary permeability (inflammation, burns, sepsis), and lymphatic obstruction (lymphedema). Treatment: address underlying cause, sodium restriction, diuretics, and mechanical compression.

\medterm{Eisenmenger Syndrome}-- Reversal of a left-to-right cardiac shunt due to irreversible pulmonary hypertension, resulting in right-to-left shunting and cyanosis. It occurs in untreated congenital heart defects with large left-to-right shunts (VSD, ASD, PDA) that cause progressive pulmonary vascular disease. Features include cyanosis, clubbing, polycythemia, and paradoxical embolism. Management is supportive; pregnancy is contraindicated. Lung transplantation may be considered in severe cases.

\medterm{Ejection Click}-- A high-pitched sound occurring at the beginning of systole, caused by opening of a stenotic semilunar valve. Aortic ejection click is heard in bicuspid aortic valve and aortic stenosis with pliable leaflets. Pulmonary ejection click is heard in pulmonary stenosis and disappears as the valve becomes more calcified and immobile. An ejection click helps distinguish valvular stenosis from subvalvular or supravascular obstruction.

\medterm{Ejection Fraction}-- The percentage of blood ejected from the left ventricle during systole, calculated as (stroke volume / end-diastolic volume) x 100. Normal EF is 55-70\%. Heart failure with reduced ejection fraction (HFrEF): EF 40\% or less; heart failure with preserved ejection fraction (HFpEF): EF 50\% or greater. EF is measured by echocardiography, cardiac MRI, or ventriculography.

\medterm{Electrical Cardioversion}-- The therapeutic delivery of a low-to-moderate-energy electrical shock synchronized to the intrinsic electrical activity of the heart to terminate organized supraventricular or ventricular tachyarrhythmias (such as atrial fibrillation, atrial flutter, atrioventricular nodal reentrant tachycardia [AVNRT], and hemodynamically unstable monomorphic ventricular tachycardia with a pulse). Unlike unsynchronized defibrillation, the cardioverter machine must be set to 'Sync' mode, which detects and tracks the patient's QRS complex, automatically delivering the electrical discharge precisely on the peak of the R wave (or downslope of the S wave). This precise synchronization ensures that electrical energy is delivered strictly outside the cardiac vulnerable period (the middle-to-late T-wave, representing inhomogeneous ventricular repolarization); delivering an unsynchronized shock on the T-wave ('R-on-T phenomenon') can precipitate fatal ventricular fibrillation. Biphasic energy doses vary: 50--100 J for atrial flutter and SVT; 120--200 J for atrial fibrillation (increasing stepwise to 200--360 J if unsuccessful); and 100 J for monomorphic VT with a pulse. In non-emergent elective cardioversion for atrial fibrillation or flutter lasting $>48$ hours, therapeutic anticoagulation is mandatory for at least 3 weeks prior and 4 weeks post-procedure (or exclusion of left atrial thrombus via transesophageal echocardiography [TEE]) to prevent post-cardioversion thromboembolic stroke upon restoration of coordinated atrial mechanical contraction ('atrial stunning'). Procedural sedation (propofol, etomidate, or midazolam) and airway support are routinely administered.

\medterm{Electrocardiogram (ECG / EKG)}-- A recording of the heart's electrical activity using electrodes placed on the body surface. A standard 12-lead ECG provides views of the heart from different angles: limb leads (I, II, III, aVR, aVL, aVF) and precordial leads (V1-V6). ECG is used to diagnose arrhythmias, conduction disturbances, myocardial ischemia/infarction, chamber hypertrophy, electrolyte abnormalities, and drug effects.

\medterm{Endocardial Cushion Defect}-- A spectrum of congenital heart malformations (also designated Atrioventricular Septal Defect [AVSD] or atrioventricular canal defect) arising from the embryological failure of the superior and inferior endocardial cushions to fuse properly during heart development (weeks 4 to 5 of gestation). The endocardial cushions are responsible for dividing the primitive common atrioventricular canal into separate tricuspid and mitral orifices and forming the lower portion of the atrial septum (ostium primum) and upper inlet portion of the ventricular septum. AVSD is strongly linked to chromosomal aneuploidy, occurring in approximately 40--45\% of children with Down Syndrome (Trisomy 21). Anatomically classified into: (1) Complete AVSD: characterized by a large ostium primum atrial septal defect, a large inlet ventricular septal defect, and a solitary, common, five-leaflet atrioventricular valve bridging both ventricles, producing massive left-to-right shunting, marked pulmonary overcirculation, congestive heart failure in early infancy, and rapid development of irreversible pulmonary vascular obstructive disease (Eisenmenger syndrome) if uncorrected; (2) Partial AVSD: consists of an ostium primum ASD and two separate AV valve orifices with a characteristic cleft in the anterior mitral valve leaflet, causing mitral regurgitation without a ventricular defect; and (3) Transitional AVSD: intermediate variants. A pathognomonic 12-lead ECG finding across all forms is Extreme Left Axis Deviation (superior axis, $-30^\circ$ to $-90^\circ$ or 'Northwest' axis) with a prolonged PR interval (first-degree AV block) due to postero-inferior displacement of the AV node and bundle of His. Complete AVSD requires primary surgical repair (closing septal defects and dividing the common valve into competent tricuspid and mitral valves) within the first 3 to 6 months of life.

\medterm{Endocarditis}-- Infection of the endocardial surface of the heart, most commonly involving the valves. Acute bacterial endocarditis: Staph aureus (most common in IV drug users and prosthetic valves), rapid destruction. Subacute: viridans group streptococci (pre-existing valve disease). Symptoms: fever, heart murmur, petechiae, splinter hemorrhages, Janeway lesions, Osler nodes, Roth spots. Diagnosis: modified Duke criteria, blood cultures (2-3 sets), echocardiography (transthoracic and transesophageal). Complications: valvular regurgitation, heart failure, embolic phenomena (stroke, septic emboli), perivalvular abscess. Treatment: prolonged IV antibiotics based on organism and sensitivity; surgical valve replacement for severe regurgitation, persistent infection, or large vegetations.

\medterm{Endocarditis Surgical Indications}-- Surgery indicated for: heart failure from valve dysfunction, uncontrolled infection (abscess, fistula, vegetation >10 mm with persistent emboli), fungal endocarditis, prosthetic valve dehiscence, and large vegetations (>15 mm) with embolic events. Timing: urgent (within days) for heart failure or uncontrolled infection; emergent (immediate) for acute aortic regurgitation with heart failure or ruptured mycotic aneurysm.

\medterm{Endocardium}-- The innermost layer of the heart wall, a thin membrane of endothelial cells and underlying connective tissue that lines the cardiac chambers and covers the heart valves. The endocardium is continuous with the endothelium of blood vessels and provides a smooth, non-thrombogenic surface. Endocardial damage can predispose to infective endocarditis, and endocardial fibrosis may occur in certain conditions.

\medterm{Epicardium}-- The outermost layer of the heart wall, which is the visceral layer of the serous pericardium. It consists of a single layer of mesothelial cells and underlying connective tissue containing blood vessels, nerves, and adipose tissue. The epicardium produces pericardial fluid that lubricates the heart surface. It can be involved in inflammatory processes such as pericarditis and may contain epicardial fat pads.

\medterm{Essential Hypertension}-- A chronic, sustained elevation of systemic arterial blood pressure (defined by 2017 ACC/AHA guidelines as systolic $\text{BP} \ge 130$ mmHg or diastolic $\text{BP} \ge 80$ mmHg, or by ESC/ESH guidelines as $\ge 140/90$ mmHg) occurring in the absence of a distinct, identifiable secondary etiology (such as renal artery stenosis, primary aldosteronism, or pheochromocytoma), accounting for approximately 90--95\% of all diagnosed cases of clinical hypertension. Essential (primary) hypertension is a multifactorial, polygenic disorder arising from complex, lifelong gene-environment interactions that disrupt normal hemodynamic homeostasis. Key pathophysiological drivers include: (1) Altered Renal Sodium Handling: a rightward shift in the renal pressure-natriuresis curve requiring elevated arterial pressure to excrete dietary sodium loads; (2) Inappropriate Neurohormonal Activation: chronic sympathetic nervous system overactivity and hyperactivation of the Renin-Angiotensin-Aldosterone System (RAAS); (3) Vascular Endothelial Dysfunction: impaired nitric oxide synthesis and increased production of endothelin-1, provoking elevated total peripheral vascular resistance (SVR); and (4) Progressive Arterial Stiffening: loss of elastin and accumulation of collagen in large conduit arteries, driving isolated systolic hypertension and elevated pulse wave velocity in elderly patients. Uncontrolled essential hypertension is a major 'silent killer', driving asymptomatic, progressive target organ damage: left ventricular hypertrophy (LVH), HFpEF, coronary atherosclerosis, arteriolosclerotic nephrosclerosis (hypertensive chronic kidney disease), ischemic and hemorrhagic stroke, and hypertensive retinopathy. First-line guideline-directed medical therapy (GDMT) comprises thiazide-like diuretics (chlorthalidone, indapamide), calcium channel blockers (amlodipine), and ACE inhibitors or ARBs, frequently administered as initial single-pill fixed-dose combination therapy.

\medterm{Ewart Sign}-- Dullness to percussion, decreased breath sounds, and bronchial breathing at the left infrascapular region caused by compressive atelectasis from a large pericardial effusion. It is a physical examination finding of pericardial effusion. The effusion compresses the adjacent left lower lobe, creating a zone of consolidation. Ewart sign is best detected with the patient in the left lateral decubitus position.

\medterm{Exercise Intolerance}-- A cardinal clinical symptom and objective physiological limitation characterized by the reduced ability of an individual to perform physical exertion or sustain physical activity at an age- and sex-predicted expected level, manifesting as early exertional fatigue, disabling breathlessness (dyspnea), heavy legs, or chest discomfort. In cardiovascular medicine, exercise intolerance is the primary clinical hallmark of chronic Heart Failure (both HFrEF and HFpEF), chronic coronary syndromes, pulmonary hypertension, and valvular heart disease. Pathophysiologically, normal exercise demands a 4- to 6-fold increase in cardiac output (achieved through augmented heart rate and stroke volume mediated by the Frank-Starling mechanism, enhanced contractility, and peripheral vasodilation). In heart failure, cardiac output reserve is severely blunted (chronotropic incompetence, systolic dysfunction, or reduced stroke volume due to impaired diastolic relaxation). Consequently, left ventricular filling pressures rise sharply during exercise, transmitting retrogradely into the pulmonary capillary bed to cause exertional pulmonary venous congestion and dyspnea. Furthermore, inadequate peripheral skeletal muscle blood flow, endothelial dysfunction, impaired skeletal muscle mitochondrial oxidative phosphorylation, and muscle fiber atrophy contribute heavily to early exertional lactic acidosis and peripheral fatigue. Objective clinical quantification is achieved via Cardiopulmonary Exercise Testing (CPET), which measures peak oxygen consumption ($\dot{V}O_{2\text{ peak}}$): a $\dot{V}O_{2\text{ peak}} < 14$ mL/kg/min (or $<12$ mL/kg/min on beta-blockers) indicates severe functional impairment and serves as an objective prognostic threshold for cardiac transplantation evaluation.

\medterm{Exercise Stress Test}-- A noninvasive cardiac diagnostic test that evaluates the cardiovascular response to graded physical exertion, primarily used to diagnose coronary artery disease, assess functional capacity, and guide prognosis. The test is most commonly performed using a treadmill (Bruce protocol: graded increase in speed and incline every 3 minutes) or stationary bicycle ergometer, following the standard Bruce or modified Bruce protocol. During the test, continuous 12-lead ECG monitoring records heart rate, rhythm, and ST-segment changes.

\medterm{Exercise-Induced Arrhythmias}-- Catecholamine-sensitive ventricular tachycardia and right ventricular outflow tract VT can be triggered by exercise. HCM-related sudden death risk increases during vigorous exercise. Exercise-induced myocardial ischemia may trigger ventricular arrhythmias in CAD. Arrhythmogenic right ventricular cardiomyopathy frequently presents with exercise-related VT. Athletes with structural heart disease or inherited arrhythmia syndromes may require activity restriction and ICD placement.

\medterm{Extracorporeal Membrane Oxygenation}-- A temporary mechanical circulatory support system that provides cardiac and/or respiratory support. VA-ECMO (veno-arterial) provides both cardiac and respiratory support; VV-ECMO (veno-venous) provides respiratory support only. Used as bridge to recovery or decision in severe cardiogenic shock, refractory cardiac arrest, or acute respiratory failure. Complications include bleeding, thromboembolism, limb ischemia, and neurological injury.

\medterm{Familial Hypercholesterolemia}-- An autosomal dominant genetic disorder causing severely elevated LDL cholesterol from birth, most commonly due to LDL receptor mutations. Heterozygous FH affects 1 in 250 people, causing premature coronary artery disease (MI in 40s-50s). Diagnosis: LDL >190 mg/dL in adults, or >160 mg/dL in children. Treatment: high-intensity statins, ezetimibe, PCSK9 inhibitors. Cascade screening of family members is essential. Lipoprotein apheresis may be needed for severe cases.

\medterm{Fast Heart Rate}-- A clinical and physiological term denoting a cardiac rhythm faster than the expected physiological baseline for a patient's age and activity level, formally designated as Tachycardia (defined in adults as a resting ventricular heart rate $>100$ beats per minute). Tachycardia may represent an appropriate physiological compensatory response (sinus tachycardia driven by sympathetic activation and vagal withdrawal in response to exercise, emotional excitement, pain, fever, anemia, dehydration, hypovolemia, sepsis, or hypoxia) or an inappropriate pathological arrhythmia. Pathological tachycardias are classified electrophysiologically on a 12-lead ECG by QRS duration (narrow-complex $<120$ ms vs. wide-complex $\ge 120$ ms) and regularity: (1) Narrow-complex regular tachycardias: sinus tachycardia, atrial flutter with fixed AV conduction, atrioventricular nodal reentrant tachycardia (AVNRT), and atrioventricular reentrant tachycardia (AVRT); (2) Narrow-complex irregular tachycardias: atrial fibrillation (irregularly irregular), atrial flutter with variable block, and multifocal atrial tachycardia (MAT); (3) Wide-complex regular tachycardias: monomorphic ventricular tachycardia (VT, accounting for $>80\%$ of wide-complex tachycardias), SVT with preexisting or rate-dependent bundle branch block (aberrancy), or antidromic AVRT; and (4) Wide-complex irregular tachycardias: atrial fibrillation with bundle branch block, pre-excited atrial fibrillation in WPW (potentially fatal if AV-nodal blockers are given), and polymorphic VT / Torsades de Pointes. Sustained excessive heart rates compromise diastolic ventricular filling time, reducing stroke volume and coronary arterial perfusion while increasing myocardial oxygen demand, potentially triggering angina, syncope, hypotension, acute pulmonary edema, or tachycardia-induced cardiomyopathy.

\medterm{Fibroma}-- A rare, benign primary cardiac neoplasm composed of differentiated fibroblasts, myofibroblasts, and an abundant extracellular collagenous matrix, representing the second most common primary cardiac tumor in infants and children (after rhabdomyoma). Cardiac fibromas are solitary, unencapsulated, firm, grayish-white intramural masses that arise predominantly within the thick ventricular myocardium, most frequently involving the anterior left ventricular free wall or the interventricular septum. They do not undergo spontaneous regression (unlike cardiac rhabdomyomas). While histologically benign, fibromas frequently produce severe, life-threatening clinical consequences due to their strategic intramural location: (1) Intractable ventricular arrhythmias (VT, VF) and sudden cardiac death, caused by disruption of the myocardial syncytium and creation of reentrant electrical circuits; (2) Conduction system disruption: septal fibromas can invade or compress the bundle of His and atrioventricular node, causing high-grade or complete heart block; and (3) Ventricular inflow or outflow tract obstruction, mimicking aortic or subaortic stenosis and provoking congestive heart failure. Approximately 3--5\% of cardiac fibromas are associated with Gorlin Syndrome (Nevoid Basal Cell Carcinoma Syndrome), an autosomal dominant condition caused by mutations in the PTCH1 tumor suppressor gene. Transthoracic echocardiography demonstrates a non-contractile, sharply demarcated intramural mass; cardiac MRI provides definitive characterization, showing a mass that is typically isointense on T1-weighted images, hypointense on T2-weighted images (reflecting dense, fibrous collagenous content), and displays intense, homogeneous Late Gadolinium Enhancement (LGE). Surgical resection is indicated for symptomatic tumors, ventricular obstruction, or recurrent arrhythmias.

\medterm{First-Degree AV Block}-- A conduction delay in which every atrial impulse reaches the ventricles but with a prolonged PR interval (>0.20 seconds). It is caused by delayed conduction through the AV node or His-Purkinje system. First-degree block is usually benign and asymptomatic. Causes include increased vagal tone, drugs (beta-blockers, calcium channel blockers, digoxin), inferior myocardial infarction, and fibrotic degeneration. Treatment of the underlying cause is usually sufficient.

\medterm{First-Pass Metabolism}-- The metabolism of a drug by the liver before it reaches systemic circulation after oral administration. Drugs with high first-pass metabolism (e.g., morphine, nitroglycerin, lidocaine) have low oral bioavailability and may require alternative routes (sublingual, transdermal, IV). First-pass metabolism can be affected by liver disease, enzyme induction, and enzyme inhibition. Prodrugs (e.g., enalapril) require hepatic activation during first-pass metabolism.

\medterm{Fractional Flow Reserve}-- An invasive physiological measurement during coronary catheterization using a pressure wire to assess the hemodynamic significance of a coronary stenosis. FFR = Pa/Pd (aortic pressure / distal coronary pressure) during maximal hyperemia (adenosine). Normal FFR is 1.0; FFR <=0.80 indicates ischemia-producing stenosis warranting revascularization. The FAME and FAME 2 trials demonstrated that FFR-guided PCI improves outcomes compared to angiography-guided PCI.

\medterm{Frank-Starling Mechanism}-- The physiological principle stating that stroke volume increases as end-diastolic volume (preload) increases, up to an optimal point. Within physiological limits, greater ventricular stretch leads to stronger contraction due to increased calcium sensitivity of myofilaments and optimal overlap of actin and myosin. This mechanism allows the heart to automatically match output to venous return. In heart failure, the curve shifts downward and rightward.

\medterm{Gallavardin Phenomenon}-- The dissociation of the murmur of aortic stenosis into two components: a musical quality heard at the apex and a harsh, rough quality heard at the right upper sternal border. This occurs because high-frequency components are transmitted to the apex while low-frequency components are heard at the base. It helps distinguish aortic stenosis from mitral regurgitation, which also presents with an apical murmur.

\medterm{Half-Life (t1/2)}-- The time required for plasma drug concentration to decrease by 50\%. It determines dosing frequency and time to steady state (4-5 half-lives). Drugs with long half-lives (e.g., amiodarone: 40-55 days, warfarin: 40 hours) can be given once daily but take longer to reach steady state and to be eliminated. Short half-lives (e.g., nitroglycerin: 3 minutes) require continuous infusion or frequent dosing. Half-life is affected by hepatic and renal function.

\medterm{HDL}-- High-density lipoprotein; the 'good' cholesterol that mediates reverse cholesterol transport, removing cholesterol from peripheral tissues and arterial walls and returning it to the liver for excretion. High HDL is associated with reduced cardiovascular risk; low HDL (<40 mg/dL men, <50 mg/dL women) is a risk factor. Levels are increased by exercise, weight loss, and moderate alcohol; pharmacologic elevation (niacin) does not reduce events. Measured in lipid panel; clinical focus is on overall risk and LDL lowering.

\medterm{Heart}-- A four-chambered, hollow muscular pump located within the middle mediastinum, enveloped by the fibroserous pericardium. The heart consists of two functional syncytia: the right heart (low-pressure pulmonary circulation) and the left heart (high-pressure systemic circulation), separated by the interatrial and interventricular septa. Deoxygenated systemic venous blood returns via the superior and inferior venae cavae into the right atrium (normal mean pressure: 2--6 mmHg), passes through the tricuspid valve into the right ventricle (systolic/diastolic pressure: 15--30 / 2--8 mmHg), and is ejected across the pulmonary valve into the pulmonary arterial tree. Oxygenated blood from the four pulmonary veins enters the left atrium (normal mean pressure: 6--12 mmHg), crosses the bicuspid mitral valve into the left ventricle (systolic/diastolic pressure: 100--140 / 3--12 mmHg), and is pumped across the aortic valve into the systemic arterial tree. Cardiac output (CO) is the product of heart rate (HR) and stroke volume (SV), averaging 4.5 to 5.5 L/min at rest (cardiac index: 2.5--4.0 L/min/m$^2$). The cardiac wall comprises three layers: inner endocardium, middle contractile myocardium, and outer epicardium (visceral pericardium). Intrinsic electrical impulse generation at the sinoatrial (SA) node coordinates rhythmic atrioventricular contraction.

\medterm{Heart Attack (Myocardial Infarction)}-- Acute myocardial ischemic necrosis caused by an abrupt, profound mismatch between myocardial oxygen supply and demand, most commonly precipitated by the acute rupture or erosion of an unstable atherosclerotic plaque in an epicardial coronary artery with superimposed occlusive thrombus formation. Classified according to the Fourth Universal Definition into five clinical types: Type 1 (spontaneous MI due to plaque rupture, ulceration, or erosion with acute intraluminal thrombosis); Type 2 (myocardial necrosis secondary to ischemic supply-demand mismatch in the absence of plaque rupture, such as coronary spasm, severe hypotension, anemia, tachyarrhythmias, or respiratory failure); Type 3 (cardiac death with symptoms of ischemia and presumed new ischemic ECG changes prior to blood sampling); Type 4 (PCI-related, Type 4a, or stent thrombosis, Type 4b); and Type 5 (CABG-related). Diagnosis mandates a rise and/or fall of cardiac troponin (cTnI or cTnT, with at least one value above the 99th percentile URL) accompanied by at least one of: ischemic chest symptoms (retrosternal pressure radiating to the left arm, neck, or jaw, diaphoresis, dyspnea), new ischemic ECG changes (ST-elevation, horizontal/downsloping ST-depression, pathological Q waves), imaging evidence of new regional wall motion abnormalities, or angiographic identification of an intracoronary thrombus. Emergency treatment requires immediate chewed aspirin, dual antiplatelet therapy, anticoagulation, and emergent reperfusion: primary PCI (door-to-balloon $<90$ min) is gold standard for STEMI, while NSTEMI undergoes risk stratification (TIMI or GRACE scores) for early invasive coronary angiography within 24 hours. Long-term secondary prevention comprises lifelong aspirin, $P2Y_{12}$ inhibitor, high-intensity statin, beta-blocker, ACE inhibitor/ARB, and cardiac rehabilitation.

\medterm{Heart Block}-- A delay or interruption in the conduction of cardiac impulses from the atria to the ventricles through the atrioventricular (AV) node and His-Purkinje system. Classification: first-degree AV block: PR interval >200 ms, each atrial impulse conducts to the ventricles, typically asymptomatic. Second-degree AV block: Mobitz type I (Wenckebach): progressive PR prolongation then a dropped QRS, usually benign, localised to AV node. Mobitz type II: constant PR interval with intermittent dropped QRS, indicates disease below the AV node (His-Purkinje), risk of progression to complete heart block. Third-degree (complete) AV block: no atrial impulses conducted to ventricles; atria and ventricles beat independently (AV dissociation). Ventricular escape rhythm (wide QRS, rate 30-50 bpm) or junctional escape (narrow QRS, rate 40-60 bpm). Causes: idiopathic fibrosis (Lenegre-Lev disease, most common), ischemic heart disease, myocardial infarction (inferior MI--AV nodal block, usually transient; anterior MI--infranodal block, often permanent, poor prognosis), cardiomyopathy, drugs (beta-blockers, calcium channel blockers, digoxin), electrolyte abnormalities (hyperkalaemia), Lyme disease, Chagas disease, endocarditis (aortic valve abscess), post-cardiac surgery, and infiltrative diseases (sarcoidosis, amyloidosis). Symptoms: syncope, presyncope, dizziness, fatigue, dyspnea, angina (see Stokes-Adams attacks). Treatment: urgent pacing for symptomatic or high-grade AV block (type II second-degree or third-degree). Temporary pacing: transcutaneous (external) or transvenous. Permanent pacemaker: dual-chamber (DDD) preferred for sinus node function, single-chamber (VVI) for atrial fibrillation with AV block. Dual-chamber pacing maintains AV synchrony and reduces atrial fibrillation and heart failure risk.

\medterm{Heart Disease}-- A broad term encompassing various conditions affecting the heart and blood vessels, most commonly coronary artery disease, heart failure, arrhythmias, and valvular heart disease. Atherosclerosis, the buildup of plaque in the arteries, is the underlying cause of most cardiovascular events. Risk factors include hypertension, hyperlipidemia, smoking, diabetes, obesity, and physical inactivity.

\medterm{Heart Failure with Preserved Ejection Fraction (HFpEF)}-- A clinical heart failure syndrome characterized by signs and symptoms of congestive heart failure (exertional dyspnea, fatigue, elevated JVP, pulmonary rales, peripheral edema) in the setting of a normal or near-normal left ventricular ejection fraction ($\text{LVEF} \ge 50\%$), accounting for approximately 50\% of all heart failure cases worldwide, predominantly affecting elderly, female, and hypertensive patients with metabolic comorbidities (obesity, type 2 diabetes mellitus, atrial fibrillation, chronic kidney disease). The primary underlying pathophysiological mechanism is left ventricular diastolic dysfunction: chronic systemic microvascular endothelial inflammation drives myocyte hypertrophy, hypophosphorylation of the structural giant protein titin (increasing passive myocyte stiffness), and extensive interstitial collagen deposition and myocardial fibrosis. Consequently, left ventricular chamber compliance is markedly reduced; during diastole, filling requires markedly elevated filling pressures (LVEDP and PCWP), which are transmitted retrogradely into the left atrium and pulmonary circulation, provoking pulmonary venous congestion and exertional dyspnea. Echocardiographic criteria include: left atrial enlargement (LAVI $>34$ mL/m$^2$), increased LV wall thickness, elevated $E/e'$ ratio ($>14$, indicating elevated filling pressures), and tricuspid regurgitation velocity $>2.8$ m/s. Serum natriuretic peptides (BNP and NT-proBNP) are elevated, though often lower than in HFrEF (and falsely suppressed by obesity). Validated diagnostic algorithms include the H2FPEF and HFA-PEFF scores. According to modern guidelines, foundational therapeutics proven to reduce cardiovascular death and hospitalizations comprise SGLT2 inhibitors (dapagliflozin, empagliflozin), mineralocorticoid receptor antagonists (spironolactone), ARNI (sacubitril/valsartan), and loop diuretics for fluid decongestion.

\medterm{Heart Failure with Reduced Ejection Fraction (HFrEF)}-- A clinical syndrome of heart failure characterized by left ventricular systolic dysfunction with a reduced ejection fraction (\(\le\)40\%). Causes: ischemic heart disease (prior MI, 60-70\%), dilated cardiomyopathy, hypertension, valvular heart disease. Pathophysiology: eccentric remodeling, progressive ventricular dilation, neurohormonal activation (RAAS, sympathetic nervous system). Clinical features: fatigue, dyspnea, orthopnea, PND, bibasilar rales, S3 gallop, elevated JVP, peripheral edema. Diagnosis: echocardiogram (LVEF \(\le\)40\%), elevated BNP/NT-proBNP. Guideline-Directed Medical Therapy (GDMT 4 pillars): ARNI (sacubitril/valsartan) or ACEi/ARB, beta-blocker (carvedilol, metoprolol succinate, bisoprolol), mineralocorticoid receptor antagonist (spironolactone, eplerenone), and SGLT2 inhibitor (dapagliflozin/empagliflozin). Device therapy: ICD for primary prevention if LVEF \(\le\)35\% despite \(\ge\)3 months GDMT; CRT for LBBB with QRS \(\ge\)150 ms.

\medterm{Heart Health}-- The state of cardiovascular well-being, characterized by efficient pumping function, clear arteries, normal blood pressure, and regular heart rhythm. Maintaining heart health involves a heart-healthy diet, regular aerobic and resistance exercise, weight management, stress reduction, and avoidance of tobacco. Preventive cardiology emphasizes early detection and management of risk factors to reduce the burden of cardiovascular disease.

\medterm{Heart Murmur}-- An audible extra sound produced by turbulent blood flow within the heart or great vessels, heard on auscultation. Murmurs are described by: timing (systolic, diastolic, continuous), intensity (Levine grading 1-6), location (apex, left sternal border, right upper sternal border), radiation (axilla, carotids, back), pitch (high or low frequency), quality (blowing, rumbling, harsh, musical), and response to maneuvers (inspiration, Valsalva, squatting, standing). Systolic murmurs: (1) Aortic stenosis (ejection systolic, crescendo-decrescendo, right upper sternal border, radiating to carotids); (2) Mitral regurgitation (holosystolic/pansystolic, blowing, apex, radiating to axilla); (3) Mitral valve prolapse (mid-systolic click + late systolic murmur); (4) Hypertrophic obstructive cardiomyopathy (ejection systolic, left sternal border, increases with Valsalva); (5) Ventricular septal defect (holosystolic, pansystolic, left lower sternal border). Diastolic murmurs: (1) Aortic regurgitation (early diastolic, decrescendo, high-pitched blowing, left sternal border); (2) Mitral stenosis (mid-diastolic, low-pitched rumbling, apex, accentuated by left lateral decubitus, presystolic accentuation). Continuous murmurs: patent ductus arteriosus (machinery murmur, left infraclavicular area). Innocent murmurs: Still murmur (vibratory, left lower sternal border, common in children), pulmonary flow murmur, venous hum. Workup: echocardiogram (transthoracic echo, bubble study, transoesophageal echo if needed) to identify valvular abnormality, assess severity, and determine need for intervention.

\medterm{Heart Rate}-- The number of cardiac contractions (heartbeats) per minute. Normal resting heart rate (sinus rhythm): 60-100 bpm in adults; higher in children (70-120 at age 1-2 years, 80-130 in neonates). Heart rate is determined by the intrinsic firing rate of the sinoatrial (SA) node (60-100 bpm), modulated by the autonomic nervous system: sympathetic activation (noradrenaline increases heart rate through beta-1 adrenergic receptors, positive chronotropy) and parasympathetic activation (vagus nerve, acetylcholine acts on M2 receptors, negative chronotropy). Tachycardia: heart rate >100 bpm at rest. Causes: sinus tachycardia (physiological response to exercise, stress, fever, pain, hypovolaemia, anaemia, hyperthyroidism, PE, MI, heart failure), re-entrant tachycardias (AVNRT, AVRT, atrial flutter, atrial fibrillation), ventricular tachycardia. Bradycardia: heart rate <60 bpm. Causes: sinus bradycardia (physiological in athletes, sleep), sick sinus syndrome, heart block (first, second, third degree), drug effect (beta-blockers, calcium channel blockers, digoxin), hypothyroidism, increased intracranial pressure, hypothermia. Assessment: pulse palpation (radial, carotid, apical), ECG (lead II rhythm strip for most accurate rate calculation: 300 / number of large squares between R waves; or number of R waves x 6 for a 10-second strip).

\medterm{Heart Sounds}-- Audible sounds produced by the heart during the cardiac cycle, normally heard as two sounds (lub-dub). S1 (lub) is caused by closure of the mitral and tricuspid valves at the beginning of systole. S2 (dub) is caused by closure of the aortic and pulmonary valves at the end of systole. Additional sounds (S3, S4) may indicate pathology. Heart sounds are auscultated at specific chest wall locations.

\medterm{Heart Surgery}-- Surgical procedures performed on the heart or great vessels to treat structural or functional cardiac conditions. Common operations include coronary artery bypass grafting (CABG), valve repair or replacement, and repair of congenital heart defects. Surgical approaches range from traditional open-heart surgery with sternotomy to minimally invasive and robot-assisted techniques.

\medterm{Heart Valve Disorders}-- Disorders affecting any of the four heart valves, including stenosis (narrowing) and regurgitation (leaking), which impair blood flow through the heart. Causes include congenital malformations, age-related degeneration, rheumatic fever, and endocarditis. Treatment ranges from medication monitoring to valve repair or replacement via surgery or transcatheter interventions.

\medterm{Heart Valves}-- The four cardiac valves ensure unidirectional blood flow through the heart. Atrioventricular valves: mitral valve (bicuspid, between left atrium and left ventricle, two leaflets — anterior and posterior, attached via chordae tendineae to papillary muscles preventing prolapse) and tricuspid valve (between right atrium and right ventricle, three leaflets). Semilunar valves: aortic valve (between left ventricle and aorta, three cusps — left coronary, right coronary, non-coronary) and pulmonary valve (between right ventricle and pulmonary artery, three cusps). Valve dysfunction: stenosis (failure to open fully, causing pressure overload) and regurgitation (failure to close fully, causing volume overload). Common pathologies: aortic stenosis (calcific, bicuspid valve, rheumatic), mitral regurgitation (mitral valve prolapse, ischemic, rheumatic), and infective endocarditis (vegetations on valve leaflets, typically Streptococcus viridans, Staphylococcus aureus). Clinical assessment: auscultation (heart sounds and murmurs), echocardiography (transthoracic and transoesophageal) for anatomy and hemodynamics.

\medterm{Heartburn}-- A subjective symptom characterized by a burning sensation or retrosternal pain arising in the epigastrium and radiating upward toward the neck, throat, or jaw, typically resulting from the reflux of acidic gastric contents across an incompetent lower esophageal sphincter (LES) into the esophagus (Gastroesophageal Reflux Disease, GERD). In clinical cardiology, heartburn is of paramount importance because gastroesophageal and cardiac sensory afferent nerve fibers share identical common spinal pathways: visceral sensory pain fibers from both the heart and the esophagus converge upon the same dorsal root ganglion cells and thoracic spinal cord segments ($T_1$ through $T_4--T_5$). Consequently, the cerebral cortex cannot reliably differentiate between myocardial ischemia and esophageal mucosal irritation or spasm. Acute myocardial ischemia, angina pectoris, and acute myocardial infarction (particularly inferior wall MI, which irritates the adjacent diaphragm) frequently masquerade clinically as 'heartburn', indigestion, or dyspepsia. Crucially, a positive clinical response to sublingual nitroglycerin or an antacid/viscous lidocaine 'GI cocktail' CANNOT be used to differentiate esophageal pain from coronary artery disease, as nitroglycerin and calcium channel blockers relax esophageal smooth muscle (relieving diffuse esophageal spasm) in addition to dilating coronary arteries, and placebo effects are common. In any patient presenting with acute retrosternal burning pain, particularly those with cardiovascular risk factors (age $>40$, diabetes, hypertension, smoking, dyslipidemia), life-threatening Acute Coronary Syndromes MUST be definitively excluded with immediate 12-lead electrocardiography and serial cardiac troponins before attributing symptoms to benign gastroesophageal reflux.

\medterm{Heart-Healthy Eating}-- A dietary approach emphasizing foods that support cardiovascular function and reduce the risk of heart disease. Key principles include limiting saturated fats, trans fats, sodium, and added sugars while prioritizing fruits, vegetables, whole grains, lean proteins, and healthy fats like those from olive oil and fatty fish. This pattern aligns closely with the Mediterranean and DASH dietary guidelines.

\medterm{Hemochromatosis Cardiomyopathy}-- Iron overload cardiomyopathy from hereditary (HFE gene mutations, most common) or secondary causes. Iron deposition in the myocardium causes restrictive then dilated cardiomyopathy. Diagnosis: elevated transferrin saturation (>45\%), ferritin, and genetic testing. Cardiac MRI T2* <20 ms indicates myocardial iron overload. Treatment: phlebotomy for primary hemochromatosis, iron chelation (deferoxamine) for secondary. Early treatment can reverse cardiac dysfunction.

\medterm{Heparin}-- A fast-acting, highly sulfated glycosaminoglycan parenteral anticoagulant widely used for the acute treatment and prophylaxis of arterial and venous thromboembolic disease. Unfractionated heparin (UFH, molecular weight 3--30 kDa) binds with high affinity to the endogenous serine protease inhibitor antithrombin (antithrombin III) via a unique pentasaccharide sequence, inducing an allosteric conformational change that accelerates antithrombin's rate of inhibition of activated factor X (FXa) and thrombin (factor IIa) by $>1,000$-fold. To inactivate thrombin, heparin must form a ternary bridge contacting both antithrombin and thrombin simultaneously (requiring $\ge 18$ saccharide units), whereas factor Xa inhibition requires only pentasaccharide binding. UFH is administered intravenously or subcutaneously and monitored using the activated partial thromboplastin time (aPTT, target 1.5--2.5 times control) or anti-factor Xa levels (0.3--0.7 IU/mL). In contrast, Low-Molecular-Weight Heparin (LMWH, e.g., enoxaparin, dalteparin, ~4.5 kDa) preferentially inactivates factor Xa over thrombin (ratio 3:1 to 4:1) with predictable pharmacokinetics that eliminate routine lab monitoring. Indications include acute coronary syndromes (PCI, STEMI, NSTEMI), pulmonary embolism, deep vein thrombosis, and cardiopulmonary bypass. Critical adverse effects include bleeding, osteoporosis with prolonged use, and Heparin-Induced Thrombocytopenia (HIT Type II), an immune-mediated limb- and life-threatening thrombophilia caused by IgG antibodies targeting platelet factor 4 (PF4)-heparin complexes, causing platelet activation and consumption ($>50\%$ drop in platelets). UFH is rapidly reversed by intravenous protamine sulfate (1 mg protamine per 100 units of heparin).

\medterm{Hepatic Drug Metabolism}-- Phase I reactions (oxidation, reduction, hydrolysis) primarily by CYP450 enzymes introduce or expose functional groups. Phase II reactions (conjugation) with glucuronide, sulfate, or acetyl groups increase water solubility for renal excretion. Hepatic metabolism can activate prodrugs, inactivate active drugs, or produce toxic metabolites. Liver disease reduces metabolism, requiring dose adjustment. Enzyme induction increases metabolism (decreasing drug levels); enzyme inhibition decreases metabolism (increasing drug levels).

\medterm{Hepatojugular Reflux}-- Sustained pressure applied over the liver for 15-30 seconds causes a sustained rise in JVP of more than 3 cm above baseline. It indicates right heart failure or fluid overload, reflecting the inability of the right ventricle to accommodate increased venous return. Also seen in cardiac tamponade, constrictive pericarditis, and tricuspid regurgitation. It is a useful bedside test for volume status assessment.

\medterm{Holosystolic Murmur Differential}-- Pansystolic (holosystolic) murmurs are plateau-shaped and extend from S1 to S2. Differential includes: mitral regurgitation (apex, radiates to axilla, increases with handgrip), tricuspid regurgitation (left lower sternal border, increases with inspiration), ventricular septal defect (left lower sternal border, may have thrill), and functional murmurs from increased flow (anemia, thyrotoxicosis, pregnancy).

\medterm{Hyperlipidemia}-- Elevated levels of lipids in the blood—total cholesterol, LDL, triglycerides, and/or reduced HDL—a major modifiable risk factor for atherosclerosis and cardiovascular disease. Primary (genetic: familial hypercholesterolemia) or secondary (hypothyroidism, diabetes, nephrotic syndrome, obesity, alcohol). Screening: fasting lipid panel; risk calculators (ASCVD) guide treatment. Management: lifestyle modification, statins (first-line), ezetimibe, PCSK9 inhibitors, fibrates/omega-3 for triglycerides. Targets depend on cardiovascular risk category.

\medterm{Hypertension}-- A chronic, highly prevalent cardiovascular condition characterized by persistently elevated systemic arterial blood pressure that imposes chronic mechanical stress on the heart and arterial tree, serving as the leading modifiable risk factor for stroke, coronary artery disease, heart failure, chronic kidney disease, and vascular dementia. According to the authoritative 2017 ACC/AHA guidelines, blood pressure categories in adults are defined based on average seated resting clinic measurements: (1) Normal: SBP $<120$ mmHg AND DBP $<80$ mmHg; (2) Elevated Blood Pressure: SBP 120--129 mmHg AND DBP $<80$ mmHg; (3) Stage 1 Hypertension: SBP 130--139 mmHg OR DBP 80--89 mmHg; (4) Stage 2 Hypertension: SBP $\ge 140$ mmHg OR DBP $\ge 90$ mmHg; and (5) Hypertensive Crisis: SBP $>180$ mmHg and/or DBP $>120$ mmHg, divided into Hypertensive Urgency (no acute end-organ damage) and Hypertensive Emergency (evidence of acute, life-threatening target organ damage, including hypertensive encephalopathy, acute aortic dissection, acute pulmonary edema, acute MI, or acute renal failure, requiring ICU admission and IV antihypertensive titration). Approximately 90--95\% of cases represent Essential (Primary) Hypertension, a multifactorial disorder arising from interacting genetic polymorphisms and lifestyle factors (excess sodium intake, obesity, physical inactivity, alcohol, stress). The remaining 5--10\% represent Secondary Hypertension, which should be investigated in young patients ($<30$ years) or resistant hypertension; causes include renal parenchymal disease (most common), renovascular hypertension (renal artery stenosis), primary hyperaldosteronism, obstructive sleep apnea, pheochromocytoma, Cushing syndrome, and coarctation of the aorta. First-line pharmacological therapy consists of four classes: thiazide diuretics (chlorthalidone, hydrochlorothiazide), ACE inhibitors, ARBs, and dihydropyridine calcium channel blockers.

\medterm{Hypertensive Crisis}-- Severe elevation in blood pressure (SBP >180 or DBP >120) requiring urgent medical attention. Two types: hypertensive emergency (with end-organ damage) and hypertensive urgency (no acute end-organ damage). Hypertensive urgency: can be managed with oral antihypertensives (captopril, clonidine, labetalol) and careful outpatient follow-up. Distinction is critical for appropriate management.

\medterm{Hypertensive Emergency}-- Severe hypertension (SBP >180 or DBP >120) with acute target organ damage. Organs involved: brain (hypertensive encephalopathy, stroke), heart (acute coronary syndrome, acute pulmonary edema, aortic dissection), kidneys (acute kidney injury), and eyes (papilledema). Requires immediate BP reduction with IV antihypertensives (nicardipine, labetalol, nitroprusside) and ICU monitoring. Goal: reduce MAP by 10-20\% in the first hour, then cautious further reduction.

\medterm{Hypertensive Heart Disease}-- Left ventricular hypertrophy and diastolic/systolic dysfunction resulting from chronic hypertension. LVH is an independent cardiovascular risk factor and precursor to heart failure. Diagnosis by ECG (voltage criteria) or echocardiography (wall thickness >1.1 cm). Treatment of hypertension can regress LVH. Complications include heart failure (HFpEF initially, then HFrEF), atrial fibrillation, and sudden cardiac death. Beta-blockers, ACE inhibitors, and ARBs can regress LVH.

\medterm{Hypertrophic Cardiomyopathy}-- A common, primary inherited genetic heart disease (prevalence ~1 in 500) characterized by unexplained left ventricular hypertrophy (LV wall thickness $\ge 15$ mm in adults, or $\ge 13$ mm in individuals with a positive family history) in the absence of another cardiac, systemic, or metabolic etiology capable of producing that magnitude of hypertrophy (such as severe aortic stenosis or uncontrolled hypertension). Transmitted in an autosomal dominant pattern with variable penetrance, HCM is predominantly caused by pathogenic variants in genes encoding cardiac sarcomeric proteins, most commonly MYH7 ($\beta$-myosin heavy chain) and MYBPC3 (myosin-binding protein C), accounting for $>70\%$ of genetically identified cases. Histological hallmarks include marked cardiomyocyte hypertrophy, extensive myocyte disarray (chaotic branching and loss of parallel cell alignment), microvascular remodeling, and replacement interstitial fibrosis. In approximately 70\% of patients, dynamic Left Ventricular Outflow Tract (LVOT) obstruction is present at rest or provoked by physiological stress (defined as a peak instantaneous LVOT gradient $\ge 30$ mmHg, with severe obstruction $\ge 50$ mmHg). The obstruction is driven by asymmetric septal hypertrophy alongside Systolic Anterior Motion (SAM) of the mitral valve, where the anterior mitral leaflet is pulled toward the septum by Venturi forces and flow drag, simultaneously causing dynamic subaortic stenosis and eccentric posterior mitral regurgitation. Auscultation reveals a harsh, crescendo-decrescendo systolic murmur at the left sternal border that uniquely intensifies with maneuvers that reduce LV preload or afterload (Valsalva strain, standing) and softens with increased preload or afterload (squatting, passive leg raise, handgrip). HCM is the single most common cause of sudden cardiac death (SCD) in young competitive athletes, provoked by ventricular tachycardia/fibrillation. Pharmacotherapy begins with non-vasodilating beta-blockers (titrated to suppress exertional gradients), non-dihydropyridine calcium channel blockers (verapamil), and the novel cardiac myosin ATPase inhibitor Mavacamten. In refractory severe symptomatic LVOT obstruction, invasive septal reduction therapies comprise surgical septal myectomy (Morrow procedure) or percutaneous Alcohol Septal Ablation (ASA). High-risk patients receive an Implantable Cardioverter-Defibrillator (ICD) for SCD prevention.

\medterm{Hypoplastic Left Heart Syndrome}-- A severe congenital heart defect where the left heart structures (mitral valve, left ventricle, aortic valve, ascending aorta) are underdeveloped and unable to support systemic circulation. The right ventricle pumps blood to both the pulmonary and systemic circulations via a patent ductus arteriosus. Presents in the first days of life with cyanosis, tachypnea, and cardiovascular collapse when the ductus closes. Diagnosis: echocardiography shows a tiny left ventricle, hypoplastic ascending aorta, and ductal-dependent systemic circulation. Staged surgical palliation: Norwood procedure (stage 1, neonatal), Glenn shunt (stage 2, 4-6 months), and Fontan completion (stage 3, 2-4 years). Prognosis has improved significantly but remains guarded. Without surgery, it is uniformly fatal.

\medterm{Hypotension}-- Abnormally low blood pressure, generally systolic <90 mmHg or mean arterial pressure <65 mmHg. Causes include hypovolemia, cardiac failure, sepsis, anaphylaxis, and medications. Orthostatic hypotension is a drop upon standing. Severe hypotension (shock) requires urgent fluid resuscitation and vasopressors. Chronic asymptomatic hypotension is usually benign.

\medterm{ICD Shock Therapy}-- Anti-tachycardia pacing (ATP) for slow VT delivers rapid pacing pulses to terminate the arrhythmia painlessly. Cardioversion shock synchronously terminates VT. Defibrillation shock treats VFib asynchronously. The MADIT-RIT trial showed that higher-rate programming (detection >180-200 bpm) reduced inappropriate shocks and mortality. S-ICD (subcutaneous ICD) avoids transvenous leads but lacks ATP and pacing capability. ICD shocks are associated with psychological morbidity.

\medterm{Idioventricular Rhythm}-- A slow, regular ventricular rhythm that arises from an ectopic pacemaker focus situated within the specialized His-Purkinje system or ventricular myocardium, operating at its intrinsic physiological rate of 20 to 40 beats per minute. An idioventricular rhythm functions as a vital tertiary safety-net escape pacemaker, emerging when higher supraventricular pacemakers (the sinoatrial node and atrioventricular junction) fail to fire (sinus arrest, extreme sinus bradycardia) or when their electrical impulses are completely blocked from reaching the ventricles (third-degree / complete AV block). Because ventricular activation originates within the ventricular myocardium outside the rapid Purkinje network and conducts slowly from cell to cell via gap junctions, the 12-lead ECG is characterized by: (1) Slow ventricular rate (20--40 bpm); (2) Wide, bizarre QRS complexes ($\ge 120$ ms); (3) T waves that are typically discordant (opposite polarity to the QRS complex); and (4) Complete atrioventricular dissociation (independent P waves march through the rhythm at a faster atrial rate, unrelated to QRS complexes). When the intrinsic ventricular rate exceeds 40 bpm but remains $<100--110$ bpm, the rhythm is termed Accelerated Idioventricular Rhythm (AIVR), a benign rhythm driven by abnormal automaticity that classically appears as a 'reperfusion arrhythmia' following successful mechanical (PCI) or pharmacological thrombolytic recanalization of an occluded coronary artery during acute STEMI. True slow idioventricular escape rhythms produce severe bradycardia, presyncope, syncope (Stokes-Adams attacks), and hypotension, requiring urgent chronotropic support (atropine, isoproterenol) and temporary or permanent transvenous pacemaker implantation.

\medterm{Impella Device}-- Percutaneous microaxial left ventricular assist device that actively pumps blood from the LV into the aorta, reducing LV workload and increasing cardiac output. Types: Impella 2.5 (up to 2.5 L/min), Impella CP (up to 3.7 L/min), Impella 5.0 (up to 5 L/min, surgically placed). Used in cardiogenic shock, high-risk PCI, and as bridge to decision. Complications: hemolysis, limb ischemia, device malposition, bleeding.

\medterm{Implantable Cardioverter-Defibrillator}-- An implanted device that monitors heart rhythm and delivers electrical therapy (antitachycardia pacing or defibrillation) for life-threatening ventricular arrhythmias. Indications include secondary prevention (survived cardiac arrest, sustained VT) and primary prevention (EF 35\% or less with NYHA II-III symptoms despite optimal medical therapy, or post-MI with EF 35\% or less after 40 days). ICDs reduce sudden cardiac death and overall mortality in selected patients.

\medterm{Implantable Loop Recorder}-- Subcutaneous device providing continuous ECG monitoring for up to 3 years. Used for evaluation of unexplained syncope, recurrent palpitations, cryptogenic stroke, and occult atrial fibrillation detection. Small size (coin-sized), inserted under local anesthesia in the chest. Automatically records and stores rhythm data during symptoms or detected arrhythmias. Increasingly important in the workup of cryptogenic stroke for detecting silent AFib.

\medterm{Infective Endocarditis}-- Infection of the endocardial surface of the heart, most commonly affecting the heart valves. It is diagnosed using the Duke criteria, which include major criteria (positive blood cultures, echocardiographic evidence of vegetations) and minor criteria (predisposing heart condition, fever, vascular phenomena, immunologic phenomena). Treatment requires prolonged IV antibiotic therapy (4-6 weeks). Surgery is indicated for complications including heart failure, persistent infection, and embolic events.

\medterm{Interatrial Septum}-- The wall separating the right and left atria. It contains the fossa ovalis, a remnant of the foramen ovale from fetal circulation. The interatrial septum has a thinner area where transseptal puncture can be performed during cardiac catheterization. Atrial septal defects (ASD) are congenital openings that allow left-to-right shunting, potentially causing right heart failure.

\medterm{Interventricular Septum}-- The thick muscular wall separating the left and right ventricles. It has two portions: the muscular (inferior, thicker) and membranous (superior, thinner) portions. The septum contains parts of the cardiac conduction system, including the bundle of His. Ventricular septal defects (VSDs) are congenital openings in the septum that allow left-to-right shunting of blood.

\medterm{Intra-Aortic Balloon Pump}-- A mechanical circulatory support device that increases myocardial perfusion and reduces afterload. The balloon inflates during diastole (augmenting coronary perfusion) and deflates during systole (reducing afterload). Inserted percutaneously into the femoral artery. Indications include cardiogenic shock, mechanical complications of MI (papillary muscle rupture, VSD), and high-risk PCI. IABP provides modest hemodynamic support and is being replaced by more effective devices.

\medterm{Intravascular lymphoma}-- A rare aggressive B-cell non-Hodgkin lymphoma characterized by proliferation of neoplastic lymphocytes within the lumina of small blood vessels, predominantly affecting the skin and central nervous system. Presents with livedo reticularis, skin nodules, neurological deficits, and fever of unknown origin.

\medterm{Intravascular Ultrasound}-- Catheter-based imaging using a miniature ultrasound probe inside coronary arteries to visualize vessel wall layers, plaque composition, and stent deployment. Provides more detailed information than angiography alone, including vessel remodeling, plaque burden, and true vessel diameter. Used to optimize stent deployment (apposition, expansion) and assess intermediate angiographic lesions. OCT (optical coherence tomography) provides higher resolution imaging.

\medterm{Ischemic Cardiomyopathy}-- Dilated cardiomyopathy resulting from extensive myocardial infarction with significant scar burden. EF is typically severely reduced (<35\%). The infarcted myocardium is replaced by fibrotic scar tissue, impairing contractile function and promoting ventricular arrhythmias. Viability testing (PET, dobutamine echo, cardiac MRI) may identify hibernating myocardium that could benefit from revascularization. Treatment includes standard heart failure therapy, ICD for arrhythmia prevention, and revascularisation when viable myocardium is present.

\medterm{Ischemic gangrene}-- Death of tissue due to prolonged severe arterial insufficiency, most commonly in the extremities. Presents as dry, black, mummified tissue with a clear line of demarcation from viable tissue. Requires revascularization when possible and surgical debridement or amputation of necrotic tissue.

\medterm{Janeway Lesions}-- Painless, erythematous or hemorrhagic macules on the palms and soles, caused by septic microemboli in infective endocarditis. They are a minor criterion in the Duke criteria. Unlike Osler nodes, Janeway lesions are painless and represent microabscesses in the dermis. They are transient and may disappear quickly. Both Osler nodes and Janeway lesions indicate active endocarditis and guide diagnosis.

\medterm{Killip Classification}-- A clinical classification system for assessing the severity of myocardial infarction based on physical examination findings. Class I: no heart failure signs. Class II: mild heart failure (crackles, S3 gallop, elevated JVD). Class III: pulmonary edema. Class IV: cardiogenic shock. The Killip class correlates with in-hospital mortality: Class I (6\%), Class II (17\%), Class III (38\%), Class IV (81\%). It guides initial management intensity.

\medterm{Kussmaul Sign}-- Paradoxical increase in jugular venous pressure with inspiration, normally the JVP decreases with inspiration. It indicates impaired right ventricular filling during inspiration. Seen in constrictive pericarditis, restrictive cardiomyopathy, severe right heart failure, cardiac tamponade, and massive pulmonary embolism. It is a reflection of ventricular interdependence and loss of normal pericardial compliance.

\medterm{LDL}-- Low-density lipoprotein; the 'bad' cholesterol that transports cholesterol to peripheral tissues and is the principal driver of atherosclerosis. Elevated LDL is a major modifiable cardiovascular risk factor; guidelines target LDL levels based on global risk (statin intensity). Measurement: direct or Friedewald equation (LDL = total cholesterol - HDL - triglycerides/5). Management: lifestyle (diet, exercise), statins, ezetimibe, PCSK9 inhibitors, bempedoic acid. Each 1 mmol/L (~39 mg/dL) LDL reduction lowers major events ~20-25 percent.

\medterm{Left Anterior Descending Artery}-- Also designated the anterior interventricular artery (LAD), the premier and most clinically critical branch of the left coronary artery, often colloquially termed the 'widow maker' due to the catastrophic mortality associated with its acute proximal thrombotic occlusion. Arising from the bifurcation of the left main coronary artery behind the pulmonary trunk, the LAD travels inferiorly along the anterior interventricular sulcus toward the cardiac apex, which it typically wraps around to anastomose with the posterior descending artery. The LAD gives off two primary sets of branches: (1) Septal perforator branches (first septal is the largest), which dive deep into the myocardium to supply the anterior two-thirds of the interventricular septum, including the bundle of His and bundle branches of the cardiac conduction system; and (2) Diagonal branches ($D_1, D_2$), which course laterally across the anterolateral free wall of the left ventricle. Overall, the LAD perfuses approximately 45--55\% of the total left ventricular myocardial mass (including the anterior wall, anteroseptal region, cardiac apex, and the anterolateral papillary muscle). Acute occlusion produces an extensive Anterior or Anteroseptal ST-Elevation Myocardial Infarction (STEMI), characterized by ST-segment elevation in precordial leads $V_1$ through $V_4$ (and limb leads I and aVL if high lateral diagonal involvement occurs). Complications of LAD infarction include cardiogenic shock, acute pulmonary edema, ventricular free-wall rupture, ventricular septal defect (VSD), ventricular apical aneurysm with mural thrombus formation, and new-onset right bundle branch block (RBBB) or bifascicular block.

\medterm{Left Atrial Myxoma}-- Most common primary cardiac tumor (75\%), typically originating from the interatrial septum in the left atrium. Presents with embolic events (25\%), intracardiac obstruction (mitral valve obstruction mimicking mitral stenosis), and constitutional symptoms (fever, weight loss, elevated ESR mimicking infection or malignancy). Echocardiography shows a mobile, pedunculated mass. Surgical excision is curative. Recurrence is rare but possible, especially in familial syndromes (Carney complex).

\medterm{Left Atrium}-- The posterior-most cardiac chamber, situated directly anterior to the esophagus and descending thoracic aorta, receiving oxygenated blood from the pulmonary circulation via four pulmonary veins (left and right superior and inferior veins) that lack valves. Normal resting left atrial pressure averages 6 to 12 mmHg (reflected clinically as the pulmonary capillary wedge pressure, PCWP). During diastole, blood flows across the open mitral valve into the left ventricle; in late diastole, coordinated left atrial contraction ('atrial kick') contributes approximately 20--30\% of left ventricular end-diastolic filling volume, which becomes critical for cardiac output in patients with non-compliant ventricles (aortic stenosis, hypertrophic cardiomyopathy). Pathological left atrial enlargement ($>34$ mL/m$^2$ volume index) develops secondary to chronic pressure overload (mitral stenosis, systemic hypertension) or volume overload (mitral regurgitation). Left atrial dilatation disrupts electrical architecture, serving as the primary anatomical substrate for the initiation and maintenance of atrial fibrillation. Furthermore, left atrial stasis, particularly within the finger-like left atrial appendage (LAA), promotes blood stagnation and thromboembolism, accounting for $>90\%$ of non-valvular atrial fibrillation-related ischemic strokes; percutaneous LAA closure (e.g., Watchman device) or surgical excision mitigates this risk.

\medterm{Left Bundle Branch}-- A continuation of the bundle of His that divides into anterior and posterior fascicles. The anterior fascicle supplies the anterior and superior left ventricular wall; the posterior fascicle supplies the inferior and posterior wall. Left bundle branch block (LBBB) causes delayed left ventricular activation and is commonly associated with hypertension, coronary artery disease, and cardiomyopathy. LBBB on ECG shows widened QRS with notched R waves in lateral leads.

\medterm{Left Main Coronary Artery}-- The primary arterial conduit (LMCA) supplying the vast majority of the left ventricular myocardium, arising directly from the ostium in the left posterior (left) aortic sinus of Valsalva, immediately above the aortic valve leaflet. The LMCA typically measures 3.5 to 5.0 mm in caliber and courses for a short distance (averaging 5 to 25 mm, though occasionally absent in congenital dual-ostium variants) between the left atrial appendage and the pulmonary trunk before bifurcating into the left anterior descending (LAD) and left circumflex (LCx) arteries (or trifurcating with an intermediate ramus branch). Because the LMCA supplies $>70--75\%$ of the total left ventricular muscle mass (and nearly 100\% in left-dominant circulations), hemodynamically significant stenosis ($\ge 50\%$ luminal diameter reduction) represents an exceptionally high-risk anatomical scenario with grave prognostic implications. Clinically, LMCA ischemia presents with severe angina, sudden acute pulmonary edema, or cardiogenic shock; on a 12-lead resting or stress ECG, acute subtotal LMCA occlusion typically manifests as widespread, deep horizontal or downsloping ST-segment depressions across multiple leads ($\ge 6--8$ leads, particularly leads I, II, and $V_4$--$V_6$) accompanied by prominent reciprocal ST-segment elevation in lead aVR ($\ge 1$ mm, often exceeding ST elevation in $V_1$). Coronary artery bypass grafting (CABG) has historically been the established gold-standard revascularization strategy, though percutaneous coronary intervention (PCI) with drug-eluting stents is increasingly performed in patients with low-to-intermediate SYNTAX scores.

\medterm{Left Ventricle}-- The primary thick-walled muscular pumping chamber of the heart that ejects oxygenated blood under high pressure into the systemic arterial circulation. In healthy adults, the left ventricular wall thickness measures 8 to 11 mm in diastole, generating peak systolic pressures of 100 to 140 mmHg against systemic afterload, with normal end-diastolic pressures of 3 to 12 mmHg. The chamber exhibits a conical geometry with a spiraling myocardial fiber orientation that produces a coordinated wringing or torsional motion during systole, ejecting a normal stroke volume of 60 to 100 mL (corresponding to a normal left ventricular ejection fraction [LVEF] of 52--72\% in men and 54--74\% in women). The LV wall is anchored by two prominent papillary muscles (anterolateral, dual blood supply from LAD and LCx; and posteromedial, single blood supply from PDA, making it highly susceptible to ischemic rupture during inferior MI). According to Laplace's Law ($\text{Wall Stress} = [P \times r] / [2h]$), chronic pressure overload (aortic stenosis, chronic hypertension) induces concentric hypertrophy (increased wall thickness $h$ with normal chamber radius $r$ to normalize systolic wall stress), whereas chronic volume overload (mitral or aortic regurgitation) induces eccentric hypertrophy (chamber dilation). Progressive systolic or diastolic LV dysfunction drives heart failure syndromes (HFrEF and HFpEF).

\medterm{Left Ventricular Noncompaction}-- A rare cardiomyopathy characterized by prominent myocardial trabeculations with deep intertrabecular recesses, thought to result from arrested myocardial compaction during embryogenesis. It may present with heart failure, arrhythmias, or thromboembolism. Diagnosis by echocardiography or cardiac MRI showing a noncompacted to compacted myocardial ratio >2.3 (echocardiography) or >2.2 (cardiac MRI). Treatment includes heart failure therapy, anticoagulation for thromboembolism prevention, and arrhythmia management.

\medterm{Left-sided superior vena cava}-- A congenital systemic thoracic venous anomaly (also designated persistent left superior vena cava, PLSVC) resulting from the embryological failure of regression of the left anterior cardinal vein and common cardinal vein (the ligament of Marshall) during early cardiac development. Occurring in approximately 0.3--0.5\% of the general healthy population and up to 3--10\% of patients with congenital heart defects (such as atrial septal defects, tetralogy of Fallot, or coarctation of the aorta), PLSVC represents the most common thoracic venous malformation. In $>90\%$ of cases, a normal right superior vena cava is also present (bilateral SVCs, connected by a bridging innominate vein). The left-sided SVC courses vertically downward anterior to the left pulmonary hilum and the aortic arch, penetrates the pericardium, and drains directly into a markedly dilated Coronary Sinus (CS), which subsequently empties into the right atrium; hemodynamically, this creates no physiological shunt and patients remain entirely asymptomatic. However, in rare variants (~10\%), the PLSVC drains anomalously into the left atrium or the left superior pulmonary vein (either directly or via an unroofed coronary sinus), producing a right-to-left shunt with desaturation, cyanosis, and heightened risk for paradoxical systemic thromboembolism or brain abscess. Clinical recognition is critical during interventional and surgical procedures: left subclavian vein cannulation leads to atypical catheter positioning descending into the left mediastinum, transvenous pacemaker or ICD lead placement via the coronary sinus is technically challenging and risks CS dissection, and during cardiac surgery requiring cardiopulmonary bypass, retrograde cardioplegia can cause ineffective myocardial arrest and massive systemic air embolization. Transthoracic echocardiography typically reveals an unexplained giant dilated coronary sinus ($>1$ cm), confirmed by agitated saline ('bubble') injection into a left arm vein showing rapid opacification of the coronary sinus before the right atrium.

\medterm{Lingual artery}-- A branch of the external carotid artery supplying the tongue, floor of mouth, and sublingual structures. It arises anterior to the hypoglossal nerve and courses deep to the hyoglossus muscle. Lingual artery aneurysm or thrombosis is rare but can cause life-threatening hemorrhage or tongue ischemia.

\medterm{Long QT Syndrome}-- A channelopathy characterized by prolonged QT interval (QTc >450 ms in men, >460 ms in women) with risk of torsades de pointes and sudden cardiac death. Acquired: drugs (antiarrhythmics, antipsychotics, antibiotics), electrolyte abnormalities (hypokalemia, hypomagnesemia), bradycardia. Congenital: LQT1 (exercise-triggered, KCNQ1), LQT2 (emotional/auditory, KCNH2), LQT3 (sleep, SCN5A). Diagnosis: ECG, stress test, genetic testing. Treatment: avoid QT-prolonging drugs, beta-blockers for LQT1/2, left cardiac sympathetic denervation, and ICD for high-risk patients.

\medterm{Marfan Syndrome}-- An autosomal dominant connective tissue disorder caused by mutations in the FBN1 gene (fibrillin-1). Cardiovascular manifestations include aortic root dilation and dissection, mitral valve prolapse, and aortic regurgitation. Skeletal features include tall stature, arachnodactyly, pectus deformity, and joint hypermobility. Aortic root >5.0 cm or rapid dilation warrants prophylactic aortic root replacement. Beta-blockers or ARBs reduce aortic dilation rate. Regular echocardiographic surveillance is recommended.

\medterm{Mean Arterial Pressure}-- The time-weighted average perfusion pressure driving systemic blood flow through the microvasculature across an entire cardiac cycle, serving as the primary physiological determinant of end-organ perfusion pressure. At resting heart rates (60--80 bpm), diastole comprises approximately two-thirds of the cardiac cycle, and systole comprises one-third; consequently, MAP is estimated clinically from non-invasive blood pressure measurements using the formula: $\text{MAP} \approx \text{DBP} + \frac{1}{3}(\text{SBP} - \text{DBP}) = \frac{2(\text{DBP}) + \text{SBP}}{3}$, where $(\text{SBP} - \text{DBP})$ represents the pulse pressure (at tachycardic rates, systole occupies a greater proportion, approaching a 50:50 ratio). Under continuous arterial line catheter monitoring, MAP is measured with greater precision by electronic integration of the area under the arterial pulse contour divided by cycle duration. Hemodynamically, MAP is governed by Ohm's Law analogue for the circulation: $\text{MAP} = (\text{CO} \times \text{SVR}) + \text{CVP}$, where CO is cardiac output, SVR is systemic vascular resistance, and CVP is central venous pressure. In critical care and resuscitation medicine (sepsis, cardiogenic shock, trauma), international guidelines mandate targeting a minimum $\text{MAP} \ge 65$ mmHg to preserve autoregulation of vital organ capillary beds (cerebral, coronary, and renal glomerular perfusion). Values consistently $<60$ mmHg result in cellular ischemia, acute tubular necrosis, and multi-organ failure.

\medterm{Mechanical Circulatory Support}-- Devices that provide temporary or permanent cardiac support. Options include intra-aortic balloon pump (IABP), Impella (percutaneous ventricular assist device), TandemHeart, extracorporeal membrane oxygenation (ECMO), and durable left ventricular assist devices (LVADs). Indications include cardiogenic shock, bridge to transplant, and bridge to recovery. LVADs (HeartMate 3) are increasingly used as destination therapy in advanced heart failure ineligible for transplant.

\medterm{Mid-Systolic Click}-- A sharp, high-pitched sound occurring in mid-to-late systole, associated with mitral valve prolapse. The click represents sudden tensing of the elongated chordae tendineae as the prolapsing leaflet billows into the left atrium. It may be followed by a late systolic murmur of mitral regurgitation. The click moves earlier with decreased preload (standing, Valsalva) and later with increased preload (squatting, leg elevation).

\medterm{Mineralocorticoid Receptor Antagonist}-- A class of potassium-sparing diuretics and neurohormonal antagonists comprising spironolactone and eplerenone (as well as the non-steroidal agent finerenone) that competitively inhibit intracellular mineralocorticoid (aldosterone) receptors in the principal cells of the renal cortical collecting duct and in extrarenal cardiovascular tissues. Under normal conditions, aldosterone binds cytoplasmic receptors to promote the synthesis and luminal insertion of epithelial sodium channels (ENaC) and basolateral $Na^+/K^+$-ATPases, stimulating sodium and water reabsorption in exchange for urinary potassium and hydrogen ion excretion. By blocking this receptor, MRAs promote mild natriuresis and diuresis while conserving potassium and hydrogen ions. Beyond the nephron, MRAs block aldosterone-mediated myocardial fibrosis, vascular endothelial dysfunction, and sympathetic autonomic activation. In heart failure with reduced ejection fraction (HFrEF, NYHA Class II--IV), addition of an MRA (spironolactone or eplerenone) reduces all-cause mortality by ~30\% and decreases hospitalizations (proven in the RALES and EPHESUS trials). MRAs are also indicated in post-MI LV dysfunction, resistant hypertension, and primary hyperaldosteronism (Conn syndrome). Spironolactone also acts as an antagonist at androgen and progesterone receptors, causing endocrine side effects in up to 10\% of men (painful gynecomastia, erectile dysfunction, decreased libido); eplerenone is highly selective for mineralocorticoid receptors and avoids these anti-androgenic effects. The primary, life-threatening adverse risk is hyperkalemia ($K^+ > 5.5$ mEq/L), particularly when combined with ACE inhibitors, ARBs, or ARNIs in patients with chronic kidney disease (eGFR $<30$ mL/min/1.73 m$^2$).

\medterm{MitraClip}-- A transcatheter device that creates a double orifice in the mitral valve by clipping the anterior and posterior leaflets together, reducing mitral regurgitation. It is indicated for patients with significant (3+ or 4+) mitral regurgitation who are at high risk for surgical mitral valve repair. The COAPT trial demonstrated benefit in selected patients with secondary mitral regurgitation. The procedure is performed under transesophageal echocardiographic guidance.

\medterm{Mitral Regurgitation}-- Incomplete closure of the mitral valve allowing blood to flow backward from the left ventricle to the left atrium during systole. Causes include mitral valve prolapse, papillary muscle dysfunction, rheumatic disease, and ischemic heart disease. Murmur: pansystolic (holosystolic) blowing murmur at the apex radiating to the axilla. Chronic MR leads to left atrial and ventricular dilation. Severe MR requires surgical repair or replacement.

\medterm{Mitral Stenosis}-- Narrowing of the mitral valve orifice obstructing diastolic blood flow from the left atrium into the left ventricle. Etiologies are predominantly post-inflammatory Rheumatic Heart Disease (responsible for $>85\%$ of cases globally, where recurrent cross-reactive anti-streptococcal autoimmunity produces fibrous thickening, commissural fusion, chordal shortening, and calcification) and degenerative severe Mitral Annular Calcification (MAC, presenting in elderly patients). Normal mitral valve orifice area is 4.0 to 6.0 cm$^2$; clinical stenosis develops when the area is reduced to $<2.5$ cm$^2$, with severe stenosis defined as a valve area $\le 1.5$ cm$^2$ (and very severe $<1.0$ cm$^2$, with mean transmitral gradient $>10$ mmHg at rest). Chronically elevated left atrial pressure causes marked left atrial enlargement, pulmonary venous hypertension, reactive pulmonary arterial vasoconstriction, and eventual right heart failure (tricuspid regurgitation, hepatomegaly, ascites, peripheral edema). Symptoms include exertional dyspnea, orthopnea, fatigue, hemoptysis (rupture of pulmonary-bronchial venous varices), and hoarseness (Ortner's syndrome, compression of the left recurrent laryngeal nerve by a giant dilated left atrium). Auscultatory findings in pliable rheumatic valves feature a loud first heart sound ($S_1$), a high-pitched Opening Snap (OS) heard in early diastole (the $A_2\text{--OS}$ interval shortens with increasing severity), and a low-pitched, rumbling mid-diastolic murmur heard best at the apex with the bell in the left lateral decubitus position, concluding with presystolic accentuation in sinus rhythm. Atrial fibrillation occurs in $>40\%$ of patients due to atrial dilation, carrying an extreme risk of systemic left atrial appendage thromboembolism and ischemic stroke; lifelong anticoagulation (warfarin is mandatory for moderate-to-severe rheumatic MS) is indicated. In patients with favorable valvular morphology (Wilkins echocardiographic score $\le 8$, evaluating leaflet mobility, thickness, calcification, and subvalvular thickening without left atrial thrombus or $>2+$ MR), Percutaneous Balloon Mitral Valvuloplasty (PBMV) is the preferred first-line intervention; otherwise, surgical mitral valve replacement is performed.

\medterm{Mitral Valve Prolapse}-- A valvular heart disease where the mitral valve leaflets prolapse into the left atrium during systole. Most patients are asymptomatic; some develop palpitations, atypical chest pain, or dyspnea. Diagnosis by echocardiography. Complications include mitral regurgitation, infective endocarditis, and arrhythmias. Most cases require only monitoring; severe regurgitation may need surgical repair.

\medterm{Mitral-valve repair}-- A surgical or transcatheter cardiac intervention aimed at restoring structural and functional competence to a regurgitant native mitral valve while completely preserving native valvular tissue, subvalvular chordal apparatus, and papillary muscle architecture. Mitral valve repair is the gold-standard, guideline-recommended therapy for severe Primary (Degenerative) Mitral Regurgitation (most commonly caused by myxomatous degeneration, Barlow disease, fibroelastic deficiency, or chordal rupture with leaflet flail). Compared to prosthetic mitral valve replacement, repair is associated with superior preservation of left ventricular systolic function (preserving the continuous fibrous continuity between the mitral annulus, chordae tendineae, and LV papillary muscles), lower operative mortality, superior long-term survival, and avoids the lifelong bleeding and thromboembolic risks of chronic systemic anticoagulation required for mechanical prostheses. Surgical techniques pioneered by Alain Carpentier include quadrangular or triangular leaflet resection, artificial Gore-Tex neochordae placement, sliding plasty, and mandatory annuloplasty (implanting a rigid or semi-rigid prosthetic annuloplasty ring to reshape and stabilize the dilated mitral annulus). In patients with severe secondary (functional) mitral regurgitation or those with primary MR who are at prohibitive or high surgical risk, Transcatheter Edge-to-Edge Repair (TEER, using the MitraClip device) is performed percutaneously via transseptal venous access, grasping and clipping the opposing anterior and posterior leaflets at the origin of the regurgitant jet to create a functional double-orifice valve (reproducing the surgical Alfieri stitch).

\medterm{Mondor’s disease}-- A rare, benign, typically self-limiting form of superficial thrombophlebitis affecting the subcutaneous veins of the anterior chest wall, breast, or epigastrium, first comprehensively described by the French surgeon Henri Mondor in 1939. The condition most frequently involves the lateral thoracic vein, the thoracoepigastric vein, or the superior epigastric vein. Patients present with the sudden appearance of a painful, tender, cord-like subcutaneous band or string-like induration on the chest wall or breast, which becomes accentuated and visible as a shallow groove or skin tethering when the arm is raised or the breast is stretched. Pathogenesis is governed by local vascular injury, stasis, and localized hypercoagulability (Virchow's triad). While many cases are idiopathic, prominent triggers include prior breast surgery (biopsy, aesthetic augmentation, reduction, or lumpectomy), local blunt chest trauma, vigorous upper-body athletics, tight compressive clothing, and underlying systemic hypercoagulable states (such as antiphospholipid syndrome or Factor V Leiden). Importantly, in a small subset of patients, Mondor disease can be the presenting cutaneous sign of an underlying breast carcinoma compressing venous drainage; consequently, bilateral diagnostic mammography and high-resolution breast ultrasound are mandatory to exclude underlying malignancy. Treatment is conservative and symptomatic, consisting of warm compresses, oral nonsteroidal anti-inflammatory drugs (NSAIDs), and supportive bras; complete clinical resolution typically occurs spontaneously over 4 to 8 weeks.

\medterm{Multifocal Atrial Tachycardia}-- A distinctive supraventricular tachyarrhythmia (MAT) characterized by an irregular ventricular rhythm driven by multiple competing ectopic pacemaker foci within the atria, firing at rates typically between 100 and 150 beats per minute. Diagnostic 12-lead electrocardiographic criteria strictly require: (1) An irregular ventricular rate $>100$ beats per minute; (2) Presence of discrete, organized P waves with at least three distinct, identifiable P-wave morphologies (varying in amplitude, axis, and contour) in a single lead; (3) Varying, irregular P-P, P-R, and R-R intervals; and (4) An isoelectric baseline between P waves (distinguishing MAT from the chaotic, baseline fibrillatory oscillations of atrial fibrillation). When the ventricular rate is $<100$ bpm with the same morphological criteria, the rhythm is termed Wandering Atrial Pacemaker (WAP). Pathophysiologically, MAT is driven by triggered activity secondary to intracellular calcium overload in atrial myocytes. Greater than 60--80\% of MAT cases occur in elderly patients hospitalized with severe underlying pulmonary disease, most notably acute exacerbations of Chronic Obstructive Pulmonary Disease (COPD) or respiratory failure, triggered by profound alveolar hypoxia, hypercapnic acidosis, acute right atrial stretch from cor pulmonale, and excessive use of $\beta_2$-agonist bronchodilators and theophylline. Severe electrolyte derangements (hypokalemia and hypomagnesemia) and sepsis frequently co-exist. Treatment focuses strictly on optimizing the underlying pulmonary pathology (supplemental oxygen, bronchodilators, non-invasive ventilation) and aggressively correcting electrolyte deficits. Electrical cardioversion is ineffective and contraindicated because the rhythm is non-reentrant; if rate control is required, non-dihydropyridine calcium channel blockers (diltiazem, verapamil) are preferred over beta-blockers, which are contraindicated due to bronchospasm.

\medterm{Murmur Classification}-- Heart murmurs are classified by timing, grade, shape, location, radiation, and pitch. Timing: systolic (between S1 and S2), diastolic (between S2 and S1), or continuous. Grade: I (barely audible) to VI (heard without stethoscope). Shape: crescendo, decrescendo, or plateau. The grading system helps distinguish pathological from physiological murmurs and guides further diagnostic workup.

\medterm{Myocardial Perfusion Imaging}-- Nuclear cardiology technique using radiotracers (Tc-99m sestamibi, Tl-201) to assess blood flow to the myocardium. Performed at rest and during stress (exercise or pharmacological) to detect reversible ischemia (stress defect with rest redistribution) or fixed defects (scar from prior MI). SPECT (single-photon emission CT) and PET (positron emission tomography) are the modalities used. PET also quantifies myocardial blood flow and can detect balanced ischemia.

\medterm{Myocarditis}-- An inflammatory disease of the myocardium diagnosed by established histological, immunological, and immunohistochemical criteria. Etiologies encompass infectious pathogens (most frequently viral infections: coxsackievirus B, adenovirus, parvovirus B19, human herpesvirus 6, SARS-CoV-2, and HIV; bacterial like Borrelia burgdorferi [Lyme disease]; and protozoal like Trypanosoma cruzi [Chagas disease]), autoimmune disorders (systemic lupus erythematosus, sarcoidosis, giant cell myocarditis, and eosinophilic myocarditis / hypersensitivity), and toxic exposures (anthracyclines, cocaine, and immune checkpoint inhibitors [ICIs, such as nivolumab and pembrolizumab]). Pathophysiologically, viral replication drives direct myocyte lysis followed by a secondary, T-cell-mediated autoimmune response that damages sarcomeric proteins and microvasculature, potentially resolving or progressing to chronic dilated cardiomyopathy. Clinical presentation is highly variable, ranging from mild chest pain (often with a pericardial component, 'myopericarditis') and flu-like symptoms to acute decompensated heart failure, cardiogenic shock, and fatal ventricular arrhythmias. Serum cardiac troponins are elevated in the acute phase. The definitive non-invasive imaging modality of choice is Cardiac Magnetic Resonance (CMR) applying the updated 2018 Lake Louise Criteria, requiring at least one T2-based criterion (myocardial edema on T2-weighted imaging or elevated T2 mapping) and at least one T1-based criterion (increased myocardial T1 relaxation time, expanded extracellular volume fraction, or patchy Late Gadolinium Enhancement [LGE] displaying a non-ischemic subepicardial or mid-wall distribution). Endomyocardial Biopsy (EMB) evaluated by the Dallas criteria (inflammatory lymphocytic infiltrate with associated myocyte necrosis) remains the diagnostic gold standard for unexplained acute fulminant heart failure, hemodynamic instability, or suspected giant cell myocarditis. Treatment is supportive: standard guideline-directed medical therapy for heart failure (ACEi/ARNI, beta-blockers, MRAs), mechanical circulatory support (Impella, VA-ECMO) for cardiogenic shock, specific immunosuppressive regimens (high-dose corticosteroids and antithymocyte globulin) for giant cell or ICI-induced myocarditis, and strict avoidance of competitive athletics and vigorous exercise for 3 to 6 months to prevent sudden cardiac death.

\medterm{Myocardium}-- The thick muscular middle layer of the heart wall composed of cardiac muscle cells (cardiomyocytes). The myocardium is responsible for the contractile function of the heart. It contains intercalated discs with gap junctions that allow electrical coupling between cells. The left ventricular myocardium is significantly thicker (8-15 mm) than the right ventricle (3-5 mm) due to higher pressure demands of systemic circulation.

\medterm{Myxoma}-- A usually benign primary cardiac tumour composed of myxoid tissue, most often arising in the left atrium from the interatrial septum. It may cause positional obstruction of the mitral valve, constitutional symptoms, or systemic embolism; echocardiography is the usual initial diagnostic test.

\medterm{NSTEMI}-- Non-ST-elevation myocardial infarction: acute myocardial necrosis (elevated cardiac troponin) without persistent ST-segment elevation on ECG, caused by subtotal occlusion or distal embolization of a coronary artery. Features: chest pain, ECG changes (ST depression, T-wave inversion), elevated troponin. Management: antiplatelet (aspirin + P2Y12 inhibitor), anticoagulant, nitrates, beta-blockers, statins; risk stratification (GRACE score)—early invasive angiography for high risk, ischemia-guided for low risk. Distinguish from STEMI (transmural, needs immediate reperfusion).

\medterm{NYHA Functional Classification}-- A four-class system grading heart failure severity based on symptoms and exercise tolerance. Class I: no limitation of physical activity. Class II: slight limitation; comfortable at rest but ordinary activity causes symptoms. Class III: marked limitation; comfortable at rest but less than ordinary activity causes symptoms. Class IV: symptoms at rest; unable to carry out any physical activity without discomfort.

\medterm{Obtuse Marginal Artery}-- One or more substantial epicardial arterial branches ($OM_1, OM_2, OM_3$) that arise from the left circumflex artery (LCx) as it traverses the left atrioventricular (coronary) sulcus. The obtuse marginal arteries course down along the rounded, obtuse lateral margin of the heart, heading toward the cardiac apex to provide arterial blood supply to the anterolateral, lateral, and posterolateral free walls of the left ventricle, as well as contributing perfusion to the anterolateral papillary muscle of the mitral valve. In coronary angiography, the size, number, and distribution of obtuse marginal branches vary significantly depending on coronary dominance: in a left-dominant circulation (where the LCx gives rise to the posterior descending artery), the marginal branches are typically extensive and robust, whereas in right-dominant systems they supply primarily the lateral LV wall. Atherosclerotic stenosis or acute thrombotic occlusion of an obtuse marginal branch produces lateral wall myocardial ischemia or lateral STEMI, manifested on a 12-lead ECG as ST-segment elevation in the high lateral limb leads (I and aVL) and/or lateral precordial leads ($V_5$ and $V_6$), often accompanied by reciprocal ST-segment depression in the inferior leads (III and aVF).

\medterm{Orthostatic Hypotension}-- Decrease in SBP >=20 mmHg or DBP >=10 mmHg within 3 minutes of standing. Causes: dehydration, medications (antihypertensives, diuretics, alpha-blockers, antidepressants), autonomic dysfunction (diabetes, Parkinson disease), and prolonged bed rest. Symptoms: dizziness, lightheadedness, syncope, falls. Evaluation: orthostatic vital signs, autonomic reflex testing. Treatment: address reversible causes, increase fluid/salt intake, compression stockings, midodrine or fludrocortisone.

\medterm{Orthostatic Intolerance}-- A clinical disorder characterized by the inability of the cardiovascular and autonomic regulatory systems to maintain adequate systemic perfusion and cerebral blood flow upon assuming an upright, standing posture, resulting in symptoms that are relieved promptly by recumbency. Upon standing, gravity causes the rapid pooling of 500 to 1,000 mL of venous blood in the lower extremities and splanchnic circulation, reducing venous return, central blood volume, and cardiac stroke volume. In healthy physiology, arterial baroreceptors in the carotid sinus and aortic arch sense the transient drop in pressure, triggering an instantaneous compensatory reflex: vagal withdrawal and sympathetic efferent activation increase heart rate (by 10--20 bpm) and induce systemic vasoconstriction to maintain mean arterial pressure. Orthostatic intolerance arises when these compensatory mechanisms fail or are inappropriately exaggerated, presenting with lightheadedness, dizziness, visual graying or dimming, presyncope, palpitations, cognitive 'brain fog', and fatigue. Major clinical subtypes include: (1) Orthostatic (Postural) Hypotension: a sustained reduction in systolic blood pressure of $\ge 20$ mmHg or diastolic blood pressure of $\ge 10$ mmHg within 3 minutes of standing (caused by volume depletion, medications [antihypertensives, diuretics], or neurogenic autonomic failure [diabetes, Parkinson disease, multiple system atrophy]); (2) Postural Orthostatic Tachycardia Syndrome (POTS): a sustained heart rate increase of $\ge 30$ bpm (or $\ge 40$ bpm in adolescents) within 10 minutes of standing in the ABSENCE of orthostatic hypotension, predominantly affecting young females; and (3) Neurally Mediated (Vasovagal) Syncope. Diagnostic evaluation relies on active standing tests and Head-Up Tilt Table Testing (HUTT).

\medterm{Osler Nodes}-- Painful, erythematous nodules on the pads of fingers and toes, caused by immune complex deposition in infective endocarditis. They are a minor criterion in the Duke criteria. Osler nodes are tender and may last hours to days. They must be differentiated from Janeway lesions (painless, erythematous/purpuric macules on palms and soles caused by septic microemboli). Both are classic peripheral stigmata of infective endocarditis.

\medterm{P Wave}-- The first deflection on the ECG, representing atrial depolarization. Normal P waves are upright in leads I and II, inverted in aVR, and less than 0.12 seconds (3 small boxes) in duration and 2.5 mm in amplitude. Abnormal P waves indicate atrial enlargement: peaked P waves (right atrial enlargement), notched P waves (left atrial enlargement). Absent P waves suggest atrial fibrillation or junctional rhythm.

\medterm{P2Y12 Inhibitor}-- A class of oral and intravenous antiplatelet therapeutics that selectively inhibit the purinergic $P2Y_{12}$ receptor on the platelet surface, forming an essential cornerstone of Dual Antiplatelet Therapy (DAPT) alongside aspirin in coronary artery disease. Under physiological conditions, adenosine diphosphate (ADP) released from platelet dense granules binds to $G_i$-coupled $P2Y_{12}$ receptors, suppressing adenylate cyclase, decreasing intracellular cAMP, and driving the inside-out conformational activation of platelet integrin $\alpha_{IIb}\beta_3$ (GPIIb/IIIa), which cross-links fibrinogen to mediate platelet aggregation. $P2Y_{12}$ inhibitors are divided into two chemical classes: (1) Thienopyridines (irreversible prodrugs requiring hepatic CYP bioactivation): Clopidogrel (a second-generation thienopyridine requiring two-step CYP2C19 oxidation, subject to pharmacogenomic resistance in CYP2C19 loss-of-function poor metabolizers) and Prasugrel (a third-generation thienopyridine achieving faster, more potent, and more consistent platelet inhibition, indicated in ACS undergoing PCI, but contraindicated in patients with prior stroke/TIA or age $\ge 75$); and (2) Non-thienopyridines (direct-acting, reversible antagonists): Ticagrelor (an oral cyclopentyltriazolopyrimidine that does not require metabolic activation, administered twice daily, which uniquely inhibits adenosine uptake via ENT1, causing transient dyspnea and ventricular pauses) and Cangrelor (an intravenous, ultra-rapid-acting, reversible analogue with an elimination half-life of 3 to 6 minutes, used as an infusion during PCI in $P2Y_{12}$-naive patients). DAPT is mandatory following percutaneous coronary intervention (PCI) with drug-eluting stents (typically for 6 to 12 months) to prevent catastrophic acute stent thrombosis.

\medterm{Pacemaker}-- An implanted electronic medical device comprising a pulse generator (containing a lithium battery and microprocessor) connected to one or more transvenous or epicardial insulated leads containing pacing electrodes anchored within the myocardium. A pacemaker continuously senses the heart's intrinsic electrical depolarizations and delivers timed electrical stimuli (current pulses above the pacing threshold) to capture the myocardium whenever intrinsic heart rate falls below a programmed minimum rate. Pacemaker function is designated by the standard five-letter NASPE/BPEG (NBG) code: Position I = Chamber paced (A = Atrium, V = Ventricle, D = Dual [A+V]); Position II = Chamber sensed (A, V, D, O = None); Position III = Response to sensing (I = Inhibited, T = Triggered, D = Dual [Inhibited + Triggered], O = None); Position IV = Rate modulation (R = Rate-responsive via accelerometer); and Position V = Multisite pacing. Key pacing modes include: (1) VVI/VVIR (single-ventricular pacing, used for atrial fibrillation with slow ventricular response); (2) DDD/DDDR (dual-chamber pacing, maintaining physiological atrioventricular synchrony and 'atrial kick' for sinus node dysfunction and complete AV block); and (3) Biventricular Pacing / Cardiac Resynchronization Therapy (CRT-P), which paces both the right and left ventricles simultaneously to correct electromechanical dyssynchrony in heart failure patients with LVEF $\le 35\%$ and wide QRS (LBBB $\ge 150$ ms). Major clinical indications include symptomatic sick sinus syndrome (tachy-brady syndrome), symptomatic chronotropic incompetence, high-grade second-degree AV block (Mobitz Type II), and third-degree (complete) AV block. Leadless pacemakers (e.g., Micra) are miniature self-contained capsules implanted directly into the right ventricular trabeculae via femoral transcatheter delivery, eliminating transvenous lead and surgical pocket complications.

\medterm{Pacemaker Complications}-- Early: pneumothorax, hematoma, cardiac perforation/tamponade, lead dislodgement. Late: lead fracture, battery depletion, infection, tricuspid regurgitation from lead interference, venous thrombosis, and pacemaker syndrome (AV dyssynchrony in VVI pacing). Electromagnetic interference from devices (cautery, TENS, certain industrial equipment). MRI-conditional pacemakers reduce restrictions but require specific programming. Regular interrogation and follow-up essential.

\medterm{Palpitation}-- A subjective, sensory awareness of the beating of the heart, commonly described by patients as a sensation of thumping, pounding, racing, fluttering, flopping, or 'skipping' a beat in the chest, throat, or neck. Although frequently benign, palpitations warrant systematic clinical evaluation to exclude underlying life-threatening cardiovascular pathology. Etiologies are categorized as: (1) Cardiac Arrhythmias: (a) Premature beats: premature atrial contractions (PACs) or premature ventricular contractions (PVCs), where the patient often feels the compensatory pause followed by a forceful beat ('thump') due to increased end-diastolic filling; (b) Supraventricular tachycardias: atrial fibrillation (irregularly irregular pounding), atrial flutter, AVNRT, or AVRT (sudden-onset, rapid, regular tachycardia, often terminating abruptly with vagal maneuvers); and (c) Ventricular arrhythmias: ventricular tachycardia or frequent non-sustained VT; (2) Structural Heart Disease: valvular regurgitation (aortic or mitral regurgitation creating large stroke volumes), atrial septal defect, hypertrophic cardiomyopathy; (3) Systemic Hyperdynamic States: hyperthyroidism / thyrotoxicosis, severe anemia, fever, pregnancy, pheochromocytoma, hypoglycemia; (4) Exogenous Stimulants and Toxins: caffeine, nicotine, alcohol (holiday heart syndrome), sympathomimetic decongestants (pseudoephedrine), amphetamines, and cocaine; and (5) Psychiatric / Psychosomatic: panic attacks, generalized anxiety disorder, somatization. Diagnostic workup begins with a resting 12-lead ECG, followed by ambulatory ECG monitoring (24--48 hour Holter monitor for daily symptoms; 14-day patch monitor or external event recorder for weekly symptoms; or implantable loop recorder [ILR] for infrequent, sporadic episodes) and echocardiography to exclude structural heart disease. Red flags include palpitations accompanied by syncope, presyncope, angina, dyspnea, or a family history of sudden cardiac death.

\medterm{Paroxysmal Supraventricular Tachycardia}-- An episodic, rapid, regular tachyarrhythmia (PSVT) arising from cardiac tissues at or above the bifurcation of the bundle of His, characterized by sudden, abrupt onset and equally abrupt termination. Ventricular heart rates typically range between 150 and 220 beats per minute with narrow QRS complexes ($<120$ ms, unless rate-dependent bundle branch block aberrancy is present). Mechanistically, PSVT is driven by reentrant electrical circuits in $>90\%$ of cases: (1) Atrioventricular Nodal Reentrant Tachycardia (AVNRT, ~60\% of cases): the reentrant circuit is localized entirely within the AV node utilizing dual AV nodal pathways (a slow pathway with a short refractory period and a fast pathway with a long refractory period); an appropriately timed PAC blocks in the fast pathway and conducts down the slow pathway, retrogradely re-entering the recovered fast pathway; retrograde P waves are buried within or immediately follow the QRS complex (pseudo-$r'$ in lead $V_1$ or pseudo-$S$ waves in inferior leads); (2) Atrioventricular Reentrant Tachycardia (AVRT, ~30\% of cases): involves an anatomically distinct extranodal accessory pathway (such as the bundle of Kent in Wolff-Parkinson-White syndrome), where orthodromic conduction travels antegrade down the AV node and retrograde up the accessory pathway, producing narrow QRS complexes with retrograde P waves visible in the ST-T segment; and (3) Focal Atrial Tachycardia (~10\%). Patients present with sudden palpitations, neck pulsations ('frog sign' caused by simultaneous atrial and ventricular contraction against a closed tricuspid valve, sending cannon $a$ waves into the jugular veins), lightheadedness, and polyuria (atrial stretch releases BNP/ANP). Acute termination of stable PSVT begins with vagal maneuvers (the modified Valsalva maneuver, bearing down with passive leg raise, or unilateral carotid sinus massage), followed by rapid intravenous push of Adenosine (6 mg, then 12 mg, which transiently blocks AV nodal conduction). Catheter ablation of the slow pathway (AVNRT) or accessory pathway (AVRT) provides a definitive cure with $>95\%$ success rates.

\medterm{Patent Ductus Arteriosus (PDA)}-- A congenital heart defect characterized by persistent patency of the ductus arteriosus, a vessel connecting the pulmonary artery to the proximal descending aorta, which normally closes within the first few days of life. The resulting left-to-right shunt leads to pulmonary overcirculation and left heart volume overload. Presents with a continuous "machinery" murmur heard best at the left upper sternal border. Management includes pharmacologic closure (NSAIDs such as indomethacin or ibuprofen), transcatheter coil or device occlusion, or surgical ligation.

\medterm{Patent Foramen Ovale}-- A normal fetal communication between the right and left atria (the foramen ovale) that fails to close after birth. Present in 25-30\% of the general population. Typically small and hemodynamically insignificant, with left-to-right shunting that is usually not clinically important. Most PFOs are incidental findings on echocardiography. However, PFO is associated with cryptogenic stroke (paradoxical embolism — thrombus traverses from the right to left atrium, bypassing the pulmonary circulation), migraine with aura, and decompression sickness in divers. Diagnosis: transthoracic or transesophageal echocardiography with bubble study (agitated saline contrast shows bubbles crossing from right to left atrium, especially with Valsalva maneuver). Treatment: antiplatelet therapy for cryptogenic stroke; percutaneous PFO closure for selected high-risk patients (large shunt, atrial septal aneurysm, recurrent stroke despite medical therapy).

\medterm{Percutaneous Coronary Intervention}-- A catheter-based procedure to open blocked coronary arteries, typically involving balloon angioplasty and stent placement. Indications include acute STEMI (primary PCI), NSTEMI (early invasive strategy), and stable CAD with significant stenosis. Drug-eluting stents have reduced restenosis rates compared to bare-metal stents. Dual antiplatelet therapy is required after stent placement. PCI can be performed via radial or femoral artery access.

\medterm{Pericardial Effusion}-- The abnormal accumulation of fluid within the potential space of the fibroserous pericardial sac (which normally contains 15 to 50 mL of serous lubricating fluid). The hemodynamic and clinical consequence is determined not merely by the absolute volume of fluid (which can exceed 1,000 to 2,000 mL in chronic slowly developing effusions without collapse) but crucially by the rate of fluid accumulation: rapid accumulation of even 150 to 200 mL can overwhelm pericardial compliance, causing a steep rise in intrapericardial pressure that compresses cardiac chambers and precipitates life-threatening Cardiac Tamponade. Etiologies comprise infectious pericarditis (viral [coxsackie, echo], tuberculosis [the predominant cause in developing countries]), neoplasms (metastatic lung cancer, breast cancer, lymphoma, leukemia), acute myocardial infarction (early post-infarction effusion or late Dressler syndrome), uremia (uremic pericarditis in ESRD), systemic autoimmune diseases (SLE, rheumatoid arthritis, scleroderma), trauma, aortic dissection (type A dissection rupturing into the pericardial sac), and iatrogenic injury during percutaneous cardiac catheterization. Diagnostic evaluation begins with a 12-lead ECG, which characteristically reveals low QRS voltages across all leads and Electrical Alternans (beat-to-beat alternating amplitude of the QRS complexes caused by the heart mechanically swinging within the fluid-filled sac). Chest radiography shows a symmetrically enlarged, globular cardiac silhouette ('water-bottle heart'). Transthoracic echocardiography is the definitive diagnostic imaging modality, visualizing an echo-free space surrounding the heart and assessing for tamponade physiology (diastolic right ventricular free wall collapse, right atrial late diastolic inversion, and respiratory inflow Doppler variations). Diagnostic and therapeutic management requires emergency Pericardiocentesis for hemodynamic collapse or large effusions, alongside fluid analysis (cytology, Gram stain, AFB, cell count) and creation of a surgical pericardial window for recurrent malignant effusions.

\medterm{Pericardial Knock}-- A high-pitched, early diastolic sound caused by sudden cessation of ventricular filling in constrictive pericarditis. It is heard at the left sternal border and may be confused with an S3 or opening snap. The knock is typically louder with inspiration and is a characteristic finding in constrictive pericarditis. It reflects the abrupt restriction of diastolic filling by the rigid pericardium.

\medterm{Pericardiocentesis}-- A procedure to aspirate fluid from the pericardial space, performed emergently for cardiac tamponade or diagnostically for pericardial effusion. The subxiphoid approach is most commonly used, with the needle directed toward the left shoulder under echocardiographic guidance. Complications include myocardial puncture, coronary artery injury, and pneumothorax. Pericardial fluid can be analyzed for cytology, culture, and chemistry to determine the cause of effusion.

\medterm{Pericarditis}-- Inflammation of the pericardium, the fibrous sac surrounding the heart. Acute pericarditis presents with sharp, pleuritic chest pain that improves with leaning forward, pericardial friction rub, and diffuse ST elevation on ECG. Causes include viral infection (most common), bacterial, uremia, MI (Dressler syndrome), radiation, and autoimmune diseases. Treatment includes NSAIDs and colchicine. Complications include pericardial effusion and cardiac tamponade.

\medterm{Peripartum Cardiomyopathy}-- An uncommon, potentially life-threatening form of dilated cardiomyopathy defined by the development of heart failure secondary to left ventricular systolic dysfunction (left ventricular ejection fraction $\text{LVEF} < 45\%$) toward the end of pregnancy (typically the last month) or in the months immediately following delivery (up to 5 months postpartum), in a woman without previously known structural heart disease and with no other identifiable etiology. Risk factors include advanced maternal age ($>30$ years), multiparity, multiple gestations (twins, triplets), preeclampsia or gestational hypertension, African ancestry, and prolonged use of tocolytic agents. The molecular pathophysiology is driven by an antiangiogenic and pro-apoptotic signaling cascade: late-gestational oxidative stress stimulates myocardial cathepsin D cleavage of the pituitary hormone prolactin into a vasotoxic, pro-apoptotic 16-kDa prolactin fragment; this fragment induces endothelial microvascular damage, cardiomyocyte apoptosis, and marked myocardial metabolic starvation. Patients present with congestive heart failure symptoms (exertional dyspnea, orthopnea, paroxysmal nocturnal dyspnea, peripheral edema, fatigue) that are frequently mistaken for normal third-trimester physiological gestational complaints, leading to hazardous diagnostic delays. Echocardiography is definitive, revealing LV dilation and depressed systolic ejection fraction ($\text{LVEF} < 45\%$). Medical management must account for fetal safety if prior to delivery: beta-1 selective blockers (metoprolol succinate), hydralazine plus nitrates, and loop diuretics are safe; ACE inhibitors, ARBs, ARNIs, and MRAs are strictly contraindicated during pregnancy due to severe fetotoxicity and renal dysgenesis, but can be initiated postpartum. The prolactin-release inhibitor Bromocriptine has emerged as a disease-specific disease-modifying therapy. Long-term prognosis is favorable in approximately 50--70\% of patients who recover normal LV systolic function ($\text{LVEF} \ge 50\%$) within 6 to 12 months. However, patients with persistent LV dysfunction face high mortality and subsequent pregnancies carry an extreme 30--50\% risk of catastrophic heart failure relapse; subsequent pregnancy is strictly contraindicated if LVEF does not fully normalize.

\medterm{Peripheral Arterial Disease}-- Atherosclerotic occlusive disease of arteries supplying the lower limbs, causing reduced blood flow. Features: claudication (exertional calf pain relieved by rest), rest pain, diminished pulses, cool pale skin, ulcers/gangrene; chronic critical limb ischemia is limb-threatening. Risk factors: smoking (strongest), diabetes, hypertension, hyperlipidemia, age. Diagnosis: ankle-brachial index (ABI <0.9), duplex ultrasound, CTA/MRA. Management: risk factor control, exercise, antiplatelets, statins, revascularization (angioplasty/stent, bypass); associated with high cardiovascular morbidity (MI, stroke).

\medterm{Peripheral Edema}-- Swelling due to accumulation of fluid in the interstitial space of the limbs, most commonly the lower extremities. Causes: increased hydrostatic pressure (heart failure, venous insufficiency, deep vein thrombosis), reduced oncotic pressure (hypoalbuminemia—liver disease, nephrotic syndrome, malnutrition), lymphatic obstruction (lymphedema), and sodium/fluid retention (renal failure). Usually bilateral with pitting in cardiac/renal disease; unilateral suggests venous or lymphatic cause. Diagnosis: history, exam (pitting vs. non-pitting), and targeted tests (BNP, albumin, renal, DVT imaging).

\medterm{Phlebitis}-- Inflammation of a vein, often affecting the lower extremities. Superficial thrombophlebitis: inflammation and thrombosis of a superficial vein (great saphenous vein most common), presenting with erythema, tenderness, warmth, and a palpable cord along the vein. Causes: intravenous cannulation (most common—chemical phlebitis from irritant drugs, mechanical irritation, bacterial infection), varicose veins, thrombophilia, pregnancy, prolonged travel, and Behçet disease. Diagnosis is clinical; doppler ultrasound should be performed to exclude concomitant deep vein thrombosis (15--20\% of cases). Treatment: warm compresses, NSAIDs, compression stockings, elevation, and ambulation. Anticoagulation is not required for isolated superficial thrombophlebitis unless it extends to the saphenofemoral junction (then treat as DVT). Septic thrombophlebitis requires IV antibiotics and occasionally surgical excision. Deep vein thrombophlebitis (DVT) is a distinct entity requiring anticoagulation.

\medterm{Phlegmasia cerulea dolens}-- A rare, catastrophic, limb-threatening vascular emergency representing the most severe manifestation of acute deep vein thrombosis, characterized by massive, near-total occlusive thrombosis of the entire deep, superficial, and collateral venous drainage system of an extremity (most commonly involving the common iliac and femoral veins of the left lower extremity). Typically preceded by Phlegmasia Alba Dolens ('milk leg', severe pale, non-cyanotic edema from deep venous thrombosis with preserved collateral drainage), Phlegmasia Cerulea Dolens ('painful blue edema') develops when collateral venous outflow is completely obliterated. The resulting complete venous outflow obstruction causes a catastrophic rise in capillary and interstitial hydrostatic pressure. Fluid extravasation causes massive, tense, brawny edema of the entire limb, culminating in an acute compartmental syndrome: interstitial tissue pressure exceeds arteriolar perfusion pressure, causing secondary collapse of arterial inflow. This triggers arterial ischemia and Venous Gangrene. Patients present with a triad of agonizing pain, massive edema, and deep violaceous cyanosis with cutaneous bullae and loss of distal arterial pulses. Over 50\% of cases occur in the setting of advanced underlying malignancy; other causes include hypercoagulable states and severe immobility. Mortality exceeds 20--40\%, and limb amputation rates range from 20 to 50\%. Immediate management requires systemic unfractionated heparin, limb elevation, and urgent limb-saving venous revascularization via catheter-directed thrombolysis (CDT) or surgical venous thrombectomy.

\medterm{Pneumopericardium}-- Air in the pericardial space, usually from trauma (penetrating chest injury), esophageal perforation, or iatrogenic causes (central line placement, cardiac surgery). Can cause cardiac tamponade if under tension. Diagnosis by chest X-ray showing air surrounding the heart (air-fluid level if combined with effusion). Treatment: observation if small and asymptomatic; pericardial drainage if causing tamponade.

\medterm{Post-Cardiac Arrest Care}-- A comprehensive, bundle-driven critical care protocol implemented immediately following the Return of Spontaneous Circulation (ROSC) after cardiac arrest, designed to mitigate the multifaceted Post-Cardiac Arrest Syndrome (PCAS). PCAS comprises four key pathological processes: (1) Post-cardiac arrest brain injury (ischemic-reperfusion cerebral edema, neuro-inflammation, seizures); (2) Post-cardiac arrest myocardial dysfunction (transient global myocardial stunning with low cardiac output and elevated filling pressures); (3) Systemic ischemia/reperfusion response (widespread systemic inflammatory cascade mimicking septic shock with vasoplegia); and (4) Persistent precipitating pathology. Immediate resuscitation priorities encompass hemodynamic optimization (maintaining $\text{MAP} \ge 65$ mmHg using norepinephrine, avoiding hypotension), mechanical ventilatory lung-protective management (targeting normocapnia with $\text{PaCO}_2$ 35--45 mmHg and normoxia with $\text{SpO}_2$ 92--98\%, strictly avoiding hyperoxia which exacerbates reactive oxygen species generation), emergent coronary angiography for all patients with ST-segment elevation on ECG or suspected acute coronary occlusion, and Targeted Temperature Management (TTM, maintaining a constant body temperature between 32$^\circ$C and 36$^\circ$C for at least 24 hours to reduce cerebral metabolic rate and suppress apoptosis, followed by controlled rewarming at 0.25--0.5$^\circ$C/h and strict fever prevention). Continuous electroencephalography (EEG) monitoring is conducted to detect non-convulsive status epilepticus. Neuroprognostication is formally withheld until at least 72 hours post-ROSC (and after normothermia and clearance of sedatives), utilizing a multimodal evaluation comprising bilateral absence of pupillary and corneal reflexes, bilateral absence of N20 cortical responses on somatosensory evoked potentials (SSEPs), serum neuron-specific enolase (NSE), and brain MRI.

\medterm{Posterior Descending Artery}-- Also termed the posterior interventricular artery (PDA), a critical coronary vessel that runs in the posterior interventricular sulcus toward the cardiac apex, providing arterial perfusion to the inferior (diaphragmatic) wall of both ventricles, the posterior one-third of the interventricular septum, and the posteromedial papillary muscle of the mitral valve. Crucially, the origin of the PDA defines clinical Coronary Dominance: in approximately 85--90\% of humans, the anatomy is Right-Dominant, meaning the PDA arises as a terminal branch of the right coronary artery (RCA) after giving off the acute marginal branches; in 7--10\%, the anatomy is Left-Dominant, with the PDA arising from the distal left circumflex artery (LCx); and in 3--5\%, it is Codominant (balanced), receiving branches from both the RCA and LCx. The PDA gives off small septal perforator branches that supply the inferior interventricular septum and, via the AV nodal artery (arising at the crux of the heart just proximal to the PDA origin in 90\% of RCA systems), supplies the atrioventricular (AV) node. Because the posteromedial papillary muscle receives a solitary, non-redundant blood supply strictly from the PDA (unlike the dual-supplied anterolateral papillary muscle), acute occlusion of the PDA or its parent vessel frequently precipitates ischemic papillary muscle dysfunction or catastrophic acute rupture, causing torrential, life-threatening acute mitral regurgitation and cardiogenic shock. Acute PDA occlusion manifests on a 12-lead ECG as an Inferior STEMI with ST elevation in leads II, III, and aVF.

\medterm{PR Interval}-- The time from the onset of the P wave to the onset of the QRS complex, representing the time for electrical conduction through the atria, AV node, and His-Purkinje system. Normal duration is 0.12-0.20 seconds (3-5 small boxes). Prolonged PR interval (>0.20 s) indicates first-degree AV block. Shortened PR interval (<0.12 s) may indicate Wolff-Parkinson-White syndrome or junctional rhythm.

\medterm{Preload}-- The end-diastolic volume or pressure in the ventricles, representing the degree of ventricular stretch before contraction. Preload is determined by venous return and is estimated by the pulmonary capillary wedge pressure (PCWP) or left ventricular end-diastolic pressure (LVEDP). According to the Frank-Starling mechanism, increased preload (within physiological limits) increases stroke volume. Pathologically elevated preload occurs in heart failure and fluid overload.

\medterm{Premature Atrial Contraction}-- A premature beat arising from an ectopic focus in the atria, producing an early P wave (often abnormal morphology) followed by a normal QRS; common and usually benign, increasing with age and caffeine/stress. ECG: early abnormal P wave, normal PR and QRS, followed by a compensatory pause. May trigger atrial fibrillation; generally asymptomatic but can cause palpitations. Treatment: reassurance, avoidance of triggers; antiarrhythmics or ablation if symptomatic or frequent.

\medterm{Premature Ventricular Contractions}-- Ectopic beats originating from the ventricles, appearing as wide, bizarre QRS complexes that occur earlier than expected. PVCs are common in the general population and are usually benign. Frequent PVCs (>10\% of total beats) may cause left ventricular dysfunction. Causes include electrolyte imbalances, stimulants, myocardial ischemia, and cardiomyopathy. Treatment is usually unnecessary unless symptomatic or causing cardiomyopathy.

\medterm{Pseudothrombocytosis due to red cell fragments}-- Falsely elevated platelet count on automated hematology analyzers caused by fragmentation of red blood cells (schistocytes) being counted as platelets. Seen in microangiopathic hemolytic anemias such as TTP, HUS, DIC, and mechanical heart valves. Manual peripheral blood smear review confirms true platelet count.

\medterm{Pulmonary Atresia}-- A congenital heart defect where the pulmonary valve is completely absent (atretic), obstructing right ventricular outflow. Two main types: pulmonary atresia with intact ventricular septum (PA-IVS) and pulmonary atresia with ventricular septal defect (PA-VSD, also called tetralogy of Fallot with pulmonary atresia). PA-IVS: the right ventricle is small and hypertensive, and pulmonary blood flow depends on a patent ductus arteriosus. PA-VSD: the right ventricle is normal-sized, and pulmonary blood flow is ductal-dependent. Both present with severe cyanosis in the first days of life as the ductus closes. Diagnosis: echocardiography shows absent pulmonary valve tissue, ductal-dependent pulmonary circulation, and RV size. Treatment: PGE1 to maintain ductal patency, followed by staged surgical palliation (BT shunt, RVOT reconstruction, or unifocalization for PA-VSD; RV decompression or Fontan for PA-IVS).

\medterm{Pulmonary Capillary Wedge Pressure}-- An indirect hemodynamic estimate of left atrial pressure (LAP) and left ventricular end-diastolic pressure (LVEDP), obtained invasively during right heart catheterization via a balloon-tipped flow-directed pulmonary artery catheter (Swan-Ganz catheter). As the catheter is floated through the right atrium, right ventricle, and main pulmonary artery into a distal branch, the balloon is inflated with 1.0 to 1.5 mL of air. The balloon wedges in a small pulmonary arteriole (2--3 mm), transiently occluding antegrade blood flow from the right ventricle. This establishes a static column of blood extending from the catheter tip through the pulmonary capillary bed and pulmonary veins directly into the left atrium; because there is no flow, vascular resistance is zero, and the pressure at the catheter tip equals left atrial pressure. Normal resting PCWP ranges from 6 to 12 mmHg, measured strictly at end-expiration (when intrathoracic pressure is closest to atmospheric). Measurement of PCWP is crucial for the hemodynamic differentiation of pulmonary edema and shock states: a PCWP $>18$ mmHg identifies cardiogenic pulmonary edema (left heart failure, severe mitral regurgitation/stenosis, acute LV volume overload), whereas a PCWP $\le 18$ mmHg in the setting of bilateral pulmonary infiltrates and severe hypoxemia identifies non-cardiogenic pulmonary edema (Acute Respiratory Distress Syndrome, ARDS, driven by alveolar-capillary membrane disruption). PCWP falsely overestimates LVEDP in mitral stenosis, left atrial myxoma, severe pulmonary veno-occlusive disease (PVOD), or when the catheter is positioned in West Zone 1 or 2 (where alveolar pressure exceeds venous pressure).

\medterm{Pulmonary Embolism (PE)}-- An acute, potentially fatal cardiovascular and pulmonary emergency caused by the mechanical obstruction of the pulmonary arterial bed by thromboembolic material originating within the venous circulation, most commonly dislodged Deep Vein Thrombosis (DVT) from the proximal deep veins of the lower extremities (femoral, popliteal, iliac veins, accounting for $>85\%$ of cases); less common emboli include tumor, fat, amniotic fluid, and air. Embolic occlusion acutely increases right ventricular afterload, precipitating acute right ventricular strain, RV dilation, tricuspid regurgitation, and right ventricular ischemia; as the dilated RV shifts the interventricular septum leftward, left ventricular diastolic filling is mechanically restricted, causing a catastrophic fall in cardiac output, systemic hypotension, and cardiogenic shock. Clinical presentation encompasses acute-onset dyspnea, pleuritic chest pain, tachypnea, tachycardia, cough, hemoptysis, presyncope, or sudden cardiac arrest. Electrocardiography classically demonstrates sinus tachycardia, incomplete or complete RBBB, right axis deviation, and the pathognomonic but insensitive $S_1Q_3T_3$ pattern (deep S wave in lead I, pathological Q wave in lead III, inverted T wave in lead III). Initial diagnostic risk stratification uses the Wells Score or revised Geneva Score, followed by high-sensitivity D-dimer testing in low/intermediate-risk patients; definitive confirmation is established by Computed Tomography Pulmonary Angiography (CTPA) or ventilation-perfusion (V/Q) scintigraphy. Bedside echocardiography identifies RV strain, McConnell's sign (akinesia of the RV mid-free wall with preserved apical motion), and pulmonary hypertension. Clinical severity is stratified into: (1) High-risk / Massive PE: sustained hypotension (systolic $\text{BP} < 90$ mmHg or cardiogenic shock), requiring emergency reperfusion via systemic intravenous thrombolysis (alteplase 100 mg IV over 2 hours) or surgical/catheter-directed embolectomy; (2) Intermediate-risk / Submassive PE: hemodynamically stable but with objective RV dysfunction on CT/echo and elevated cardiac biomarkers (troponin, BNP), managed with therapeutic anticoagulation and consideration of catheter-directed thrombolysis (CDT); and (3) Low-risk PE: treated with oral Direct Oral Anticoagulants (DOACs: apixaban, rivaroxaban) for at least 3 months for provoked PE and indefinitely for unprovoked or recurrent events. Inferior vena cava (IVC) filters are reserved strictly for patients with absolute contraindications to anticoagulation.

\medterm{Pulmonary Hypertension}-- A severe, progressive hemodynamic and pathophysiological condition characterized by an abnormal elevation of pressure within the pulmonary arterial circulation, leading to right ventricular pressure overload, progressive right ventricular hypertrophy and dilation, and eventual right heart failure. According to the updated 2022 ESC/ERS guidelines, Pulmonary Hypertension (PH) is hemodynamically defined via invasive Right Heart Catheterization (RHC) at rest as a Mean Pulmonary Arterial Pressure ($\text{mPAP} > 20$ mmHg). PH is classified clinically into five major World Health Organization (WHO) groups: (1) Group 1: Pulmonary Arterial Hypertension (PAH), a primary vasculopathy of small pulmonary arterioles ($<500$ $\mu$m) characterized by endothelial dysfunction, plexiform lesions, smooth muscle proliferation, and in situ thrombosis, defined by $\text{mPAP} > 20$ mmHg, Pulmonary Capillary Wedge Pressure $\text{PCWP} \le 15$ mmHg, and Pulmonary Vascular Resistance $\text{PVR} > 2$ Wood Units (etiologies include idiopathic, heritable [BMPR2 mutations], connective tissue disease [systemic sclerosis], portal hypertension, congenital systemic-to-pulmonary shunts [Eisenmenger syndrome], and drugs); (2) Group 2: PH secondary to Left Heart Disease (the most common cause, $\text{PCWP} > 15$ mmHg, resulting from HFrEF, HFpEF, or mitral/aortic valvular disease); (3) Group 3: PH secondary to Chronic Lung Diseases and Hypoxia (COPD, interstitial lung disease, obstructive sleep apnea); (4) Group 4: Chronic Thromboembolic Pulmonary Hypertension (CTEPH), caused by organized, fibrotic thromboemboli occluding pulmonary arteries, curable by surgical pulmonary endarterectomy (PEA) or balloon pulmonary angioplasty (BPA); and (5) Group 5: PH with unclear or multifactorial mechanisms (sarcoidosis, hematologic disorders). Symptoms include progressive exertional dyspnea, fatigue, angina, syncope, and signs of right heart failure (elevated JVP with prominent $v$ waves, tricuspid regurgitation murmur, hepatomegaly, ascites, peripheral edema). Targeted pharmacotherapies for Group 1 PAH target three core molecular pathways: the Endothelin pathway (endothelin receptor antagonists: ambrisentan, macitentan), the Nitric Oxide pathway (PDE5 inhibitors: sildenafil, tadalafil; soluble guanylate cyclase stimulators: riociguat), and the Prostacyclin pathway (prostacyclin analogs: epoprostenol, treprostinil; prostacyclin receptor agonists: selexipag).

\medterm{Pulmonary Stenosis}-- A right ventricular outflow tract obstruction at the valvular, subvalvular, or supravalvular level. Valvular pulmonary stenosis (most common) is often congenital, frequently associated with Noonan syndrome. Symptoms depend on severity: mild is asymptomatic; moderate presents with dyspnea, fatigue, and syncope; severe presents with cyanosis and right heart failure. Physical exam: harsh systolic ejection murmur at the left upper sternal border, radiating to the back, with a prominent RV heave and a widely split S2 with a diminished P2. Diagnosis: echocardiography with Doppler measures the peak gradient (mild <30 mmHg, moderate 30-60 mmHg, severe >60 mmHg). Treatment: balloon pulmonary valvuloplasty is the procedure of choice for severe valvular stenosis (gradient >40-50 mmHg). Surgical valvotomy or valve replacement if balloon fails or for dysplastic valves (Noonan syndrome).

\medterm{Pulmonary Vascular Resistance}-- The total frictional opposition to blood flow through the entire pulmonary vascular bed, representing the primary afterload against which the right ventricle must contract. PVR is calculated invasively during right heart catheterization using the hydraulic resistance formula: $\text{PVR} = \frac{\text{mPAP} - \text{PCWP}}{\text{CO}} \times 80$, where mPAP is mean pulmonary arterial pressure (normal: 10--20 mmHg), PCWP is pulmonary capillary wedge pressure (left-sided filling pressure), and CO is cardiac output (L/min). The transpulmonary gradient $(\text{mPAP} - \text{PCWP})$ reflects the pressure drop across the pulmonary arterial resistance vessels. PVR is reported in Wood units (mmHg$\cdot$min/L) or converted to absolute CGS units (dynes$\cdot$sec$\cdot$cm$^{-5}$) by multiplying by 80. Normal PVR is $<2.0$ Wood units ($<160$ dynes$\cdot$sec$\cdot$cm$^{-5}$). According to international consensus definitions, Pulmonary Hypertension (PH) is confirmed when $\text{mPAP} > 20$ mmHg at rest; Pulmonary Arterial Hypertension (PAH, Group 1 PH) is specifically defined as pre-capillary pulmonary hypertension characterized by a $\text{PCWP} \le 15$ mmHg and an elevated $\text{PVR} > 2.0$ Wood units ($>160$ dynes$\cdot$sec$\cdot$cm$^{-5}$). Chronic elevation of PVR forces progressive right ventricular hypertrophy, dilation, and fatal right heart failure (cor pulmonale).

\medterm{Pulmonary-artery dissection}-- Rare but life-threatening tear in the pulmonary artery wall, most commonly associated with pulmonary hypertension, Marfan syndrome, or prior pulmonary artery catheterization. Presents with acute chest pain, dyspnea, and hemodynamic collapse. Diagnosis by CT angiography or echocardiography; management requires emergent surgical repair.

\medterm{Pulseless Electrical Activity}-- A cardiac arrest rhythm organized electrical activity without a detectable pulse. PEA has many reversible causes (5 H's and 5 T's): hypovolemia, hypoxia, hydrogen ion (acidosis), hypo/hyperkalemia, hypothermia; tension pneumothorax, tamponade (cardiac), toxins, thrombosis (pulmonary), thrombosis (coronary). Treatment focuses on identifying and reversing the underlying cause while performing CPR and administering epinephrine.

\medterm{Pulsus Paradoxus}-- A decrease in systolic blood pressure greater than 10 mmHg during inspiration, normally less than 10 mmHg. Exaggerated in cardiac tamponade (due to ventricular interdependence in the confined pericardial space), severe asthma, COPD exacerbation, and constrictive pericarditis. It is measured by noting the Korotkoff sounds that appear only during expiration on sphygmomanometry. Kussmaul sign (paradoxical JVD increase with inspiration) is a related finding.

\medterm{Purkinje Fibers}-- The terminal branches of the cardiac conduction system, distributed throughout the ventricular myocardium. They conduct impulses rapidly (up to 4 m/s) to ensure coordinated ventricular contraction from apex to base. Purkinje fibers have an intrinsic firing rate of 20-40 beats per minute. They are resistant to ischemia but can become triggered pacemakers in disease states.

\medterm{Purulent Pericarditis}-- Bacterial infection of the pericardial space, most commonly caused by S. aureus, S. pneumoniae, or gram-negative organisms. Often occurs as a complication of pneumonia, endocarditis, chest surgery, or immunosuppression. Presents with fever, pleuritic chest pain, and rapidly accumulating effusion with tamponade physiology. High mortality without treatment. Requires emergent pericardiocentesis or surgical drainage plus prolonged IV antibiotics.

\medterm{QRS Complex}-- The deflection representing ventricular depolarization, consisting of a Q wave (initial downward deflection), R wave (upward deflection), and S wave (downward deflection after R). Normal duration is 0.06-0.10 seconds (1.5-2.5 small boxes). Widened QRS (>0.12 s) indicates bundle branch block, ventricular rhythm, or hyperkalemia. Abnormal Q waves indicate prior myocardial infarction.

\medterm{QT Interval}-- The time from the onset of the QRS complex to the end of the T wave, representing total ventricular depolarization and repolarization. Normal QT is 0.35-0.44 seconds. The corrected QT (QTc) adjusts for heart rate using Bazett's formula: QTc = QT / square root of RR interval. QTc prolongation (>450 ms in men, >470 ms in women) increases the risk of torsades de pointes. QTc is shortened in hyperkalemia and digoxin effect.

\medterm{Radial artery thrombosis}-- Formation of a thrombus in the radial artery, often at the site of arterial catheterization or in the setting of a hypercoagulable state or vasculitis. May cause hand ischemia with pain, pallor, and decreased pulse. Treatment includes anticoagulation and surgical intervention for limb-threatening ischemia.

\medterm{Raynaud Phenomenon}-- Exaggerated vasospastic response of digital arteries to cold or stress, causing well-demarcated color changes (white-blue-red) in fingers or toes. Primary Raynaud: idiopathic, symmetric, no tissue damage, normal capillaries. Secondary Raynaud: associated with connective tissue diseases (scleroderma, SLE), vascular disease, medications (beta-blockers), or occupational vibration. Diagnosis: clinical, nailfold capillaroscopy, ANA testing. Treatment: avoid cold triggers, calcium channel blockers (nifedipine), and prostacyclin analogs for severe cases.

\medterm{Redundant mitral valve}-- Excess leaflet tissue of the mitral valve, most commonly the posterior leaflet, which may prolapse into the left atrium during systole. Often associated with myxomatous degeneration and mitral valve prolapse. May be asymptomatic or cause mild mitral regurgitation detected incidentally on echocardiography.

\medterm{Regadenoson Pharmacological Stress}-- Selective A2A adenosine receptor agonist used as a pharmacological vasodilator for nuclear myocardial perfusion imaging. Single IV bolus injection produces coronary vasodilation with fewer side effects than adenosine (less bronchospasm, AV block). Contraindicated in severe asthma or COPD requiring oxygen. Advantages: selective coronary vasodilation, no need for dose adjustment, and simpler administration than adenosine infusion.

\medterm{Regurgitation}-- The backward flow of blood through a heart valve that fails to close properly (valvular incompetence or insufficiency). Regurgitation can affect any of the four heart valves: mitral regurgitation (most common—systolic murmur at apex radiating to axilla, caused by mitral valve prolapse, ischemic heart disease, rheumatic heart disease, dilated cardiomyopathy), aortic regurgitation (diastolic murmur at left sternal edge—chronic causes: bicuspid aortic valve, aortic root dilation, rheumatic, syphilis; acute causes: infective endocarditis, aortic dissection), tricuspid regurgitation (systolic murmur at left lower sternal edge—functional from right ventricular dilation, pulmonary hypertension, or primary from endocarditis, carcinoid), and pulmonic regurgitation (diastolic murmur at left upper sternal edge—usually benign from pulmonary hypertension). Chronic regurgitation allows compensatory ventricular dilation and eccentric hypertrophy, maintaining cardiac output for years before symptoms develop (exertional dyspnea, fatigue, palpitations). Acute severe regurgitation (e.g., flail mitral leaflet, chordal rupture) causes sudden volume overload, pulmonary edema, and cardiogenic shock. Diagnosis: echocardiography (colour Doppler—quantifies severity by jet area, vena contracta, PISA radius, and regurgitant volume). Treatment: medical therapy for chronic disease (vasodilators, diuretics), surgical repair or replacement for severe symptomatic regurgitation or when ventricular dysfunction develops.

\medterm{Renal Drug Excretion}-- The elimination of drugs by the kidney through glomerular filtration, active tubular secretion, and passive tubular reabsorption. Drugs eliminated renally (e.g., aminoglycosides, vancomycin, lithium, metformin) require dose adjustment in renal impairment. Renal clearance is affected by urine pH, tubular flow rate, and drug protein binding. Renal function is estimated by creatinine clearance or eGFR, guiding dosing decisions.

\medterm{Restrictive Cardiomyopathy}-- A cardiomyopathy characterized by impaired ventricular relaxation and filling due to stiff, noncompliant myocardium, with preserved systolic function and normal or reduced ventricular volumes. Causes include amyloidosis, sarcoidosis, hemochromatosis, endomyocardial fibrosis, and radiation. Patients present with signs of heart failure despite normal EF. Prognosis is often poor. Treatment is directed at underlying cause and managing volume overload with diuretics.

\medterm{Retinal ischemia}-- Reduced blood flow to the retina causing visual disturbance, most commonly from ophthalmic artery or central retinal artery occlusion. Presents as painless, sudden monocular vision loss. Emergency treatment includes ocular massage, intra-arterial thrombolysis, and reduction of intraocular pressure to restore perfusion.

\medterm{Rhabdomyoma}-- A benign cardiac tumour composed of striated muscle cells, occurring predominantly in infants and children and strongly associated with tuberous sclerosis. It commonly involves the ventricular myocardium and may regress spontaneously, although large or multiple tumours can cause obstruction or arrhythmias.

\medterm{Rheumatic Heart Disease}-- Valvular damage resulting from acute rheumatic fever, an autoimmune response to Group A streptococcal pharyngitis. Most commonly affects the mitral valve (stenosis/regurgitation), followed by aortic. Prevention through prompt treatment of streptococcal pharyngitis. Management: valve repair/replacement, anticoagulation for atrial fibrillation, and secondary prophylaxis with penicillin.

\medterm{Right Atrium}-- The thin-walled anterior cardiac chamber that collects systemic venous return from the superior vena cava (SVC), inferior vena cava (IVC), and coronary sinus. Internally, the chamber is divided into two distinct developmental zones separated by the crista terminalis: the posterior, smooth-walled sinus venarum (derived from the embryonic sinus venosus) and the anterior, muscular pectinate wall containing the right auricle. The interatrial septum bears the oval fossa (fossa ovalis), the remnant of the fetal foramen ovale. The sinoatrial (SA) node sits in the epicardium at the junction of the SVC and right atrium, while the atrioventricular (AV) node resides within the triangle of Koch (bounded by the tendon of Todaro, the tricuspid valve septal leaflet, and the coronary sinus ostium). Normal central venous pressure (CVP / right atrial pressure) ranges from 2 to 6 mmHg. Right atrial pressure tracings display classic waveforms: the $a$ wave (atrial contraction), $c$ wave (tricuspid bulging during early RV systole), $x$ descent (atrial relaxation), $v$ wave (atrial venous filling against closed tricuspid valve), and $y$ descent (rapid passive emptying upon tricuspid opening). Pathological giant $a$ waves occur in tricuspid stenosis or pulmonary hypertension; cannon $a$ waves occur in AV dissociation (complete heart block, VT); and prominent $v$ waves occur in severe tricuspid regurgitation.

\medterm{Right Bundle Branch}-- A continuation of the bundle of His that travels along the right side of the interventricular septum to the right ventricle. It supplies the right ventricular myocardium and Purkinje fibers. Right bundle branch block (RBBB) occurs when this branch is damaged, causing delayed right ventricular activation. RBBB is commonly associated with pulmonary embolism, right ventricular hypertrophy, and aging.

\medterm{Right Coronary Artery}-- Arises from the right aortic sinus and travels in the right atrioventricular groove. It supplies the right atrium, right ventricle, SA node (60\% of individuals), AV node (85\% of individuals), and the inferior wall of the left ventricle. RCA occlusion causes inferior myocardial infarction, which may be associated with bradycardia and heart block due to AV nodal involvement.

\medterm{Right Heart Failure}-- Failure of the right ventricle to adequately pump blood to the lungs, causing systemic venous congestion. Causes include left heart failure (most common), pulmonary hypertension, right ventricular infarction, pulmonary embolism, and constrictive pericarditis. Signs include elevated JVD, peripheral edema, hepatomegaly, and ascites. Treatment focuses on addressing the underlying cause and managing volume overload with diuretics. Inotropes (dobutamine, milrinone) may be needed for RV failure.

\medterm{Right Ventricle}-- The anterior-most, crescent-shaped, thin-walled (3--5 mm) cardiac chamber that receives deoxygenated blood from the right atrium across the tricuspid valve and propels it into the low-resistance, high-compliance pulmonary circulation. Normal RV peak systolic pressure is 15 to 30 mmHg, and end-diastolic pressure is 2 to 8 mmHg. Structurally, the RV wraps around the convex septum of the left ventricle and is divided into an inflow tract (heavily trabeculated with the moderator band / trabecula septomarginalis carrying the right bundle branch), an apical trabecular zone, and a smooth-walled muscular outflow tract (conus arteriosus / infundibulum) supporting the pulmonary valve. Because of its thin wall and crescentic geometry, the right ventricle functions predominantly as a volume pump that is exquisitely sensitive to acute increases in afterload (such as massive pulmonary embolism or acute respiratory distress syndrome), which can rapidly provoke acute cor pulmonale, RV dilation, tricuspid annular dilatation, and leftward septal shifting that impairs LV filling. Right ventricular myocardial infarction occurs in up to 40\% of acute inferior STEMIs (due to proximal right coronary artery occlusion), presenting with the classic triad of hypotension, clear lung fields, and elevated jugular venous pressure, confirmed by $>1$ mm ST elevation in right-sided ECG lead $V_4R$.

\medterm{Roth Spots}-- Retinal hemorrhages with pale centers, seen on fundoscopic examination in infective endocarditis. They are a minor criterion in the Duke criteria. Roth spots are caused by immune complex-mediated vasculitis of retinal vessels. They can also be seen in other conditions: severe anemia, leukemia, collagen vascular diseases, and diabetic retinopathy. Their presence in the appropriate clinical context supports endocarditis diagnosis.

\medterm{S3 Gallop Pathology}-- A third heart sound in adults over 40 indicates ventricular dysfunction. It represents rapid ventricular filling against a dilated, noncompliant ventricle. S3 is a hallmark of systolic heart failure (HFrEF) and is associated with increased mortality. It is best heard at the apex with the bell of the stethoscope in the left lateral decubitus position. Common causes: dilated cardiomyopathy, acute MI with volume overload, severe mitral regurgitation.

\medterm{S4 Gallop Pathology}-- A fourth heart sound indicates atrial contraction against a stiff ventricle (decreased compliance). Associated with left ventricular hypertrophy from hypertension, aortic stenosis, and hypertrophic cardiomyopathy. S4 is best heard at the apex with the bell, and is accentuated by exertion, expiration, and isorhythmic AV dissociation. It is a sign of diastolic dysfunction and may precede symptomatic heart failure.

\medterm{Scimitar Syndrome}-- A congenital cardiopulmonary anomaly involving partial or total anomalous pulmonary venous return from the right lung to the inferior vena cava, often with right-lung hypoplasia and dextroposition of the heart. It may cause recurrent respiratory infection, pulmonary hypertension, or heart failure.

\medterm{Second-Degree AV Block}-- Intermittent failure of AV conduction. Mobitz Type I (Wenckebach): progressive PR prolongation then dropped QRS; usually benign (AV nodal, narrow QRS). Mobitz Type II: constant PR interval with sudden dropped QRS; more serious (infranodal, wide QRS), risk of progression to complete heart block and Stokes-Adams attacks. Treatment: Type I often no treatment needed; Type II requires pacemaker implantation due to risk of high-grade block.

\medterm{Serpentine supravenous hyperpigmentation}-- A dermatological finding characterized by dark, wavy, serpiginous hyperpigmentation along the course of superficial veins, most commonly on the upper chest and neck. Associated with venous stasis, thrombophlebitis, and certain malignancies. The pigmentation results from hemosiderin deposition in the perivenous dermis.

\medterm{SGLT2 Inhibitor}-- Sodium-glucose cotransporter-2 inhibitors (gliflozins: empagliflozin, dapagliflozin, canagliflozin), a revolutionary class of metabolic and cardiovascular therapeutics originally developed as oral antihyperglycemic agents for type 2 diabetes mellitus that have fundamentally transformed the management of heart failure and chronic kidney disease. SGLT2 is a high-capacity, low-affinity cotransporter located in the early S1/S2 segments of the renal proximal convoluted tubule, responsible for reabsorbing approximately 90\% of filtered glucose alongside sodium. Pharmacological inhibition blocks glucose and sodium reabsorption, inducing glucosuria (excreting 60--80 g of glucose/day) and osmotic natriuresis, lowering blood glucose, reducing systemic blood pressure, and promoting modest weight loss. In the cardiovascular system, SGLT2 inhibitors provide profound clinical benefits independent of glycemic control: they restore tubuloglomerular feedback (delivering sodium to the macula densa, constricting afferent glomerular arterioles, and reducing intraglomerular hypertension, halting diabetic and non-diabetic CKD progression), shift myocardial fuel energetics toward ketone body oxidation, improve myocardial ion homeostasis (inhibiting $Na^+/H^+$ exchanger-1, NHE1), reduce preload and afterload, and attenuate cardiac remodeling. Major landmark trials (DAPA-HF, EMPEROR-Reduced, DELIVER, EMPEROR-Preserved) established that dapagliflozin and empagliflozin reduce cardiovascular death and heart failure hospitalizations across the ENTIRE spectrum of heart failure (both HFrEF and HFpEF), earning a Class I guideline recommendation regardless of diabetes status. Adverse effects include mycotic genital infections, euglycemic diabetic ketoacidosis (euDKA), volume depletion, and urinary tract infections.

\medterm{Shone Complex}-- A rare congenital cardiac anomaly (also termed Shone syndrome or Shone anomaly) first described in 1963 by John Shone, characterized by a developmental cluster of four sequential, obstructive malformations affecting the left side of the heart, leading to progressive inflow and outflow tract obstruction. The classic complete complex consists of four components: (1) Supravalvular Mitral Ring: a circumferential ridge of dense fibrous connective tissue encircling the left atrial aspect of the mitral valve, obstructing inflow into the valve; (2) 'Parachute' Mitral Valve: all mitral valve chordae tendineae insert into a solitary, anomalous papillary muscle, resulting in restricted leaflet mobility and severe mitral valve stenosis; (3) Subaortic Stenosis: a discrete fibromuscular membrane or diffuse fibromuscular tunnel situated within the left ventricular outflow tract (LVOT) below the aortic valve; and (4) Coarctation of the Aorta: a severe focal narrowing of the aortic isthmus near the ductus arteriosus. The complete form with all four lesions is rare; incomplete ('forme fruste') variants exhibiting two or three lesions are more common. Hemodynamically, the serial obstructions synergistically impede pulmonary venous return and systemic cardiac output, resulting in severe pulmonary venous hypertension, pulmonary arterial hypertension, and left ventricular hypertrophy in infancy or early childhood. Infants present with tachypnea, poor feeding, failure to thrive, and pulmonary edema. Comprehensive echocardiography and cardiac MRI delineate the multilevel obstructions. Management requires complex, staged surgical intervention, including resection of the supravalvular ring, mitral valvuloplasty, resection of subaortic membranes, and aortic coarctation repair.

\medterm{Sick Sinus Syndrome}-- Sinus node dysfunction causing a spectrum of arrhythmias: sinus bradycardia, sinus arrest, sinoatrial block, and tachycardia-bradycardia syndrome. Most common in elderly patients. Presents with syncope, fatigue, palpitations, and exercise intolerance. Diagnosis: ECG, Holter monitor (pauses >3 seconds). Treatment: permanent pacemaker implantation for symptomatic bradycardia. Anticoagulation needed for tachy-brady syndrome (atrial fibrillation component).

\medterm{Single Ventricle}-- A group of congenital heart defects where one functional ventricle supports both the systemic and pulmonary circulations. Includes tricuspid atresia, mitral atresia, double inlet left ventricle, hypoplastic left heart syndrome, and unbalanced AV canal. Pulmonary and systemic blood flow must be balanced by restrictive or non-restrictive shunts. Presents with cyanosis, heart failure, or a combination depending on pulmonary blood flow. All single ventricle patients undergo staged Fontan palliation: stage 1 (neonatal: BT shunt, PA band, or Norwood), stage 2 (4-6 months: Glenn shunt — SVC to PA), and stage 3 (2-4 years: Fontan completion — IVC to PA). Long-term complications: Fontan failure, protein-losing enteropathy, plastic bronchitis, thromboembolism, and arrhythmias.

\medterm{Sinoatrial Node}-- The heart's natural pacemaker, located at the junction of the superior vena cava and right atrium. It consists of specialized cells that spontaneously depolarize at a rate of 60-100 beats per minute, setting the normal sinus rhythm. The SA node receives autonomic innervation: sympathetic stimulation increases heart rate, while parasympathetic (vagal) stimulation decreases it. SA node dysfunction causes sick sinus syndrome.

\medterm{Splinter Hemorrhages}-- Linear hemorrhages under the nail bed, oriented longitudinally. In infective endocarditis, they are caused by septic microemboli to the nail bed capillaries. They are a minor Duke criterion. Splinter hemorrhages can also result from trauma, psoriasis, and systemic vasculitis. Multiple splinter hemorrhages in several nails are more suggestive of endocarditis than a single lesion from trauma.

\medterm{ST Segment}-- The isoelectric period between ventricular depolarization (QRS) and repolarization (T wave). It corresponds to the plateau phase of ventricular action potentials. ST segment elevation (>1 mm in limb leads or >2 mm in precordial leads) suggests acute myocardial infarction, pericarditis, or early repolarization. ST depression indicates ischemia, digoxin effect, or hypokalemia. ST segment changes are critical for diagnosing acute coronary syndromes.

\medterm{Statin}-- 3-Hydroxy-3-methylglutaryl-coenzyme A (HMG-CoA) reductase inhibitors (e.g., atorvastatin, rosuvastatin, simvastatin, pravastatin), the premier class of lipid-lowering drugs and the cornerstone of primary and secondary atherosclerotic cardiovascular disease (ASCVD) prevention. Statins act as reversible, competitive, structural analogues of HMG-CoA, binding with high nanomolar affinity to the catalytic domain of hepatic microsomal HMG-CoA reductase, the committed rate-limiting enzyme in de novo cholesterol biosynthesis. Inhibition halts the conversion of HMG-CoA to mevalonate, depleting intracellular hepatic cholesterol pools. In response, hepatocytes activate sterol regulatory element-binding protein-2 (SREBP-2), which transcriptionally upregulates the synthesis and cell-surface expression of LDL receptors (LDLR). Increased LDL receptors massively accelerate the endocytic clearance of circulating atherogenic ApoB-containing lipoproteins (reducing plasma LDL-C by 30\% to $>50\%$ with high-intensity therapy [atorvastatin 40--80 mg or rosuvastatin 20--40 mg], accompanied by modest triglyceride reductions and slight HDL-C elevations). Beyond lipid lowering, statins exert multifaceted 'pleiotropic' vascular effects mediated by depletion of downstream isoprenoid intermediates (farnesyl pyrophosphate and geranylgeranyl pyrophosphate): improving endothelial nitric oxide (eNOS) bioavailability, stabilizing vulnerable atherosclerotic fibrous caps, dampening vascular inflammation (lowering hs-CRP), and inhibiting platelet thrombus formation. Adverse effects include Statin-Associated Muscle Symptoms (SAMS, ranging from common benign myalgia to rare rhabdomyolysis with myoglobinuria), transient asymptomatic transaminase elevations, slightly increased risk of incident new-onset type 2 diabetes mellitus (far outweighed by cardiovascular event reduction), and rare autoimmune necrotizing myopathy associated with anti-HMGCR antibodies.

\medterm{STEMI}-- ST-elevation myocardial infarction: transmural myocardial necrosis from acute total occlusion of a coronary artery, an emergency requiring immediate reperfusion. ECG: ST-segment elevation in contiguous leads (or new LBBB); troponin elevated. Complications: cardiogenic shock, arrhythmias (VF), pump failure, papillary muscle rupture, pericarditis. Management: primary PCI within 90 minutes (or fibrinolysis within 30 min if PCI unavailable), antiplatelet, anticoagulant, nitrates, beta-blockers, statins, oxygen if hypoxic; time is myocardium. Outcomes improved by early reperfusion.

\medterm{Stevens-Johnson Syndrome}-- A severe, life-threatening mucocutaneous reaction characterized by widespread skin erythema, necrosis, and bullae formation with mucosal involvement (oral, ocular, genital). It is most commonly caused by drug reactions (sulfonamides, anticonvulsants, allopurinol, NSAIDs) and infections (Mycoplasma pneumoniae). It is a medical emergency requiring ICU care, wound management, infection prevention, and often burn center transfer. Mortality is 5-10\%.

\medterm{Strain Imaging}-- Echocardiographic or MRI technique measuring myocardial deformation (shortening and thickening). Global longitudinal strain (GLS) by speckle tracking is more sensitive than ejection fraction for detecting subclinical myocardial dysfunction. GLS < -18\% is normal. Used for early detection of chemotherapy cardiotoxicity, differentiating cardiomyopathies, and prognosis in heart failure. GLS may detect dysfunction before EF declines.

\medterm{Stroke Volume}-- The volume of blood ejected by the left ventricle with each heartbeat, normally 60-100 mL. It is determined by three factors: preload, afterload, and myocardial contractility. Stroke volume equals end-diastolic volume minus end-systolic volume. Reduced stroke volume occurs in heart failure, hypovolemia, and cardiogenic shock. Stroke volume can be increased by exercise, positive inotropes, or fluid administration.

\medterm{ST-Segment Elevation}-- Elevation of the ECG ST segment above the isoelectric baseline, reflecting abnormal ventricular repolarization. It may indicate acute coronary occlusion but can also occur in pericarditis, early repolarization, ventricular aneurysm, or other conditions; interpretation requires clinical context and serial testing.

\medterm{Subclavian vein thrombosis}-- Formation of a blood clot in the subclavian vein, often associated with central venous catheterization, thoracic outlet syndrome, or Paget-Schroetter syndrome (effort-related thrombosis in young, active individuals). Presents with arm swelling, pain, and cyanosis. Treatment: anticoagulation; catheter-directed thrombolysis for severe cases.

\medterm{Superior mesenteric artery syndrome}-- Compression of the duodenum between the aorta and the superior mesenteric artery, causing proximal duodenal obstruction. Presents with postprandial nausea, vomiting, and epigastric pain, often after significant weight loss that narrows the SMA-aortic angle. Treatment includes nutritional rehabilitation and surgical duodenal mobilization.

\medterm{Superior vena caval obstruction}-- Compression or thrombosis of the superior vena cava impairing venous return from the head, neck, and upper extremities. Most commonly caused by malignancy (lung cancer, lymphoma) or indwelling central venous devices. Presents with facial swelling, plethora, distended neck veins, and upper extremity edema.

\medterm{Supraventricular Tachycardia}-- A rapid rhythm (150-250 bpm) originating above the ventricles, characterized by narrow QRS complexes. Most common forms are AV nodal reentrant tachycardia (AVNRT) and AV reentrant tachycardia (AVRT). Episodes are typically paroxysmal with sudden onset and offset. Vagal maneuvers and adenosine are diagnostic and therapeutic. Recurrent SVT may require catheter ablation of the reentrant circuit.

\medterm{Syncope (Fainting)}-- A clinical syndrome defined as a transient, self-terminating loss of consciousness (TLOC) characterized by rapid onset, brief duration (typically $<1--2$ minutes), and spontaneous, complete, prompt recovery without focal neurological deficits, caused by sudden, temporary global cerebral hypoperfusion (mean arterial pressure dropping below the cerebral autoregulatory threshold, typically a cessation of cerebral blood flow for $>6--8$ seconds). Syncope is classified etiologically into three broad categories: (1) Neurally Mediated (Reflex) Syncope: the most common form, driven by inappropriate autonomic reflexes triggering profound vagal bradycardia (cardioinhibitory response) and/or sympathetic vasodilation (vasodepressor response), encompassing vasovagal syncope (provoked by fear, pain, emotional distress, prolonged standing, preceded by a classic prodrome of diaphoresis, pallor, nausea, and warmth), carotid sinus hypersensitivity, and situational syncope (micturition, defecation, coughing); (2) Orthostatic Hypotension: failure of compensatory vasoconstriction upon standing, defined as a sustained drop in SBP $\ge 20$ mmHg or DBP $\ge 10$ mmHg within 3 minutes of standing (caused by intravascular volume depletion, antihypertensive drugs, or autonomic neuropathies like diabetes or Parkinson disease); and (3) Cardiac Syncope: carries the highest mortality and requires immediate evaluation, caused either by structural cardiovascular disease (aortic stenosis, hypertrophic cardiomyopathy, atrial myxoma, massive pulmonary embolism) or primary cardiac arrhythmias (bradyarrhythmias: complete heart block, sick sinus syndrome; or tachyarrhythmias: VT, VF, long QT syndrome, Brugada syndrome, Wolff-Parkinson-White syndrome). Red flags indicating cardiac syncope include syncope during exertion, syncope in the supine position, sudden onset without prodrome, palpitations immediately prior, family history of premature sudden cardiac death, and abnormal resting 12-lead ECG.

\medterm{Systemic Vascular Resistance}-- The resistance against which the left ventricle must pump, calculated as (MAP - CVP) / CO x 80. Normal: 800-1200 dynes/sec/cm5. Elevated SVR occurs in cardiogenic shock (compensatory vasoconstriction), aortic stenosis, and systemic hypertension. Decreased SVR occurs in septic shock, anaphylaxis, and cirrhosis. SVR guides vasopressor and vasodilator therapy in critical care settings.

\medterm{Systolic Murmurs}-- Murmurs occurring between S1 and S2. Ejection systolic murmurs (crescendo-decrescendo) are caused by flow across semilunar valves (aortic stenosis, pulmonary stenosis) or increased flow across normal valves. Pansystolic (holosystolic) murmurs (plateau) are caused by regurgitation across AV valves (mitral regurgitation, tricuspid regurgitation) or left-to-right shunts (VSD). Timing with carotid pulsation helps distinguish aortic from mitral murmurs.

\medterm{T Wave}-- The deflection representing ventricular repolarization. Normal T waves are upright in leads I, II, and V3-V6, and inverted in aVR. They should be at least 10\% of the QRS amplitude in limb leads and 25\% in precordial leads. Abnormal T waves may indicate ischemia (inversion), hyperkalemia (peaked), or pericarditis (diffuse elevation). T wave changes are nonspecific and require clinical correlation.

\medterm{Tachycardia}-- A heart rate exceeding 100 beats per minute in adults at rest. Tachycardia is a sign, not a diagnosis, and requires investigation to determine the underlying rhythm and cause. Sinus tachycardia: a physiological response to increased demand—exercise, fever, anxiety, pain, dehydration, anaemia, infection, hyperthyroidism, pulmonary embolism, heart failure, or drugs ($\beta$-agonists, caffeine, cocaine). The P wave is present and normal, rate 100--150 (up to 200 in children/young adults). Treatment: address the underlying cause. Narrow QRS tachycardia (QRS <0.12 s): supraventricular tachycardia—AV nodal reentrant tachycardia, AV reentrant tachycardia, atrial fibrillation, atrial flutter, atrial tachycardia. Wide QRS tachycardia (QRS >0.12 s): ventricular tachycardia (most common), SVT with aberrancy, pre-excited atrial fibrillation (Wolf–Parkinson–White). Stable tachycardia: 12-lead ECG, vagal maneuvers, adenosine (for SVT), amiodarone/procainamide (for VT). Unstable tachycardia (hypotension, altered consciousness, chest pain, acute heart failure, shock): synchronized cardioversion. Palpitations are the most common symptom of tachycardia; syncope suggests hemodynamic compromise.

\medterm{Takayasu’s arteritis}-- A chronic granulomatous vasculitis of the aorta and its major branches, predominantly affecting young women of Asian descent. Causes arterial stenosis, occlusion, or aneurysm formation. Presents with limb claudication, absent pulses, bruit, hypertension, and constitutional symptoms. Treatment: corticosteroids and immunosuppressives.

\medterm{Takotsubo Cardiomyopathy}-- An acute, typically reversible syndrome of transient left ventricular systolic dysfunction characterized by apical and mid-ventricular hypokinesia or akinesia with compensatory basal hyperkinesis, resulting in a unique apical ballooning appearance during systole resembling a traditional Japanese octopus trapping pot (tako-tsubo). Also widely designated stress cardiomyopathy, broken heart syndrome, or ampulla cardiomyopathy. Greater than 90\% of cases occur in postmenopausal women, almost universally triggered by acute, severe emotional stress (grief, catastrophic news) or physical stress (major surgery, acute respiratory distress, severe pain, intracranial hemorrhage). Pathogenesis is driven by a massive, sudden surge in circulating and myocardial catecholamines (epinephrine, norepinephrine): supraphysiological catecholamine concentrations cause direct myocyte calcium overload, diffuse microvascular coronary spasm, and a switch in $\beta_2$-adrenergic receptor coupling from stimulatory $G_s$ to inhibitory $G_i$ pathways (which is densest at the cardiac apex, explaining the selective apical stunning). Patients present with symptoms indistinguishable from acute STEMI: acute substernal chest pain, dyspnea, ischemic ST-segment elevations (predominantly in anterior leads $V_2$--$V_5$), T-wave inversions, QT prolongation, and modest cardiac troponin elevations disproportionately low relative to the extensive regional wall motion abnormality. Definitive diagnosis (modified Mayo Clinic criteria) requires: (1) Transient LV apical/mid-ventricular wall motion abnormalities extending beyond a single coronary vascular territory; (2) Absence of angiographic obstructive coronary artery disease or acute plaque rupture; (3) New ECG abnormalities or modest troponin rise; and (4) Absence of pheochromocytoma or myocarditis. Complications include cardiogenic shock, dynamic left ventricular outflow tract (LVOT) obstruction from basal hyperkinesis (where inotropes like dobutamine are strictly contraindicated), apical mural thrombus, and ventricular arrhythmias. LV systolic function characteristically recovers completely within 4 to 8 weeks.

\medterm{Targeted Temperature Management}-- Controlled therapeutic hypothermia for comatose survivors of cardiac arrest. Target: 32-36 degrees C for 24-72 hours followed by gradual rewarming. Mechanism: reduces metabolic demand, limits reperfusion injury, and reduces free radical generation. TTM2 trial showed 33 degrees C was not superior to 36 degrees C. Active cooling methods: cold IV saline, surface cooling devices, endovascular cooling catheters. Complications: infection, electrolyte disturbances, coagulopathy.

\medterm{TCAs (Tricyclic Antidepressants)}-- Antidepressants that block reuptake of serotonin and norepinephrine, with additional anticholinergic, antihistaminic, and alpha-adrenergic blocking effects. Examples include amitriptyline, nortriptyline, imipramine, and desipramine. They are effective for depression, neuropathic pain, and migraine prophylaxis. Side effects include sedation, weight gain, dry mouth, urinary retention, constipation, orthostatic hypotension, and cardiac conduction abnormalities. Lethal in overdose due to cardiac sodium channel blockade and QT prolongation.

\medterm{Tetralogy of Fallot}-- The most common cyanotic congenital heart defect, characterized by four classic anatomical features: a large ventricular septal defect (VSD), right ventricular outflow tract obstruction (infundibular or valvular pulmonic stenosis), an overriding aorta, and right ventricular hypertrophy. The degree of right-to-left shunting across the VSD depends on the severity of right ventricular outflow tract obstruction, which determines the severity of cyanosis. Patients can experience acute hypoxemic spells ("tet spells") triggered by exertion or feeding, which are temporarily relieved by squatting. Treatment is surgical repair, typically performed in infancy.

\medterm{Third-Degree AV Block}-- Complete heart block with no conduction between atria and ventricles. Atria and ventricles beat independently (AV dissociation). Ventricular escape rhythm: junctional (40-60, narrow QRS) or ventricular (20-40, wide QRS). Causes: idiopathic fibrosis, MI (inferior: transient AV nodal; anterior: infranodal), drugs, cardiomyopathy, Lyme disease, post-surgical. Symptoms: syncope, dyspnea, fatigue, heart failure. Treatment: immediate temporary pacing (transcutaneous or transvenous) and definitive permanent pacemaker implantation.

\medterm{Thrombosed intracranial aneurysm}-- An intracranial arterial aneurysm that has undergone partial or complete thrombosis of its lumen, which may reduce rupture risk but can cause mass effect on adjacent brain structures. Presents with headache, cranial nerve palsies, or seizures depending on location. Diagnosis: CT, MRI, or conventional angiography.

\medterm{Thrombosis of the inferior vena cava}-- Formation of a blood clot within the inferior vena cava, which may occur in the setting of IVC filters, malignancy, hypercoagulable states, or extension from iliofemoral deep vein thrombosis. May cause bilateral lower extremity edema, abdominal pain, and hepatic or renal dysfunction. Requires anticoagulation and possible thrombectomy.

\medterm{Thrombotic vasculopathy}-- A spectrum of disorders characterized by thrombus formation within blood vessels of various sizes, often associated with antiphospholipid syndrome, disseminated intravascular coagulation, thrombotic thrombocytopenic purpura, or malignancy. Manifestations depend on the vascular bed involved and range from microangiopathic hemolytic anemia to organ ischemia.

\medterm{Torsades de Pointes}-- A specific, potentially fatal form of polymorphic ventricular tachycardia occurring in the setting of congenital or acquired prolongation of the QT interval on a 12-lead electrocardiogram (corrected QT interval [QTc] $>460$ ms in women, $>450$ ms in men, with risk rising exponentially at $\text{QTc} > 500$ ms). Described by Dessertenne (French for 'twisting of the points'), the ECG is characterized by wide QRS complexes that continuously twist their polarity and morphology around the isoelectric baseline in an undulating, sinusoidal fashion at rates of 160 to 250 beats per minute. Pathophysiology is initiated by early afterdepolarizations (EADs): delayed myocardial repolarization (caused by reduced outward $I_{Kr}$ potassium currents or increased late inward $I_{Na}/I_{Ca}$ currents) allows L-type calcium channels to recover from inactivation during phase 2 or 3 of the action potential, firing an triggered beat. When an EAD falls on the peak or downslope of the preceding T-wave ('R-on-T phenomenon'), it triggers Torsades de Pointes, which is maintained by transmural dispersion of repolarization and functional reentry. TdP is typically initiated by a characteristic 'short-long-short' ventricular cycle sequence (a premature beat followed by a compensatory pause, which markedly lengthens the subsequent QT interval, followed by a second premature beat). Etiologies include severe electrolyte disturbances (hypokalemia, hypomagnesemia, hypocalcemia), severe bradycardia, congenital long QT syndromes (LQTS 1--3), and an extensive list of QT-prolonging pharmaceuticals (antiarrhythmics like sotalol, dofetilide, procainamide; antipsychotics like haloperidol; macrolides, fluoroquinolones, methadone, and ondansetron). Episodes are often self-limiting, causing recurrent syncope, but can degenerate into ventricular fibrillation and sudden cardiac death. Immediate acute pharmacotherapy is intravenous magnesium sulfate (1 to 2 g IV push over 1--2 minutes, even if serum magnesium is normal, which suppresses early afterdepolarizations), alongside temporary transvenous pacing or isoproterenol infusion to increase heart rate (shortening the QT interval) and urgent defibrillation if hemodynamic collapse ensues.

\medterm{Total Anomalous Pulmonary Venous Return}-- A cyanotic congenital heart defect where all four pulmonary veins drain into the systemic venous circulation (right atrium or its tributaries) rather than the left atrium. Supracardiac (most common, 45\%): pulmonary veins drain via a vertical vein into the left innominate vein or SVC. Cardiac (25\%): drain directly into the coronary sinus or right atrium. Infracardiac (20\%): drain via a descending vein into the portal vein or IVC (most obstructed type). Mixed (10\%). Obstruction presents with severe cyanosis and respiratory distress in the newborn period. Non-obstructed TAPVR presents with mild cyanosis and right heart volume overload. An obligatory right-to-left shunt via an ASD is required for survival. Diagnosis: echocardiography shows pulmonary veins not connecting to the left atrium. Treatment: surgical reimplantation of the pulmonary veins to the left atrium, closure of the ASD. Emergency surgery for obstructed TAPVR.

\medterm{Toxic Epidermal Necrolysis}-- A severe, life-threatening drug reaction characterized by widespread full-thickness epidermal necrosis and detachment, affecting >30\% of body surface area. It represents a more severe form of Stevens-Johnson syndrome. Common causative drugs include anticonvulsants (carbamazepine, phenytoin), sulfonamides, allopurinol, and NSAIDs. Treatment is supportive, similar to burn care, with ICU admission and multidisciplinary care. Mortality is 30-50\%.

\medterm{Transcatheter Aortic Valve Replacement (TAVR)}-- A minimally invasive, catheter-based structural cardiovascular procedure that has revolutionized the treatment of severe, symptomatic aortic stenosis across all surgical risk categories (prohibitive, high, intermediate, and low risk). Developed as an alternative to conventional open-heart surgical aortic valve replacement (SAVR) requiring median sternotomy and cardiopulmonary bypass, TAVR involves deploying a bioprosthetic bovine pericardial or porcine tissue valve mounted within an expandable metal stent directly inside the diseased native calcified aortic valve. Procedural approaches include the primary transfemoral route (via the femoral artery, utilizing local anesthesia and conscious sedation), with alternative vascular access (transaxillary, transcarotid, transcaval, or transapical) reserved for patients with severe peripheral arterial disease. Two primary valve design platforms are utilized: balloon-expandable valves (cobalt-chromium frame deployed during rapid ventricular pacing to transiently halt cardiac output) and self-expanding valves (nitinol shape-memory frame deployed without pacing). TAVR restores normal aortic valve area ($>1.5--2.0$ cm$^2$) and eliminates transvalvular pressure gradients. Key potential complications include paravalvular regurgitation (leakage around the stent frame, minimized by precise CT annular sizing), new-onset high-grade atrioventricular block requiring permanent pacemaker implantation (caused by stent compression of the adjacent bundle of His in the membranous septum, higher in self-expanding valves), annular rupture, vascular access bleeding, stroke, and coronary ostial obstruction.

\medterm{Transposition of the Great Arteries}-- A congenital cardiac defect in which the aorta arises from the right ventricle and the pulmonary artery from the left ventricle, creating parallel rather than serial circulations. Survival after birth requires mixing through an atrial, ventricular, or ductal communication until urgent medical and surgical management.

\medterm{Tricuspid Regurgitation}-- Incomplete closure of the tricuspid valve allowing blood to flow backward from the right ventricle to the right atrium during systole. Causes include right ventricular dilation (from pulmonary hypertension or left heart failure), endocarditis, rheumatic disease, and carcinoid syndrome. Murmur: pansystolic murmur at the left lower sternal border increasing with inspiration. Severe TR may require surgical repair or replacement.

\medterm{Tricuspid Stenosis}-- Narrowing of the tricuspid valve opening, most commonly caused by rheumatic heart disease. It is often associated with mitral stenosis. Symptoms include fatigue, edema, and hepatic congestion. Murmur: low-pitched mid-diastolic rumble at the left lower sternal border that increases with inspiration. Diagnosis is confirmed by echocardiography. Severe TS may require surgical commissurotomy or valve replacement.

\medterm{Troponins}-- Cardiac-specific proteins (troponin I and T) released into the bloodstream during myocardial cell injury. They are the most sensitive and specific biomarkers for myocardial infarction. High-sensitivity troponin assays can detect very low levels, allowing early diagnosis. Troponins rise within 3-6 hours of injury, peak at 12-24 hours, and remain elevated for 7-14 days. Serial measurements are used to diagnose acute MI.

\medterm{Truncus Arteriosus}-- A rare, complex cyanotic congenital heart defect (accounting for ~1--2\% of congenital cardiac anomalies) characterized by the complete embryological failure of the conotruncus (truncus arteriosus and bulbus cordis) to divide into the aorta and pulmonary artery during weeks 5 to 7 of fetal development. Consequently, a single, solitary great arterial trunk arises from the heart, overriding both ventricles and supplying the systemic, pulmonary, and coronary circulations. A non-restrictive, large subtruncal ventricular septal defect (VSD) is universally present. The single truncal valve is dysplastic and malformed in $>50\%$ of cases, frequently quadricuspid or tricuspid, causing varying degrees of truncal valve regurgitation or stenosis. Truncus arteriosus is strongly associated with 22q11.2 deletion syndrome (DiGeorge syndrome, present in up to 35--40\% of infants). Anatomically categorized according to the Collett-Edwards classification into Types I through IV based on the origin of the pulmonary arteries (Type I: short main pulmonary trunk arising from the truncus; Type II: separate pulmonary ostia arising closely from the posterior truncus; Type III: separate pulmonary arteries arising from lateral aspects). Neonates present with mild cyanosis at birth; as pulmonary vascular resistance naturally drops over the first weeks of life, massive left-to-right shunting drives severe pulmonary overcirculation and heart failure (tachypnea, poor feeding, diaphoresis, cardiomegaly). Without surgical intervention, $>80\%$ of infants die in the first year of life from intractable heart failure or early Eisenmenger syndrome. Complete surgical repair (VSD closure, detachment of pulmonary arteries, and placement of an RV-to-pulmonary artery valved conduit / Rastelli-type repair) is performed within the first month of life.

\medterm{U Wave}-- A small deflection following the T wave, best seen in precordial leads V2-V3. U waves are thought to represent repolarization of the Purkinje fibers or papillary muscles. Prominent U waves are associated with hypokalemia, bradycardia, and certain drugs (digitalis, quinidine). Inverted U waves may indicate ischemia or left ventricular hypertrophy. U waves are normally small and may be absent.

\medterm{Uhl Anomaly}-- An exceedingly rare, non-familial congenital cardiac malformation characterized by the partial or near-total absence of the myocardium of the right ventricular free wall, leaving the right ventricle as a paper-thin, parchment-like translucent sac composed solely of endocardium and epicardium devoid of intervening muscle fibers. The interventricular septum, right ventricular outflow tract, and left ventricular myocardium are completely spared and structurally normal. First described by Henry Uhl in 1952, the condition arises from an abnormal, selective apoptosis of right ventricular cardiac myocytes during late cardiac embryogenesis. Because the right ventricular free wall lacks contractile myocytes, it is completely akinetic or paradoxical in its motion, functioning merely as a passive, non-contractile reservoir. Consequently, antegrade pulmonary blood flow depends entirely on the pumping force of the right atrium and the mechanical suction created by the contracting left ventricle. Uhl anomaly must be strictly distinguished clinically and histologically from Arrhythmogenic Right Ventricular Cardiomyopathy (ARVC, which is genetically inherited, exhibits fibrofatty myocardial replacement rather than total absence, and typically manifests in young adults with ventricular arrhythmias). Infants with Uhl anomaly present in early infancy or childhood with massive right-sided congestive heart failure, severe hepatomegaly, ascites, peripheral edema, cyanosis (due to right-to-left shunting across a patent foramen ovale), and low cardiac output. Echocardiography and cardiac MRI establish the diagnosis by demonstrating an enormous, paper-thin, akinetic right ventricular free wall. Without surgical palliation (such as partial right ventricular exclusion, the Starnes procedure, or conversion to a Fontan univentricular circulation), the prognosis is grim, culminating in early death.

\medterm{Unstable Angina}-- Myocardial ischemia at rest, new-onset severe angina, or a crescendo pattern (more frequent, longer, lower-threshold) without evidence of myocardial necrosis (normal troponin). Pathophysiology: acute plaque rupture with thrombosis causing transient severe coronary obstruction. Part of acute coronary syndrome; features similar to NSTEMI but without biomarker elevation. Management: hospitalization, antiplatelets (aspirin, P2Y12), anticoagulation, nitrates, beta-blockers, statins; risk stratification and coronary angiography with revascularization as indicated.

\medterm{Valsalva's Manoeuvre}-- Forced expiratory effort against a closed airway (glottis closed or mouth and nose pinched shut), typically performed by bearing down as if straining during a bowel movement. The manoeuvre has four phases: Phase I (onset of straining): increased intrathoracic pressure compresses the aorta, causing a transient rise in blood pressure. Phase II (sustained strain): decreased venous return to the heart reduces cardiac output and blood pressure, with compensatory reflex tachycardia and peripheral vasoconstriction (baroreceptor-mediated). Phase III (release): sudden drop in intrathoracic pressure causes a transient fall in blood pressure. Phase IV (overshoot): cardiac output returns but peripheral vessels remain constricted, causing blood pressure to overshoot above baseline, with reflex bradycardia. Clinical uses: assessing heart murmurs (dynamic auscultation — most murmurs decrease with Valsalva, except hypertrophic obstructive cardiomyopathy [increases] and mitral valve prolapse [moves earlier]), terminating supraventricular tachycardia (via vagal stimulation increasing vagal tone to the AV node), autonomic function testing (heart rate and blood pressure responses), detecting abdominal hernias, and diagnosing pelvic organ prolapse. Contraindications: severe aortic stenosis, recent MI, glaucoma, and retinopathy.

\medterm{Varicose Veins}-- Dilated, tortuous superficial veins, most commonly in the lower extremities. Caused by venous valve incompetence leading to venous reflux. Risk factors: family history, prolonged standing, obesity, pregnancy, and older age. Symptoms: leg heaviness, swelling, aching. Complications: skin changes, ulceration, thrombophlebitis. Treatment: compression stockings, sclerotherapy, endovenous ablation, and surgical ligation.

\medterm{Vasovagal Syncope}-- Reflex syncope mediated by increased vagal tone and withdrawal of sympathetic activity, causing bradycardia and/or vasodilation resulting in transient cerebral hypoperfusion. Triggers: prolonged standing, heat, pain, emotional stress. Preceded by prodrome (lightheadedness, nausea, diaphoresis, visual tunneling). Characterized by brief loss of consciousness with spontaneous recovery. Diagnosis: clinical history, tilt-table testing. Treatment: education, lifestyle modifications, and rarely, pacemaker implantation.

\medterm{Vein-graft aneurysms}-- Focal dilation of a venous bypass graft, most commonly a saphenous vein graft used in coronary artery bypass surgery. Rare but may cause graft thrombosis, distal embolization, compression of adjacent structures, or rupture. Diagnosis: coronary angiography or CT angiography. Treatment: re-intervention or medical management.

\medterm{Venous Thrombosis}-- The formation of a blood clot (thrombus) within a vein, most commonly in the deep veins of the lower extremities (deep vein thrombosis, DVT). Pathophysiology (Virchow triad): venous stasis (immobility, surgery, long-haul travel, obesity), hypercoagulability (pregnancy, malignancy, thrombophilia — factor V Leiden, prothrombin gene mutation, antithrombin deficiency, protein C/S deficiency, antiphospholipid syndrome), and endothelial injury (trauma, surgery, central venous catheters). DVT presents with unilateral leg pain, swelling, warmth, erythema, and Homan sign (calf pain on dorsiflexion, unreliable). Complications: pulmonary embolism (PE, the most dangerous complication, causing dyspnea, chest pain, haemoptysis, and sudden death) and post-thrombotic syndrome (chronic venous insufficiency from venous valve damage). Diagnosis: D-dimer (elevated, sensitive but not specific) and compression ultrasound (non-compressible vein is diagnostic). Ventilation-perfusion (V/Q) scan or CT pulmonary angiography for PE. Treatment: anticoagulation (heparin/LMWH bridging to warfarin or DOAC — rivaroxaban, apixaban, edoxaban, dabigatran) for minimum 3-6 months; compression stockings to prevent post-thrombotic syndrome. Duration: unprovoked DVT may require indefinite anticoagulation. Prevention (prophylaxis): early mobilisation, compression devices, and prophylactic anticoagulation in hospitalised at-risk patients.

\medterm{Ventricle}-- One of the two primary muscular pumping chambers of the human heart responsible for propelling blood through the pulmonary (right ventricle) and systemic (left ventricle) vascular beds during ventricular systole. The right ventricle operates under low pressures (15--30 mmHg systolic) to circulate blood through the pulmonary capillary bed for gas exchange, whereas the left ventricle generates high pressures (100--140 mmHg systolic) to perfuse all peripheral organ systems. The mechanical performance of both ventricles is governed by three foundational physiological determinants: preload (the end-diastolic wall tension, described by the Frank-Starling mechanism where increased stretch optimizes actin-myosin overlap to enhance contractility), afterload (the impedance against ventricular ejection during systole, dictated by vascular resistance and Laplace wall stress), and inotropy (the intrinsic myocardial contractility independent of loading conditions, mediated by intracellular calcium cycling). Coordinated electromechanical contraction is mediated by the His-Purkinje conduction network, ensuring rapid, synchronized apex-to-base ventricular depolarization. Structural, ischemic, or genetic disorders affecting the ventricles lead to systolic heart failure, diastolic dysfunction, ventricular arrhythmias, and cardiogenic shock.

\medterm{Ventricular Assist Device}-- A mechanical pump that supports heart function in patients with advanced heart failure. LVADs (left ventricular assist devices) are most common, supporting the failing left ventricle. The HeartMate 3 is the current standard with fully magnetically levitated rotors. Indications include bridge to transplant (BTT), destination therapy (DT), and bridge to decision. Complications include bleeding, infection, thromboembolism, and device malfunction.

\medterm{Ventricular Fibrillation}-- Rapid, chaotic ventricular electrical activity resulting in no effective cardiac output, leading to cardiac arrest. ECG shows irregular, disorganized waveforms with no identifiable QRS complexes. VFib is a medical emergency requiring immediate defibrillation. Causes include acute myocardial infarction, electrolyte imbalances, drug toxicity, and electrical injury. Survival decreases approximately 10\% for each minute without defibrillation.

\medterm{Ventricular Septal Defect}-- The most common congenital heart defect (30\% of all CHD), characterized by an abnormal opening in the interventricular septum allowing left-to-right shunting. Types: perimembranous (most common, 70\%, near the membranous septum), muscular (5-20\%, completely surrounded by muscle), supracristal (5-7\%, above the crista supraventricularis, associated with aortic regurgitation), and inlet (5-8\%, near the tricuspid valve, common in Down syndrome). Small VSDs may close spontaneously (especially muscular). Large VSDs present with heart failure, failure to thrive, and frequent respiratory infections. On examination: holosystolic murmur at the left lower sternal border, hyperdynamic precordium, and widely split S2. Diagnosis: echocardiography shows shunt and estimates RV pressure. Treatment: small/asymptomatic — observation; large/symptomatic — surgical patch closure or transcatheter device closure. Pulmonary vascular disease develops if unrepaired (Eisenmenger syndrome).

\medterm{Ventricular Tachycardia}-- A wide-complex tachycardia (QRS >0.12 s) originating from the ventricles, with a rate of 100-250 bpm. Monomorphic VT has uniform QRS complexes; polymorphic VT has varying QRS morphology. VT is life-threatening and may degenerate to ventricular fibrillation. Causes include coronary artery disease, cardiomyopathy, electrolyte abnormalities, and drugs. Treatment includes cardioversion for unstable VT, amiodarone for stable VT, and implantable cardioverter-defibrillator for prevention.

\medterm{Werner Syndrome}-- An autosomal-recessive progeroid disorder caused most often by pathogenic variants in the WRN gene. It typically presents after puberty with premature ageing, short stature, bilateral cataracts, scleroderma-like skin changes, thin limbs, diabetes, osteoporosis, vascular disease, and increased risk of malignancy.

\medterm{Wolff-Parkinson-White Syndrome}-- A pre-excitation syndrome caused by an accessory pathway (bundle of Kent) connecting atria and ventricles, bypassing the AV node. ECG shows short PR interval (<0.12 s), delta wave (slurred upstroke of QRS), and widened QRS. Patients are at risk for supraventricular tachycardia and, if atrial fibrillation occurs, rapid ventricular response potentially degenerating to ventricular fibrillation. Treatment includes catheter ablation of the accessory pathway.
